\documentclass[11pt,a4paper]{article}

% --- Lingua e codifica ---
\usepackage[T1]{fontenc}
\usepackage[utf8]{inputenc}
\usepackage[italian]{babel}
\usepackage{lmodern}
\usepackage{microtype}

% --- Layout ---
\usepackage[margin=2.5cm]{geometry}
\usepackage{parskip}            % paragrafi separati da spazio, senza rientro

% --- Matematica, tabelle, figure ---
\usepackage{amsmath,amssymb}
\usepackage{graphicx}
\graphicspath{{img/}}
\usepackage{tabularx}           % tabelle con colonne a larghezza automatica (X)
\usepackage{booktabs}           % tabelle "professionali" (\toprule, \midrule, \bottomrule)
\usepackage{siunitx}            % numeri e unità: \num{1e6}, \SI{3.2}{\giga\flops}
\sisetup{output-decimal-marker={,}}
\usepackage{caption}
\usepackage{subcaption}
\usepackage{float}
\usepackage{tikz}
\usetikzlibrary{positioning,arrows.meta}
\usepackage{pgfplots}
\pgfplotsset{compat=1.17}
\usepgfplotslibrary{groupplots}

% --- Codice sorgente ---
\usepackage{xcolor}
\usepackage{listings}
\lstdefinestyle{codice}{
  language=C,
  basicstyle=\ttfamily\small,
  keywordstyle=\color{blue!70!black}\bfseries,
  commentstyle=\color{green!45!black}\itshape,
  stringstyle=\color{red!60!black},
  numbers=left, numberstyle=\tiny\color{gray},
  frame=single, rulecolor=\color{gray!50},
  breaklines=true, showstringspaces=false, tabsize=4,
  morekeywords={__global__,__device__,__shared__,__syncthreads,
                blockIdx,threadIdx,blockDim,gridDim,
                MPI_Comm,MPI_Init,MPI_Finalize,MPI_Send,MPI_Recv,
                MPI_Bcast,MPI_Scatter,MPI_Gather,MPI_Allreduce}
}
\lstset{style=codice}
\lstset{literate={à}{{\`a}}1 {è}{{\`e}}1 {é}{{\'e}}1 {ì}{{\`i}}1
  {ò}{{\`o}}1 {ù}{{\`u}}1 {È}{{\`E}}1}
\renewcommand{\lstlistingname}{Listato}

% --- Link e riferimenti (va caricato per ultimo) ---
\usepackage[hidelinks]{hyperref}
\usepackage[italian]{cleveref}
\crefname{listing}{listato}{listati}
\Crefname{listing}{Listato}{Listati}

% --- Promemoria visibili nel PDF (da togliere nella versione finale) ---
\newcommand{\darevisionare}[1]{\par\noindent\textcolor{red}{\textbf{[DA REVISIONARE:} #1\textbf{]}}\par}

% --- Repository del codice (aggiornare quando il repo esiste) ---
\newcommand{\repourl}{\url{https://github.com/carloni2021/SCPA}}

% --- Dati del documento ---
\title{Prodotto tra Matrice Densa e Multivettore}
\author{Luca Carloni}
\date{\today}

\begin{document}

\begin{titlepage}
  \centering
  {\large Università degli Studi di Roma ``Tor Vergata''\par}
  \vspace{0.3em}
  {\normalsize Corso di Laurea Magistrale in Ingegneria Informatica\par}

  \vspace{4cm}
  {\Large Sistemi di Calcolo Parallelo e Applicazioni\par}
  \vspace{1em}
  \rule{\textwidth}{0.4pt}\par\vspace{0.8em}
  {\huge\bfseries Prodotto tra Matrice Densa e Multivettore\par}
  \vspace{0.5em}
  \rule{\textwidth}{0.4pt}\par

  \vspace{3cm}
  {\large Luca Carloni\par}
  \vspace{0.5em}
  {\small Codice: \repourl\par}

  \vfill
  {\normalsize Anno Accademico 2025/2026\par}
\end{titlepage}
\tableofcontents
\newpage

% =====================================================================
\section{Traccia e impostazione}\label{sec:traccia}

\subsection{Il nucleo di calcolo}\label{sec:nucleo}

Il progetto richiede di realizzare un nucleo di calcolo per il prodotto
\[
  Y \leftarrow A\,X,
\]
dove $A$ è una matrice densa di dimensione $M \times N$, mentre $X$ e $Y$ sono
due \emph{multivettori}, cioè matrici dense con poche colonne: $X$ ha
dimensione $N \times k$ e $Y$ ha dimensione $M \times k$, con $k$ piccolo. Ogni
elemento del risultato vale
\[
  y_{ij} = \sum_{l=0}^{N-1} a_{il}\, x_{lj},
  \qquad 0 \le i < M,\quad 0 \le j < k .
\]
Il problema sta a metà strada tra due nuclei classici: per $k = 1$ si riduce
al prodotto matrice--vettore, mentre per $k$ paragonabile a $M$ e $N$ diventa
un prodotto tra matrici. La traccia chiede che il codice funzioni per $k$
generico e che sia collaudato almeno per $k = 3, 6, 8, 20, 32$. Nel programma
$k$ è un parametro scelto a tempo di esecuzione.

Il valore di $k$ determina il comportamento del nucleo. Ogni elemento di $A$
partecipa a esattamente $k$ moltiplicazioni-somme: con $k$ piccolo si
eseguono poche operazioni per ogni byte letto dalla memoria, e le prestazioni
sono limitate dalla banda più che dalla capacità di calcolo. Al crescere di
$k$ aumenta il lavoro che si può svolgere su un dato già caricato. Sfruttare
questo riuso è l'idea su cui si basano tutti i kernel presentati nelle
\cref{sec:gpu,sec:cpu}, e il suo effetto è analizzato in \cref{sec:roofline}.

Anche le dimensioni $M$ e $N$ sono scelte dall'utente a tempo di esecuzione.
Per le misure la traccia chiede di far crescere le dimensioni mantenendo
almeno tre rapporti fissati tra $M$ e $N$, tra cui il caso quadrato. Si sono
scelte tre forme:
\begin{itemize}
  \item \emph{quadrata}: $M = N$;
  \item \emph{alta}: $M = 3N$;
  \item \emph{larga}: $N = 2M$.
\end{itemize}
La forma cambia il rapporto tra la dimensione di $X$ ($N \times k$) e quella
di $Y$ ($M \times k$), e quindi la quantità di dati che i processi devono
scambiarsi (\cref{sec:mpi}).

\subsection{Metrica di prestazione}\label{sec:metrica}

Come richiesto dalla traccia, le prestazioni si esprimono in operazioni in
virgola mobile al secondo:
\[
  \text{FLOPS} = \frac{2\,M\,N\,k}{T},
\]
dove $M$, $N$ e $k$ sono le dimensioni del problema globale e $T$ è il tempo
medio di una invocazione del nucleo. Il fattore $2$ conta una moltiplicazione e
una somma per ciascuno degli $M N k$ termini del prodotto. I risultati sono
riportati in GFLOPS.

Il nucleo viene invocato più volte di seguito sugli stessi dati e $T$ è la
media sulle invocazioni misurate, precedute da una fase di riscaldamento non
registrata. Il modo in cui si definisce il tempo di una singola invocazione
quando i processi sono più di uno, e la necessità del riscaldamento, sono
discussi in \cref{sec:misure}.

La traccia stabilisce che il tempo includa solo il calcolo del prodotto.

\paragraph{Escluso dal tempo (misurato a parte):}
\begin{itemize}
  \item Il preprocessamento: allocazione e generazione dei blocchi di $A$,
        padding delle righe e copia di $A$ sulla GPU. È eseguito una sola
        volta, prima delle invocazioni.
  \item La copia di $X$ dalla memoria della CPU a quella della GPU prima di
        ogni invocazione e, per i kernel che lo richiedono, la sua
        trasposizione.
  \item La copia di $Y$ dalla GPU alla CPU dopo ogni invocazione.
\end{itemize}

\paragraph{Incluso nel tempo (del nucleo):}
\begin{enumerate}
  \item La distribuzione di $X$ ai processi che ne hanno bisogno
        (\texttt{MPI\_Bcast}).
  \item Il calcolo locale del prodotto parziale su ogni processo.
  \item La raccolta e la somma dei risultati parziali in $Y$
        (\texttt{MPI\_Reduce}).
\end{enumerate}

Le comunicazioni MPI sono quindi comprese nel tempo misurato. In un prodotto
distribuito, distribuire $X$ e combinare i contributi parziali sono parte
integrante del calcolo: senza di essi il prodotto non è completo. Tra gli
spostamenti di dati, la traccia esclude esplicitamente solo quelli da e verso
la scheda CUDA, mentre le comunicazioni tra processi non sono menzionate.

Oltre al tempo del nucleo, il programma riporta separatamente il tempo di
ciascuna fase. In questo modo si può valutare quanto pesa la comunicazione
rispetto al calcolo e quanto peserebbero i trasferimenti esclusi se venissero
contati.

\subsection{Collaudo}\label{sec:collaudo}

La traccia chiede di collaudare il codice confrontandolo con
un'implementazione seriale minimale. Il riferimento usato è un prodotto a tre
cicli annidati, senza ottimizzazioni se non l'ordine dei cicli: per ogni riga
di $A$, ciascun elemento viene letto una sola volta e moltiplicato per l'intera
riga corrispondente di $X$.

Il confronto avviene sull'intero risultato $Y$, misurando l'errore relativo in
norma (di Frobenius) rispetto a una tolleranza che dipende dalla precisione e
da $N$. La verifica non raccoglie mai $A$ o $Y$ su un solo processo: ogni
processo della colonna~0 ricalcola le proprie righe di $Y$ con il riferimento
seriale, rigenerando $A$ una riga alla volta dagli indici globali, e una sola
\texttt{MPI\_Allreduce} somma i contributi all'errore. In memoria resta solo
una riga di $A$ per volta, più $X$ intera.

Il riferimento accumula sempre in doppia precisione, anche nella versione in
singola: deve essere più accurato del codice che giudica, così l'errore
misurato è quello del risultato e non la somma degli errori di due calcoli
approssimati.

% =====================================================================
\section{Server e hardware}\label{sec:hardware}

\subsection{Il server \texttt{calcpar}}\label{sec:calcpar}

Tutte le misure sono state raccolte sul server di dipartimento
\texttt{calcpar}, come richiesto dalla traccia. Si tratta di un singolo nodo
con due processori e una scheda grafica, descritti in
\cref{sec:gpu-hw,sec:cpu-hw}, a cui si accede via SSH.

L'ambiente di sviluppo è gestito con il sistema dei \emph{moduli}: ogni modulo
rende disponibile nella shell una versione precisa di uno strumento. Il
progetto usa:
\begin{itemize}
  \item \texttt{gnu/13.3.0}: compilatore GCC~13.3.0 per il codice C;
  \item \texttt{openmpi/5.0.9}: Open MPI~5.0.9, che fornisce \texttt{mpicc}
        e \texttt{mpirun};
  \item \texttt{cuda/12.8}: CUDA Toolkit~12.8 (\texttt{nvcc} V12.8.61), che
        fornisce il compilatore per la GPU e la libreria cuBLAS.
\end{itemize}

Due caratteristiche della piattaforma condizionano il progetto.

\paragraph{Un solo nodo.} Tutti i processi MPI girano sulla stessa macchina:
le comunicazioni tra processi passano per la memoria condivisa e non per una
rete, e i processi si contendono le stesse risorse, cioè la banda di memoria
della CPU e, nella versione CUDA, l'unica GPU. Le conseguenze sulle misure
sono discusse in \cref{sec:misure}.

\paragraph{MPI non CUDA-aware.} L'installazione di MPI non sa leggere o
scrivere direttamente la memoria della GPU, tutti i dati comunicati
devono quindi trovarsi nella memoria della CPU. Per questo, a ogni
invocazione, $X$ viene copiato sulla GPU dopo il broadcast e $Y$ viene
riportato sulla CPU prima della riduzione. La matrice
$A$, che non viene mai comunicata, è invece caricata sulla GPU una sola volta.

\subsection{GPU: Quadro RTX 5000}\label{sec:gpu-hw}

La GPU contiene 48 \emph{Streaming Multiprocessor} (SM), che eseguono i thread. Ogni SM esegue i thread
a gruppi di 32, detti \emph{warp}, e ospita fino a 1024 thread attivi
contemporaneamente (32 warp). Ogni SM dispone di 64 unità aritmetiche per la singola precisione ma di sole
2 per la doppia. Ogni unità esegue una FMA, cioè 2 operazioni in virgola
mobile, per ciclo: con 48 SM, alla frequenza Boost dichiarata di
\SI{1815}{MHz}, i picchi teorici sono quindi di circa \SI{11.2}{TFLOPS} in
singola precisione e \SI{348}{GFLOPS} in doppia. La
frequenza dichiarata non è però un massimo: se temperatura e consumo lo
permettono la GPU sale oltre, questo ha
reso necessario un riscaldamento a tempo prima delle misure
(\cref{sec:riscaldamento}).
I picchi effettivi, misurati, sono riportati in \cref{sec:banda}.

\paragraph{Gerarchia di memoria.} Dalla più lenta e capiente alla più veloce e
piccola:
\begin{itemize}
  \item \emph{memoria globale}: \SI{16}{GiB} GDDR6 su bus a 256 bit: con \SI{14}{Gbit/s} per pin la banda teorica è
        $14 \times 256 / 8 = \SI{448}{GB/s}$. Qui risiedono $A$, $X$ e $Y$;
  \item \emph{cache L2}: \SI{4}{MiB}, condivisa da tutti gli SM;
  \item \emph{cache L1 e memoria shared}: \SI{96}{KiB} per SM, ripartiti tra le
        due in modo configurabile. La memoria shared è gestita esplicitamente
        dal programmatore, e ogni blocco di thread può usarne al massimo
        \SI{64}{KiB};
  \item \emph{registri}: \num{65536} registri da 32 bit per SM, divisi tra i
        thread attivi. Sono la memoria più veloce, ma più registri usa ogni
        thread, meno thread possono essere attivi insieme
        (\cref{sec:gpu}).
\end{itemize}

\subsection{CPU: 2 $\times$ Xeon Silver 4210}\label{sec:cpu-hw}

Il server monta due processori Intel Xeon Silver 4210 su due \emph{socket} distinti.

\paragraph{Core.} Ogni processore ha 10 core fisici, per un totale di
\textbf{20 core}. Con l'\emph{hyper-threading} ogni core si presenta al sistema
operativo come due core logici, quindi 40 in tutto. I due core logici di uno stesso core
fisico condividono però le stesse unità di calcolo e le stesse cache: per un
nucleo che tiene costantemente occupate le unità di calcolo e la memoria, un
secondo thread sullo stesso core non porterebbe vantaggi. Per questo si usa al più
un processo per core fisico (\cref{sec:cpu}). Ogni processore ha una cache L3 condivisa tra i suoi 10 core.

\paragraph{Due nodi NUMA.} Ogni processore ha la propria memoria, collegata
direttamente ai suoi controller: \SI{32}{GiB} per socket, \SI{64}{GiB} in
totale. Un core può leggere anche la memoria dell'altro processore, ma
attraverso il collegamento tra i due socket, più lento. Il sistema si presenta
quindi come due \emph{nodi NUMA} (\emph{Non-Uniform Memory Access}),
\texttt{numactl -H} riporta una distanza relativa di 10 verso la memoria locale e
di 21 verso quella remota, cioè un accesso remoto costa circa il doppio,
perché le prestazioni siano buone, ogni processo deve lavorare sulla memoria
del proprio nodo, il progetto lo ottiene facendo generare a ogni processo il
proprio blocco di $A$ (\cref{sec:binding}) e fissando i
processi ai core (\cref{sec:cpu}).

\subsection{Banda e picco di calcolo misurati}\label{sec:banda}

Per $k$ piccolo il nucleo è limitato dalla velocità con cui legge $A$ dalla
memoria. Per valutare i kernel serve quindi il valore
\emph{reale} della banda, oltre a quello nominale. La banda è stata misurata
con due programmi dedicati (\texttt{bench/banda/}), che leggono la memoria
nello stesso modo in cui il nucleo legge $A$.

\paragraph{Banda della GPU (\texttt{banda\_gpu.cu}):} Un kernel legge un vettore di \SI{2}{GiB} di \texttt{double},
molto più grande della cache L2 da \SI{4}{MiB}, così che ogni lettura arrivi
davvero dalla memoria globale. I warp leggono elementi contigui (letture
coalescenti) e ogni thread mantiene più somme indipendenti, per tenere sempre
molte letture in corso. Il risultato è di \SI{407}{GB/s}, il \SI{91}{\percent}
dei \SI{448}{GB/s} nominali. Media e massimo quasi coincidono
(\SI{406.8}{GB/s} e \SI{407.1}{GB/s}): la misura è stabile.

\paragraph{Banda della CPU (\texttt{banda\_cpu.c}):} Si usano 20 processi MPI, uno per core fisico e legato al
proprio core, come nella versione CPU del nucleo. Ogni processo inizializza e
poi legge un proprio pezzo del vettore da \SI{2}{GiB}, così che i dati si
trovino nella memoria del suo nodo NUMA. Anche qui ogni processo usa più
somme indipendenti: altrimenti ogni somma attenderebbe la precedente e si
misurerebbe la latenza dell'addizione invece della banda. Il risultato è di
\SI{71}{GB/s} come massimo e di \SI{70.4}{GB/s} in media. Come tetto si
assume il massimo, che un kernel non può superare.

\paragraph{Perché la CPU si ferma a \SI{71}{GB/s}.} Ogni processore dispone di
6 canali di memoria, 12 in tutto, e ogni canale con un modulo DDR4-2400
installato fornisce \SI{19.2}{GB/s}. Il sistema riporta però solo quattro
moduli da \SI{16}{GiB}, uno per controller di memoria: sono occupati solo 4
canali su 12. Ne segue:
\begin{itemize}
  \item banda nominale con 4 canali: $4 \times 19.2 = \SI{76.8}{GB/s}$, di cui
        la misura raggiunge il \SI{93}{\percent};
  \item banda nominale con tutti i 12 canali occupati: \SI{230.4}{GB/s}.
\end{itemize}
La banda della CPU è quindi limitata dalla configurazione della memoria del
server, non dai processori. È un dato da tenere presente nel confronto tra le
due versioni: la GPU dispone di circa 5,7 volte la banda della CPU.

\paragraph{Picco di calcolo della GPU.} Un secondo programma
(\texttt{bench/banda/picco\_gpu.cu}) misura il picco di calcolo: ogni thread
esegue solo moltiplicazioni-somme sui propri registri, senza accedere alla
memoria, con più catene indipendenti per tenere sempre occupate le unità. Il
risultato è di \SI{374}{GFLOPS} in \texttt{double} e \SI{11.9}{TFLOPS} in
\texttt{float}, corrispondenti in entrambi i casi a una frequenza di circa
\SI{1.94}{\giga\hertz}. Sono questi i valori usati come tetto del calcolo in
\cref{sec:roofline}.

\subsection{Il modello roofline}\label{sec:roofline}

Per capire quanto un kernel è vicino al
massimo possibile si usa il modello \emph{roofline}~\cite{williams2009}. Le
prestazioni raggiungibili sono limitate da due tetti, si considera sempre il minore:
\[
  \text{GFLOPS}_{\text{massimi}} = \min\bigl(\text{picco di calcolo},\ \text{banda} \times \underbrace{\text{intensità}}_{\text{definita tra poco}}\bigr).
\]
\begin{itemize}
  \item Il \emph{tetto della memoria} è banda $\times$ intensità: quante
        operazioni si possono fare con i dati che la memoria riesce a fornire.
  \item Il \emph{tetto del calcolo} è il \emph{picco} dell'hardware: quante
        operazioni al secondo le unità di calcolo possono eseguire, qualunque
        sia la quantità di dati.
\end{itemize}
La banda si misura in byte al secondo e l'intensità in operazioni per byte: il
loro prodotto è quindi in operazioni al secondo, rappresenta quini la velocità a cui il
kernel andrebbe se il suo unico limite fosse la memoria. Il picco è invece la
velocità a cui andrebbe se fosse limitato solamente dal calcolo. Il kernel deve
rispettare entrambi i limiti contemporaneamente, quindi conta il più basso dei
due: da qui il minimo.

\paragraph{Quante operazioni per ogni byte.} L'intensità indica quanto lavoro
si fa con ogni byte letto:
\[
  \text{intensità}
  = \frac{\text{operazioni}}{\text{byte letti dalla memoria}}
  = \frac{\underbrace{k}_{\text{moltiplicazioni}} + \underbrace{k}_{\text{somme}}}
         {\text{byte per elemento}} .
\]
È la \emph{operational intensity} di~\cite{williams2009}: contano solo i byte
letti dalla memoria principale, dopo il filtro delle cache. Nel nucleo quasi
tutti i byte letti appartengono ad $A$, di gran lunga la
matrice più grande. Ogni suo elemento viene letto una volta e usato in $k$
moltiplicazioni e $k$ somme (con k numero di colonne di $X$), cioè $2k$ operazioni, dividendo per i byte
occupati da un elemento si ottiene:
\begin{itemize}
  \item in \texttt{double} (8 byte): $2k/8 = k/4$ operazioni per byte;
  \item in \texttt{float} (4 byte): $2k/4 = k/2$ operazioni per byte.
\end{itemize}
L'intensità cresce con $k$: più colonne ha $X$, più lavoro si fa con lo stesso
elemento di $A$.

\paragraph{I due tetti della GPU.} Per il tetto della memoria si usa la banda
misurata, \SI{407}{GB/s}, perché quella nominale non è raggiungibile
(\cref{sec:banda}), per il tetto del calcolo il picco misurato,
\SI{374}{GFLOPS} in \texttt{double} e \SI{11.9}{TFLOPS} in \texttt{float}
(\cref{sec:banda}). Inseriti nella
formula, il tetto della memoria diventa $407 \cdot \underbrace{k/4}_{\text{intensità}}$~GFLOPS in
\texttt{double} e $407 \cdot k/2$~GFLOPS in \texttt{float} (\cref{tab:tetti}).

\paragraph{Dove si incontrano i due tetti.} Il limite passa dalla memoria al
calcolo quando i due tetti sono uguali:
\[
  \text{banda} \times I = \text{picco}
  \quad\Rightarrow\quad
  I = \frac{\text{picco}}{\text{banda}} .
\]
In \texttt{double}
$I = 374/407 \approx 0{,}92$ operazioni per byte e, poiché $I = k/4$,
$k \approx 3{,}7$, mentre in \texttt{float} $I = \num{11876}/407 \approx 29{,}2$ e,
poiché $I = k/2$, $k \approx 58$.

\emph{In \texttt{double}}, con $k = 3$ la memoria consente al massimo
$407 \cdot 3/4 \approx 305$~GFLOPS, sotto il picco: il limite è la memoria. Già
con $k = 6$ la memoria consentirebbe 611~GFLOPS, ma la GPU in \texttt{double}
non va oltre 374: da lì in poi il limite è il calcolo, e aumentare $k$ non serve
più.

\emph{In \texttt{float}} il picco è circa 32 volte più alto, \num{11900}~GFLOPS. Anche
con $k = 32$ la memoria si ferma a $407 \cdot 32/2 \approx 6500$~GFLOPS, ancora
sotto il picco: in \texttt{float} il limite è sempre la memoria, per tutti i
$k$ della traccia.

\begin{table}[H]
  \centering
  \caption{Tetto della roofline della GPU per i $k$ della traccia, in GFLOPS.}\label{tab:tetti}
  \begin{tabular}{lccccc}
    \toprule
    $k$ & 3 & 6 & 8 & 20 & 32 \\
    \midrule
    Tetto in \texttt{double} & 305 & 374 & 374 & 374 & 374 \\
    Tetto in \texttt{float}  & 611 & 1221 & 1628 & 4070 & 6512 \\
    \bottomrule
  \end{tabular}
\end{table}

\Cref{fig:roof-double,fig:roof-float} mostrano, per ciascuna precisione, i due
tetti separati e la roofline che ne risulta, cioè il loro minimo.

\begin{figure}[H]
  \centering
\begin{tikzpicture}
\begin{axis}[
  width=0.92\textwidth, height=0.5\textwidth,
  xmode=log, ymode=log, log basis x=2, log basis y=10,
  xmin=0.125, xmax=16, ymin=40, ymax=8000,
  grid=both, major grid style={gray!20}, minor grid style={gray!8},
  xlabel={intensità $I$ (operazioni/byte)}, ylabel={GFLOPS},
  xtick={0.125,0.25,0.5,1,2,4,8,16}, xticklabels={0{,}125,0{,}25,0{,}5,1,2,4,8,16},
  xticklabel style={font=\scriptsize}, yticklabel style={font=\scriptsize},
  legend style={font=\scriptsize, at={(0.02,0.98)}, anchor=north west, cells={anchor=west}},
]
\addplot[gray!60!black, dashed, line width=1pt, domain=0.125:16, samples=2] {407*x};
\addlegendentry{tetto della memoria: $407 \cdot I$}
\addplot[red!65!black, dashed, line width=1pt, domain=0.125:16, samples=2] {374};
\addlegendentry{tetto del calcolo: \num{374}}
\addplot[blue!70!black, line width=2.6pt, opacity=0.85, domain=0.125:16, samples=500] {min(374,407*x)};
\addlegendentry{roofline: $\min(\num{374},\ 407 \cdot I)$}
\addplot[only marks, mark=o, black, mark size=3.2pt, forget plot] coordinates {(0.919,374)};
\node[font=\scriptsize, anchor=north west, align=left] at (axis cs:0.97,350) {$I = 0{,}92$\\($k \approx 3{,}7$)};
\node[font=\scriptsize, gray!40!black, anchor=south, rotate=29] at (axis cs:0.3,140) {limite di memoria};
\node[font=\scriptsize, red!55!black, anchor=south] at (axis cs:6,374) {limite di calcolo};
\end{axis}
\end{tikzpicture}
  \caption{Roofline della GPU in \texttt{double}: tetto della memoria
    (\SI{407}{GB/s}), tetto del calcolo (\SI{374}{GFLOPS}) e loro minimo. I
    due tetti si incontrano a $I = 0{,}92$, cioè $k \approx 3{,}7$.}\label{fig:roof-double}
\end{figure}

\begin{figure}[H]
  \centering
\begin{tikzpicture}
\begin{axis}[
  width=0.92\textwidth, height=0.5\textwidth,
  xmode=log, ymode=log, log basis x=2, log basis y=10,
  xmin=0.5, xmax=128, ymin=150, ymax=60000,
  grid=both, major grid style={gray!20}, minor grid style={gray!8},
  xlabel={intensità $I$ (operazioni/byte)}, ylabel={GFLOPS},
  xtick={0.5,1,2,4,8,16,32,64,128}, xticklabels={0{,}5,1,2,4,8,16,32,64,128},
  xticklabel style={font=\scriptsize}, yticklabel style={font=\scriptsize},
  legend style={font=\scriptsize, at={(0.02,0.98)}, anchor=north west, cells={anchor=west}},
]
\addplot[gray!60!black, dashed, line width=1pt, domain=0.5:128, samples=2] {407*x};
\addlegendentry{tetto della memoria: $407 \cdot I$}
\addplot[red!65!black, dashed, line width=1pt, domain=0.5:128, samples=2] {11876};
\addlegendentry{tetto del calcolo: \num{11876}}
\addplot[orange!85!black, line width=2.6pt, opacity=0.85, domain=0.5:128, samples=500] {min(11876,407*x)};
\addlegendentry{roofline: $\min(\num{11876},\ 407 \cdot I)$}
\addplot[only marks, mark=o, black, mark size=3.2pt, forget plot] coordinates {(29.18,11876)};
\node[font=\scriptsize, anchor=north west, align=left] at (axis cs:31,10800) {$I = 29{,}2$\\($k \approx 58$)};
\node[font=\scriptsize, gray!40!black, anchor=south, rotate=29] at (axis cs:3,1300) {limite di memoria};
\node[font=\scriptsize, red!55!black, anchor=south] at (axis cs:70,11876) {limite di calcolo};
\end{axis}
\end{tikzpicture}
  \caption{Roofline della GPU in \texttt{float}: tetto della memoria
    (\SI{407}{GB/s}), tetto del calcolo (\SI{11876}{GFLOPS}) e loro minimo. I
    due tetti si incontrano a $I = 29{,}2$, cioè $k \approx 58$: per tutti i
    $k$ della traccia il limite è la memoria.}\label{fig:roof-float}
\end{figure}

In breve, in \texttt{double} la GPU fatica a \emph{calcolare}, in
\texttt{float} fatica a \emph{leggere i dati}. Il grafico con i due tetti e i
risultati misurati è riportato in \cref{sec:gpu-riepilogo}.

% =====================================================================
\section{Ambiente e strumenti}\label{sec:ambiente}

\subsection{Sviluppo e compilazione}\label{sec:compilazione}

La compilazione è gestita da un \texttt{Makefile} con due parametri scelti da
riga di comando:
\begin{itemize}
  \item \texttt{BACKEND} = \texttt{cuda} (default) oppure \texttt{cpu}: sceglie
        la parallelizzazione interna al processo, sulla GPU oppure su un core
        per processo;
  \item \texttt{PREC} = \texttt{double} (default) oppure \texttt{float}: sceglie
        la precisione, fissata in compilazione (\cref{sec:architettura}).
\end{itemize}
Ne risultano quattro eseguibili distinti, ad esempio
\texttt{bin/mxm\_cuda\_double} o \texttt{bin/mxm\_cpu\_float}. Ogni
combinazione ha una propria cartella di file oggetto, così che cambiando
backend o precisione non si mescolino file compilati con opzioni diverse.

Il codice C è compilato con \texttt{mpicc}, quello CUDA con \texttt{nvcc} per
l'architettura \texttt{sm\_75}. Le opzioni principali sono riassunte nella
tabella~\ref{tab:opzioni}.

\begin{table}[H]
  \centering
  \small
  \begin{tabularx}{\textwidth}{@{}lX@{}}
    \toprule
    \textbf{Opzione} & \textbf{Effetto} \\
    \midrule
    \texttt{-O3} & Massimo livello di ottimizzazione. \\
    \addlinespace
    \texttt{-march=native} & Genera codice per il processore su cui si
      compila, cioè lo Xeon del server con le istruzioni AVX-512: per questo
      il codice va compilato sul server stesso. \\
    \addlinespace
    \texttt{-ffp-contract=fast} & Consente di fondere \texttt{a*b + c} in
      un'unica istruzione FMA. Con \texttt{-std=c11} GCC la disattiva per
      default, e senza di essa il picco di calcolo della CPU si dimezza;
      \texttt{nvcc} la applica già per default, quindi le due versioni
      risultano confrontabili. L'FMA arrotonda una sola volta invece di due,
      con un risultato anche leggermente più accurato. \\
    \addlinespace
    \texttt{-DUSE\_FLOAT} & Usata solo con \texttt{make PREC=float}: definisce
      il tipo degli elementi \texttt{real\_t} come \texttt{float} invece di
      \texttt{double}, e di conseguenza il tipo MPI scambiato tra i processi
      e la chiamata cuBLAS di riferimento (\texttt{Sgemm} invece di
      \texttt{Dgemm}). \\
    \bottomrule
  \end{tabularx}
  \caption{Opzioni di compilazione principali.}
  \label{tab:opzioni}
\end{table}

Il \texttt{Makefile} compila automaticamente tutti i sorgenti comuni e quelli
del backend scelto: aggiungere un nuovo kernel, cioè un nuovo file in
\texttt{src/cuda/} o \texttt{src/cpu/}, non richiede modifiche al
\texttt{Makefile}.

\subsection{Struttura del progetto}\label{sec:struttura}

Il codice completo è nella repo all'indirizzo \repourl{} ed è organizzato
nelle cartelle seguenti.

\paragraph{\texttt{include/}} Le interfacce della parte comune:
\begin{itemize}
  \item \texttt{common.h}: il tipo numerico \texttt{real\_t} e la macro di
        indicizzazione \texttt{IDX}, condivisi da C e CUDA;
  \item \texttt{distrib.h}: la griglia di processi e la distribuzione a
        blocchi degli indici;
  \item \texttt{matgen.h}: la generazione riproducibile delle matrici;
  \item \texttt{prep.h}: le funzioni di preprocessamento (allocazione
        allineata e padding);
  \item \texttt{backend.h}: l'interfaccia comune ai backend per il prodotto
        locale;
  \item \texttt{mxm.h}: l'interfaccia del nucleo distribuito;
  \item \texttt{k\_dispatch.h}: la suddivisione di $k$ in fette, comune ai
        kernel;
  \item \texttt{serial.h} e \texttt{stats.h}: il riferimento seriale e le
        statistiche sui tempi.
\end{itemize}

\paragraph{\texttt{src/}} Le implementazioni:
\begin{itemize}
  \item nella cartella principale, la parte comune ai due backend: il
        nucleo distribuito con la logica MPI (\texttt{mxm.c}), il programma
        principale che lo invoca e ne misura i tempi (\texttt{main.c}), la
        griglia, la generazione, il preprocessamento, il riferimento seriale e
        le statistiche;
  \item \texttt{src/cuda/}: il backend GPU, con un file per ogni kernel e uno
        per la gestione della memoria del dispositivo;
  \item \texttt{src/cpu/}: il backend CPU, con un file per ogni kernel e uno
        che implementa l'interfaccia del backend (\texttt{backend\_cpu.c}).
\end{itemize}

\paragraph{\texttt{test/}} Lo script di collaudo \texttt{run\_tests.sh}
(\cref{sec:verifica}) e un piccolo programma che stampa la griglia e i blocchi
assegnati a ogni processo, usato per verificare la distribuzione degli indici.

\paragraph{\texttt{bench/}} Gli strumenti per le misure: lo script
\texttt{campagna.sh} con i file di configurazione delle campagne
(\texttt{config/}), i programmi di misura della banda (\texttt{banda/}) con
lo script che li lancia (\texttt{banda.sh}), lo script \texttt{mem\_locale.sh},
che vincola la memoria di ogni processo CPU al proprio socket
(\cref{sec:binding}), e gli script per le verifiche più specifiche (budget dei registri della CPU,
frequenza della GPU), descritti in \cref{sec:campagne}.

\paragraph{\texttt{results/}} Per ogni campagna, il file CSV dei risultati,
una copia della configurazione e il log degli errori.

% =====================================================================
\section{Dati e preprocessamento}\label{sec:dati}

La traccia richiede funzioni ausiliarie per il preprocessamento dei dati della
matrice. Nel progetto sono raccolte in \texttt{matgen} (generazione dei dati) e
in \texttt{prep} (allocazione allineata e padding), e sono eseguite una sola
volta prima delle invocazioni del nucleo, fuori dal tempo misurato
(\cref{sec:metrica}).

\subsection{Layout in memoria}\label{sec:layout}

Tutte le matrici ($A$, $X$ e $Y$) sono memorizzate in ordine
\emph{row-major}, quello naturale del C: gli elementi di una riga occupano
posizioni consecutive e le righe si susseguono una dopo l'altra. Per accedere
agli elementi si usa un'unica macro, definita in \texttt{common.h}:
\[
  \mathrm{IDX}(i, j, \mathit{ld}) = i \cdot \mathit{ld} + j .
\]
L'elemento $(i, j)$ si trova $i$ righe intere più avanti dell'inizio, più $j$
posizioni. Il parametro $\mathit{ld}$ (\emph{leading dimension}) è la distanza
in elementi tra l'inizio di una riga e l'inizio della successiva che può variare
dal numero di colonne nel caso di matrici con padding, come per $A$ (\emph{padding},
\cref{sec:padding}). Usare ovunque la stessa macro, con $\mathit{ld}$
esplicito, evita gli errori di indicizzazione che nascerebbero dal riscrivere
il conto a mano in ogni punto del codice.

Il prodotto $i \cdot \mathit{ld}$ è calcolato nel tipo \texttt{size\_t}
(intero senza segno a 64 bit) e non in \texttt{int}: a 32 bit il prodotto
supererebbe il massimo rappresentabile, circa $2{,}1$ miliardi, per matrici
con più di $2^{31}$ elementi, e l'errore non verrebbe segnalato dunque il programma
leggerebbe dalla posizione sbagliata.

Il layout row-major si adatta bene a questo nucleo per tre motivi:
\begin{itemize}
  \item $A$ viene letta riga per riga, e ogni riga è un'area contigua di
        memoria: è il modo più efficiente di leggerla, sia sulla GPU sia sulla
        CPU (\cref{sec:gpu,sec:cpu});
  \item $X$ e $Y$ hanno solo $k$ colonne: in row-major i $k$ valori di una riga
        sono adiacenti, e ogni elemento di $A$ si moltiplica proprio per
        un'intera riga di $X$;
  \item il blocco di $X$ necessario a un processo è un insieme di righe
        consecutive, che in row-major formano un'unica area contigua, e lo stesso vale per il blocco di $Y$. Per questo \texttt{MPI\_Bcast} e
        \texttt{MPI\_Reduce} operano direttamente sui buffer, senza copiare
        prima i dati in un buffer temporaneo (\cref{sec:mpi}).
\end{itemize}

\paragraph{L'eccezione:} sulla GPU i kernel \texttt{warp\_row\_xt} e
\texttt{warp\_row\_smem} leggono $X$ in modo più efficiente se trasposta: per
loro si prepara una copia aggiuntiva $X^T$ di dimensione $k \times N$, cioè $X$
in column-major, prima del calcolo e fuori dal tempo misurato
(\cref{sec:gpu}). La $X$ originale, usata da MPI e dalla CPU, resta
row-major. La libreria cuBLAS, usata come riferimento, lavora invece in
column-major, ma non richiede copie: una matrice row-major, letta come
column-major, è la sua trasposta, e il prodotto si riformula di conseguenza
(\cref{sec:gpu}).

\subsection{Generazione riproducibile dei dati}\label{sec:generazione}

La traccia suggerisce di generare le matrici di test con una formula o con un
generatore pseudocasuale, a condizione che la generazione sia ripetibile.

Si è deciso di usare un generatore \emph{counter-based}~\cite{salmon2011}, in cui ogni
elemento è una funzione pura della sua posizione globale (\cref{lst:fill}):
\[
  a_{ij} = \mathrm{gen}(\mathit{seme}, i, j).
\]
In questo modo, il valore dipende solo dal seme e dagli indici globali, non dalle dimensioni
della matrice, dal numero di processi o dall'ordine delle chiamate.

\begin{lstlisting}[caption={Generazione di un blocco di $A$ a partire dagli indici globali (\texttt{matgen.c}).}, label={lst:fill}]
void fill_block(real_t *B, int ld, long row0, int nrows,
                long col0, int ncols, uint64_t seed)
{
    for (int r = 0; r < nrows; r++)
        for (int c = 0; c < ncols; c++)
            /* r, c indici locali; row0 + r, col0 + c indici globali */
            B[IDX(r, c, ld)] = (real_t)gen_value(seed, row0 + r, col0 + c);
}
\end{lstlisting}

La funzione di mescolamento è quella del generatore
SplitMix64~\cite{steele2014}, con le costanti di Stafford~\cite{stafford2011}
(\cref{lst:mix}):
prende 64 bit e li rimescola in modo che, cambiando un solo bit in ingresso,
cambino in media metà dei bit in uscita. Indici vicini, come $(i, j)$ e
$(i, j+1)$, producono così valori privi di relazioni apparenti.

\begin{lstlisting}[caption={Funzione di mescolamento SplitMix64 e calcolo del valore di un elemento (\texttt{matgen.c}).}, label={lst:mix}]
static inline uint64_t mix64(uint64_t z)
{
    z += 0x9E3779B97F4A7C15ULL;               /* "golden gamma" */
    z = (z ^ (z >> 30)) * 0xBF58476D1CE4E5B9ULL;
    z = (z ^ (z >> 27)) * 0x94D049BB133111EBULL;
    return z ^ (z >> 31);
}

double gen_value(uint64_t seed, long i, long j)
{
    /* prima seme e riga, poi la colonna: (i, j) e (j, i) differiscono */
    uint64_t h = mix64(mix64(seed ^ (uint64_t)i) + (uint64_t)j);
    /* 53 bit alti -> [0, 1) -> [-1, 1) */
    return (double)(h >> 11) * 0x1.0p-53 * 2.0 - 1.0;
}
\end{lstlisting}

I 64 bit ottenuti si convertono in un numero in $[-1, 1)$: si tengono i 53 bit
più significativi, tanti quanti ne ha la mantissa di un \texttt{double}, e li
si scala~\cite{vigna}. Valori di entrambi i segni, a media nulla, fanno sì che
le somme non crescano sistematicamente con $N$ e che compaiano cancellazioni,
che mettono alla prova la precisione del calcolo. $A$ e $X$ usano due semi
diversi.

Il valore è generato in \texttt{double} e arrotondato a \texttt{real\_t} al
momento della memorizzazione. In singola precisione il riferimento seriale
applica lo stesso arrotondamento, così che entrambi partano esattamente dagli
stessi dati.

\subsection{Padding e allineamento}\label{sec:padding}

La memoria non viene letta un byte alla volta, ma a blocchi di dimensione
fissa, allineati in memoria: sulla CPU l'unità è la \emph{linea di cache} da
\SI{64}{byte}, mentre sulla GPU l'unità minima è il \emph{settore} da
\SI{32}{byte}, e quattro settori formano una \emph{linea} da \SI{128}{byte}.
Se una riga di $A$ comincia a metà di uno di questi blocchi, la sua lettura
tocca un blocco in più del necessario. Per evitarlo, il preprocessamento fa sì
che \emph{ogni riga di $A$ inizi a un indirizzo allineato}, con due
accorgimenti.

\paragraph{Allocazione allineata.} $A$ è allocata con
\texttt{posix\_memalign}, che, a differenza di \texttt{malloc} (che garantisce
solo \SI{16}{byte}), restituisce un indirizzo multiplo dell'allineamento
richiesto: la prima riga è quindi allineata.

\paragraph{Padding delle righe.} La leading dimension \texttt{lda} è
arrotondata per eccesso al multiplo dell'allineamento: ogni riga è seguita da
alcune posizioni vuote, che non contengono dati e non vengono mai lette. Così
anche tutte le righe successive iniziano allineate.

\medskip
L'allineamento è scelto in base all'unità di lettura del backend, e quindi dipende dal tipo di eseguibile:
\begin{itemize}
  \item \emph{CPU}, \SI{64}{byte}: ogni riga inizia all'inizio di una linea di
        cache;
  \item \emph{GPU}, \SI{128}{byte}: in \texttt{double} un warp legge 32
        elementi consecutivi di una riga, cioè \SI{256}{byte}. Con la riga
        allineata tocca il minimo possibile, 8 settori e 2 linee, mentre disallineata
        ne tocca uno in più di ciascuno, leggendo circa il \SI{12}{\percent} di
        byte in più.
\end{itemize}

Sulla GPU, $A$ mantiene la stessa disposizione che ha in memoria host, con la
stessa \texttt{lda}, già arrotondata a \SI{128}{byte} perché l'allineamento
dipende dal backend compilato. Sul dispositivo si alloca spazio per $M$ righe
di \texttt{lda} elementi, padding compreso, e la copia avviene con
\texttt{cudaMemcpy2D}: questa trasferisce le righe una alla volta, copiando
solo gli $N$ elementi utili di ciascuna e saltando le posizioni di padding,
che sul dispositivo restano non inizializzate perché non vengono mai lette.
Le righe risultano allineate anche sulla GPU per due motivi:
\begin{itemize}
  \item \texttt{cudaMalloc} restituisce un indirizzo multiplo di
        \SI{256}{byte}, quindi la prima riga è allineata;
  \item il passo tra una riga e la successiva, \texttt{lda}, è multiplo di
        \SI{128}{byte}, quindi sono allineate anche tutte le righe successive.
\end{itemize}
È lo stesso ragionamento fatto sull'host, con \texttt{cudaMalloc} al posto di
\texttt{posix\_memalign}.

$X^T$, invece, non esiste sull'host: la costruisce direttamente la GPU
trasponendo $X$. Le sue righe sono lunghe $N$ come quelle di $A$, quindi
ricevono lo stesso padding, con leading dimension arrotondata a \SI{128}{byte}.

Il costo è al più di 63 o 127 byte per riga, trascurabile rispetto a righe di
decine di migliaia di byte. Il padding si può disattivare con l'opzione
\texttt{-p 0}, e il suo effetto sulle prestazioni è misurato dalla campagna~F
(\cref{sec:campagne}).

% =====================================================================
\section{Architettura del codice}\label{sec:architettura}

\subsection{Terminologia e kernel disponibili}\label{sec:terminologia}

Due termini ricorrono in tutta la relazione e conviene distinguerli subito.

Il \emph{nucleo di calcolo} è il termine della traccia: l'intera operazione
$Y \leftarrow A\,X$ sulla matrice globale. Nella versione parallela comprende
tre fasi:
\begin{itemize}
  \item la distribuzione di $X$ ai processi;
  \item il calcolo locale su ogni processo;
  \item la raccolta di $Y$.
\end{itemize}
Il suo tempo è quello usato per i GFLOPS (\cref{sec:metrica}).

Un \emph{kernel} è invece una specifica implementazione del \emph{prodotto
locale}, cioè del calcolo che ogni processo esegue sul proprio blocco di $A$.
Il termine viene da CUDA, dove indica una funzione eseguita sulla GPU, e nel progetto lo si usa per estensione anche per le varianti della CPU, che
svolgono lo stesso ruolo.

Il backend è fissato in compilazione (\cref{sec:compilazione})
il kernel si sceglie invece a tempo di esecuzione con
l'opzione \texttt{-v}, tra quelli del backend
(\cref{tab:kernel}). Senza \texttt{-v} si usa il kernel di default, il più
veloce. Un valore non valido termina il programma con un errore, prima di
allocare la memoria del backend.

\begin{table}[H]
  \centering
  \caption{Kernel disponibili per ciascun backend, con il valore da passare
    all'opzione \texttt{-v}.}\label{tab:kernel}
  \begin{tabular}{l p{0.75\textwidth}}
    \toprule
    Backend & Kernel \\
    \midrule
    CUDA & \texttt{thread\_row} (0), \texttt{warp\_row} (1),
           \texttt{warp\_row\_xt} (2, default), \texttt{warp\_row\_smem} (3),
           \texttt{cublas} ($-1$, solo riferimento) \\
    CPU  & \texttt{cpu\_row} (0), \texttt{cpu\_block} (1, default) \\
    \bottomrule
  \end{tabular}
\end{table}

Nel backend CUDA il file di ciascun kernel (\cref{sec:struttura}) contiene
anche una piccola funzione \emph{lanciatore}, che sceglie le dimensioni della
griglia e dei blocchi di thread e avvia il kernel, mentre \texttt{backend\_cuda.cu}
gestisce memoria, tempi e scelta del kernel, e chiama solo i lanciatori.

In sintesi, il nucleo è sempre lo stesso, mentre i kernel sono i diversi modi
di eseguirne la parte locale. Poiché la scelta avviene a tempo di esecuzione,
i kernel si confrontano con lo stesso eseguibile e sugli stessi dati,
cambiando un solo parametro: è così che sono condotti i confronti delle
\cref{sec:gpu,sec:cpu} e le campagne di misura (\cref{sec:campagne}).

\subsection{Interfaccia dei backend}\label{sec:interfaccia}

Come già introdotto, il codice è diviso in due parti con responsabilità
distinte:
\begin{itemize}
  \item la \emph{parte MPI} (\texttt{mxm.c}) decide \emph{quali dati} ha ogni
        processo e li fa comunicare;
  \item il \emph{backend} decide \emph{come calcolare} il prodotto locale del
        proprio blocco.
\end{itemize}
Tra le due c'è un'interfaccia, definita in \texttt{backend.h},
(\cref{tab:interfaccia}).

\begin{table}[H]
  \centering
  \caption{Funzioni principali dell'interfaccia dei backend.}\label{tab:interfaccia}
  \begin{tabular}{l l p{0.36\textwidth} p{0.16\textwidth}}
    \toprule
    Funzione & Quando & Cosa fa & Nel tempo \\
    \midrule
    \texttt{be\_create}  & una volta     & prepara $A$ (GPU: copia sul dispositivo)
                         & no (preprocessamento) \\
    \texttt{be\_set\_x}  & ogni invocazione & rende $X$ disponibile al calcolo
                           (GPU: copia ed eventuale trasposizione) & no \\
    \texttt{be\_compute} & ogni invocazione & calcola $Y = A\,X$ e ritorna a
                           calcolo terminato & \textbf{sì} \\
    \texttt{be\_get\_y}  & ogni invocazione & rende $Y$ disponibile a MPI
                           (GPU: copia verso la CPU) & no \\
    \texttt{be\_destroy} & una volta     & libera la memoria & --- \\
    \bottomrule
  \end{tabular}
\end{table}

Questa divisione delle funzioni corrisponde proprio alle fasi da misurare
(\cref{sec:metrica}):
\begin{itemize}
  \item \texttt{mxm.c} misura il calcolo solo attorno a \texttt{be\_compute},
        mentre i trasferimenti (\texttt{be\_set\_x} e \texttt{be\_get\_y}) sono
        cronometrati a parte;
  \item sul backend CPU i dati si trovano già nella memoria della CPU, quindi
        \texttt{be\_set\_x} e \texttt{be\_get\_y} costano poco o nulla.
\end{itemize}
Completano l'interfaccia alcune funzioni di servizio: l'allineamento preferito
per le righe di $A$ (\cref{sec:padding}), la scelta della GPU
(\cref{sec:gpu-assign}), il nome e l'elenco dei kernel e il kernel di default.

\paragraph{Tipo opaco.} \texttt{mxm.c} usa il backend solo tramite un puntatore
a \texttt{backend\_t}, dichiarato ma non definito in \texttt{backend.h}: i
campi (memoria della GPU, eventi CUDA, kernel scelto) restano nascosti
nell'implementazione, che si può quindi modificare senza toccare la parte MPI.

\paragraph{\texttt{extern "C"}.} Il backend CUDA è compilato da \texttt{nvcc}
come C++, che per gestire l'overloading aggiunge al nome di ogni funzione i
tipi dei parametri (\emph{name mangling}): \texttt{be\_create} diventa, ad
esempio, \texttt{\_Z9be\_createiiiPKdiii}. Il C invece usa il nome così
com'è, quindi \texttt{mxm.c} non troverebbe la funzione al collegamento.
\texttt{extern "C"} fa usare al compilatore C++ i nomi in stile C.

\subsection{Precisione}\label{sec:precisione}

Il programma lavora in doppia o in singola precisione, con scelta effettuata in
compilazione (\texttt{make PREC=double|float}). Tutto il codice usa un unico
tipo numerico, \texttt{real\_t}, definito in \texttt{common.h}: senza opzioni
è \texttt{double}, mentre con \texttt{PREC=float} il \texttt{Makefile} definisce la
macro \texttt{USE\_FLOAT} e \texttt{real\_t} diventa \texttt{float}. Dalla
stessa scelta discendono le costanti collegate:
\begin{itemize}
  \item \texttt{REAL\_EPS}, l'\emph{epsilon di macchina} della precisione
        scelta, usato dalla tolleranza della verifica (\cref{sec:verifica});
  \item il tipo MPI corrispondente, \texttt{MPI\_DOUBLE} o
        \texttt{MPI\_FLOAT}, usato nelle comunicazioni;
  \item il nome della precisione, riportato nell'output delle misure.
\end{itemize}
Scegliendo la precisione in compilazione, ogni kernel è compilato per un solo
tipo, senza scelte da fare durante il calcolo.

Alcune parti restano sempre in doppia precisione, indipendentemente da
\texttt{real\_t}: la \emph{generazione} dei valori, poi arrotondati a
\texttt{real\_t} al momento della memorizzazione (\cref{sec:generazione}), e
gli \emph{accumulatori del riferimento seriale}, che deve essere più accurato
del codice che verifica (\cref{sec:collaudo}).

Passare a \texttt{float} dimezza i byte per elemento: a parità di dati letti
raddoppia l'intensità, cioè le operazioni per byte. Insieme al
picco di calcolo 32 volte più alto della GPU in singola precisione
(\cref{sec:gpu-hw}), questo spiega perché in \texttt{float} la GPU sia molto
più veloce a $k$ piccolo. Il confronto tra le due precisioni, velocità contro
errore, è discusso in \cref{sec:misure}.

% =====================================================================
\section{Parallelizzazione con MPI}\label{sec:mpi}

\subsection{Concetti di base}\label{sec:mpi-base}

Un programma MPI è eseguito da più processi indipendenti, ciascuno con la
propria memoria, che condividono dati solo scambiando messaggi. I processi
sono raggruppati in \emph{comunicatori}, all'interno dei quali ognuno ha un
identificativo, il \emph{rank}. All'avvio esiste solo
\texttt{MPI\_COMM\_WORLD}, che li contiene tutti, e il programma può crearne
altri su sottoinsiemi di processi.

Le comunicazioni usate nel progetto sono tutte \emph{collettive}: coinvolgono
tutti i processi di un comunicatore, che chiamano insieme la stessa funzione.
Il progetto ne usa due:
\begin{itemize}
  \item \texttt{MPI\_Bcast} (\emph{broadcast}): un processo, la \emph{radice},
        invia gli stessi dati a tutti gli altri del proprio comunicatore;
  \item \texttt{MPI\_Reduce}: ogni processo contribuisce con un proprio
        vettore e la radice riceve la loro combinazione elemento per elemento,
        in questo progetto la somma.
\end{itemize}
Il centro del ragionamento è \emph{definire i comunicatori} e scegliere \emph{su quali comunicatori} eseguire queste due
operazioni, così che ogni messaggio raggiunga solo i processi che ne hanno
bisogno (\cref{sec:griglia,sec:tre-passi}).

\subsection{Griglia cartesiana dei processi}\label{sec:griglia}

Come richiesto dalla traccia, i processi sono disposti su una griglia
bidimensionale $P_r \times P_c$, con dimensioni scelte dall'utente tramite le
opzioni \texttt{-r} e \texttt{-c}. Il programma verifica anche che $P_r \cdot P_c$
coincida con il numero di processi lanciati e, in caso contrario, termina con
un errore.

La griglia è costruita in \texttt{grid\_create} con le \emph{topologie
cartesiane} di MPI:
\begin{itemize}
  \item \texttt{MPI\_Cart\_create} crea un nuovo comunicatore in cui ogni
        processo ha una posizione nella griglia. La griglia è \emph{non
        periodica} (\texttt{periods = 0}), poiché non servono comunicazioni
        tra bordi opposti, e \emph{senza riordino} (\texttt{reorder = 0}): MPI
        non rinumera i processi, quindi il rank resta quello assegnato da
        \texttt{mpirun} e il legame tra processi e core rimane prevedibile
        (\cref{sec:cpu});
  \item \texttt{MPI\_Cart\_coords} ricava le coordinate $(r, c)$ di ogni
        processo, in ordine row-major: $\mathit{rank} = r \cdot P_c + c$;
  \item \texttt{MPI\_Cart\_sub} estrae due famiglie di sotto-comunicatori:
        quelli di \emph{riga} (processi con lo stesso $r$, rank interno pari a
        $c$), usati per la riduzione di $Y$, e quelli di \emph{colonna}
        (processi con lo stesso $c$, rank interno pari a $r$), usati per il
        broadcast di $X$.
\end{itemize}

Le chiamate corrispondenti in \texttt{grid\_create} sono riportate nel
\cref{lst:grid}.

\begin{lstlisting}[caption={Costruzione della griglia e dei comunicatori di riga e di colonna (\texttt{distrib.c}).}, label={lst:grid}]
int dims[2]    = { Pr, Pc };
int periods[2] = { 0, 0 };
MPI_Cart_create(comm, 2, dims, periods, 0, &g->grid);
/* ... */
int coords[2];
MPI_Cart_coords(g->grid, g->rank, 2, coords);
g->myr = coords[0];
g->myc = coords[1];
/* ... */
int keep_cols[2] = { 0, 1 };
int keep_rows[2] = { 1, 0 };
MPI_Cart_sub(g->grid, keep_cols, &g->row);
MPI_Cart_sub(g->grid, keep_rows, &g->col);
\end{lstlisting}

In alternativa si potrebbe usare \texttt{MPI\_Comm\_split}: ogni processo
indica un numero, il proprio \emph{color}, e i processi con lo stesso colore
finiscono nello stesso nuovo comunicatore, usando un colore per riga e uno
per colonna.
Con le topologie cartesiane, invece, è MPI a ricavare coordinate e
gruppi a partire dalla forma della griglia.

\texttt{grid\_create} riceve come parametro il comunicatore su cui costruire
la griglia, invece di usare direttamente \texttt{MPI\_COMM\_WORLD}: la funzione
è così riutilizzabile anche in un'applicazione più grande, che dedichi al
prodotto solo una parte dei propri processi.

\subsection{Distribuzione dei dati}\label{sec:distribuzione}

\paragraph{Matrice $A$.} È divisa in blocchi secondo la griglia: le $M$ righe
sono ripartite in $P_r$ gruppi, le $N$ colonne in $P_c$ gruppi, e il processo
$(r, c)$ possiede il blocco $A_{rc}$, formato dalle righe del gruppo $r$ e
dalle colonne del gruppo $c$.

La ripartizione è \emph{bilanciata}: dividendo $n$ elementi in $p$ parti,
ognuna riceve $\lfloor n/p \rfloor$ elementi e gli $n \bmod p$ elementi
avanzati vanno uno ciascuno alle prime parti. Le parti differiscono quindi al
più di un elemento,
 Due funzioni forniscono, per la parte $i$-esima, la dimensione
(\texttt{blk\_size}) e l'indice globale del primo elemento
(\texttt{blk\_off}): ad esempio, 10 righe su 3 processi diventano
$4 + 3 + 3$, con prime righe nelle posizioni 0, 4 e 7.

\paragraph{$A$ non si sposta mai.} Ogni processo genera da solo il proprio
blocco a partire dagli indici globali (\cref{sec:generazione}): non è necessario
un processo che costruisce la matrice e la distribuisce, e $A$, la struttura più
grande, non viene mai comunicata.

\paragraph{Multivettori $X$ e $Y$.} Si dividono seguendo il prodotto:
\begin{itemize}
  \item $X$ come le \emph{colonne} di $A$: il processo $(r, c)$ ha bisogno solo
        delle righe di $X$ corrispondenti alle colonne del suo blocco, cioè di
        $X_c$; $X$ è quindi divisa in $P_c$ parti;
  \item $Y$ come le \emph{righe} di $A$: il blocco $A_{rc}$ contribuisce solo
        alle righe di $Y$ del gruppo $r$, cioè a $Y_r$: $Y$ è quindi divisa in
        $P_r$ parti.
\end{itemize}
La traccia consente di assumere che $X$ e $Y$ occupino una sola riga o colonna
della griglia. Si è scelto che $X$ stia sulla riga~0, dove il processo
$(0, c)$ possiede $X_c$, e che $Y$ si trovi alla fine sulla colonna~0, dove il
processo $(r, 0)$ possiede $Y_r$ (\cref{fig:distribuzione}). Come visto in
\cref{sec:layout}, in row-major $X_c$ e $Y_r$ sono aree contigue, quindi le
comunicazioni operano direttamente su di esse.

\begin{figure}[H]
  \centering
  \begin{tikzpicture}[
    blk/.style={draw, minimum width=2.2cm, minimum height=1.2cm, fill=gray!8},
    xv/.style={draw, minimum width=2.2cm, minimum height=0.55cm, fill=blue!12},
    yv/.style={draw, minimum width=0.7cm, minimum height=1.2cm, fill=orange!18},
    lab/.style={font=\footnotesize, text=black!70}]
    \foreach \r in {0,1} \foreach \c in {0,1,2}
      \node[blk] (A\r\c) at (2.2*\c, -1.2*\r) {$A_{\r\c}$};
    \foreach \c in {0,1,2}
      \node[xv, above=0.35cm of A0\c] (X\c) {$X_{\c}$};
    \foreach \r in {0,1}
      \node[yv, left=0.35cm of A\r0] (Y\r) {$Y_{\r}$};
    \foreach \c in {0,1,2}
      \node[lab, above=0.05cm of X\c] {$c=\c$};
    \foreach \r in {0,1}
      \node[lab, left=0.05cm of Y\r] {$r=\r$};
    \node[lab, below=0.35cm of A11, align=center]
      {processo $(r,c)$: blocco $A_{rc}$ \quad
       $X_c$ sul processo $(0,c)$ \quad
       $Y_r$ (alla fine) sul processo $(r,0)$};
  \end{tikzpicture}
  \caption{Distribuzione dei dati su una griglia $2 \times 3$: il processo
    $(r, c)$ possiede il blocco $A_{rc}$, $X_c$ risiede sul processo $(0, c)$
    della riga~0, $Y_r$ alla fine si trova sul processo $(r, 0)$ della
    colonna~0.}\label{fig:distribuzione}
\end{figure}

\subsection{Schema di calcolo in tre passi}\label{sec:tre-passi}

Con i dati distribuiti come in \cref{fig:distribuzione}, ogni invocazione del
nucleo si svolge in tre passi (\cref{fig:tre-passi}).

\paragraph{Passo 1: broadcast di $X$.} Su ogni colonna $c$ della griglia, il
processo della riga~0 invia $X_c$ a tutti i processi della stessa colonna con
\texttt{MPI\_Bcast}, sul comunicatore di colonna e con radice la riga~0. Dopo
questo passo ogni processo $(r, c)$ dispone del pezzo di $X$ che gli serve.

\paragraph{Passo 2: calcolo locale.} Ogni processo calcola sul proprio backend
il prodotto parziale
\[
  Y_{rc} = A_{rc}\, X_c ,
\]
una matrice di dimensione $M_r \times k$, dove $M_r \approx M / P_r$ è il
numero di righe del gruppo $r$ (\cref{sec:distribuzione}). È l'unico passo
che svolge operazioni in virgola mobile, ed è eseguito da uno dei kernel delle
\cref{sec:gpu,sec:cpu}.

\paragraph{Passo 3: riduzione di $Y$.} Il blocco $A_{rc}$ contiene solo una
parte delle colonne di $A$, quindi $Y_{rc}$ è solo un \emph{contributo} alle
righe di $Y$ del gruppo $r$. Il risultato completo è la somma dei contributi di
tutti i blocchi della riga:
\[
  Y_r = \sum_{c=0}^{P_c - 1} A_{rc}\, X_c .
\]
La somma è eseguita da \texttt{MPI\_Reduce} con operazione \texttt{MPI\_SUM},
sul comunicatore di riga e con radice la colonna~0: alla fine il processo
$(r, 0)$ possiede $Y_r$.

Nel codice i tre passi sono la funzione \texttt{mxm\_apply} di
\texttt{mxm.c} (\cref{lst:tre-passi}).

\begin{lstlisting}[caption={I tre passi di un'invocazione del nucleo (\texttt{mxm.c}, modalità a turni omessa).}, label={lst:tre-passi}]
/* Prodotto locale Y_rc = A_rc * X_c, con i tempi delle sue tre fasi. */
static void local_product(mxm_t *h, const real_t *Xc, double *t)
{
    const double t0 = MPI_Wtime();
    be_set_x(h->be, Xc);
    const double t1 = MPI_Wtime();
    t[MXM_T_COMPUTE_DEV] = be_compute(h->be);   /* tempo del device */
    const double t2 = MPI_Wtime();
    be_get_y(h->be, h->Ypart);
    const double t3 = MPI_Wtime();

    t[MXM_T_H2D]     = t1 - t0;
    t[MXM_T_COMPUTE] = t2 - t1;
    t[MXM_T_D2H]     = t3 - t2;
}

void mxm_apply(mxm_t *h, const real_t *X, real_t *Y, double *times)
{
    /* riga 0: X_c del chiamante; altre righe: X_c ricevuta col bcast */
    real_t *Xc = (g->myr == 0) ? (real_t *)X : h->Xrecv;

    const double t0 = MPI_Wtime();
    /* passo 1: broadcast di X_c lungo la colonna (root = riga 0) */
    MPI_Bcast(Xc, h->Nloc * h->k, MPI_REAL_T, 0, g->col);
    const double t1 = MPI_Wtime();
    /* passo 2: prodotto locale */
    local_product(h, Xc, t);
    const double t2 = MPI_Wtime();
    /* passo 3: somma di Y_rc lungo la riga (root = colonna 0) */
    MPI_Reduce(h->Ypart, Y, h->Mloc * h->k, MPI_REAL_T, MPI_SUM,
               0, g->row);
    const double t3 = MPI_Wtime();
    /* ... */
    t[MXM_T_CORE] = t[MXM_T_BCAST] + t[MXM_T_COMPUTE] + t[MXM_T_REDUCE];
}
\end{lstlisting}

Ogni broadcast coinvolge solo i $P_r$ processi di una colonna e ogni riduzione
solo i $P_c$ processi di una riga: le operazioni sulle diverse colonne, e
quelle sulle diverse righe, avvengono quindi in parallelo tra loro, ognuna sul
proprio comunicatore. Il tempo del nucleo è la somma dei tre passi
(\cref{sec:metrica}).

\begin{figure}[H]
  \centering
  \begin{tikzpicture}[
    blk/.style={draw, minimum width=2.6cm, minimum height=1.4cm, fill=gray!8, align=center},
    xv/.style={draw, minimum width=2.6cm, minimum height=0.55cm, fill=blue!12},
    yv/.style={draw, minimum width=0.7cm, minimum height=1.4cm, fill=orange!18},
    bc/.style={-{Stealth[length=2.2mm]}, thick, blue!60!black},
    rd/.style={-{Stealth[length=2.2mm]}, thick, orange!70!black},
    lab/.style={font=\footnotesize, text=black!70}]
    \foreach \r in {0,1} \foreach \c in {0,1,2}
      \node[blk] (A\r\c) at (3.6*\c, -2.0*\r)
        {$A_{\r\c}$\\[1pt]{\footnotesize $Y_{\r\c} = A_{\r\c}X_{\c}$}};
    \foreach \c in {0,1,2}
      \node[xv, above=0.6cm of A0\c] (X\c) {$X_{\c}$};
    \foreach \r in {0,1}
      \node[yv, left=0.9cm of A\r0] (Y\r) {$Y_{\r}$};
    \foreach \c in {0,1,2} {
      \draw[bc] ([xshift=-0.45cm]X\c.south) -- ([xshift=-0.45cm]A0\c.north);
      \draw[bc] ([xshift=-0.45cm]A0\c.south) -- ([xshift=-0.45cm]A1\c.north);
    }
    \foreach \r in {0,1} {
      \draw[rd] (A\r2.west) -- (A\r1.east);
      \draw[rd] (A\r1.west) -- (A\r0.east);
      \draw[rd] (A\r0.west) -- (Y\r.east) node[midway, above, font=\scriptsize] {$\Sigma$};
    }
    \node[lab, below=0.45cm of A11, align=left]
      {\textcolor{blue!60!black}{\textbf{1}}~\texttt{MPI\_Bcast} di $X_c$ lungo la colonna (radice: riga 0)\\
       \textbf{2}~calcolo locale del prodotto parziale\\
       \textcolor{orange!70!black}{\textbf{3}}~\texttt{MPI\_Reduce} (somma) lungo la riga (radice: colonna 0)};
  \end{tikzpicture}
  \caption{I tre passi di un'invocazione del nucleo su una griglia
    $2 \times 3$. Le frecce indicano il verso in cui si spostano i dati, non
    l'ordine dei messaggi: le collettive sono realizzate da MPI ad albero
    (\cref{sec:costo-comm}).}\label{fig:tre-passi}
\end{figure}

\subsection{Costo delle comunicazioni}\label{sec:costo-comm}

Il calcolo locale di ogni processo non dipende dagli altri: l'unico costo che la
parallelizzazione aggiunge è la comunicazione, cioè il broadcast di $X$ e la
riduzione di $Y$ (\cref{sec:tre-passi}). Quanto pesa dipende da come si divide
$A$, ed è questo che spiega la scelta di una griglia bidimensionale.

\paragraph{Che cosa dividere.} $A$ ha $MN$ elementi, mentre $X$ e $Y$ ne hanno
solo $Nk$ e $Mk$: con $k$ piccolo, $A$ è di gran lunga la matrice più grande.
Copiarla su ogni processo sarebbe la scelta peggiore, per due ragioni
indipendenti: ogni processo dovrebbe contenerla per intero, e la rileggerebbe
tutta, mentre è proprio la lettura di $A$ a determinare il tempo
(\cref{sec:roofline}). \textbf{$A$ va quindi divisa}, mentre $X$ e $Y$, piccole, possono
invece essere copiate dove servono.

\paragraph{Le due divisioni a una dimensione.} Le divisioni più semplici tagliano
$A$ lungo una sola direzione. Nello schema in tre passi (\cref{sec:tre-passi})
corrispondono alle due griglie estreme, in cui resta una sola delle due
comunicazioni:
\begin{itemize}
  \item \emph{per righe} (griglia $P \times 1$): ogni processo possiede righe
        intere di $A$ e calcola righe di $Y$ già definitive, senza somme finali, ma in cambio gli serve tutta $X$, che va inviata a tutti i $P$
        processi;
  \item \emph{per colonne} (griglia $1 \times P$): a ogni processo basta il
        proprio pezzo di $X$, ma il suo risultato è parziale e riguarda tutte
        le righe di $Y$, quindi va sommato con quello di tutti gli altri
        processi.
\end{itemize}
Le due soluzioni sono speculari e hanno lo stesso difetto: \textbf{ogni
comunicazione coinvolge tutti i processi, su un oggetto grande quanto l'intera
$X$ o l'intera $Y$}. Aggiungere processi aumenta il numero dei destinatari, ma
non riduce la quantità di dati che ciascuno deve ricevere o inviare.

\paragraph{La divisione a griglia.} Con una griglia $P_r \times P_c$ ogni
processo possiede un rettangolo di $A$. Servono sia il broadcast sia la
riduzione, ma ciascuno riguarda solo una parte dei processi: il broadcast i
$P_r$ processi di una colonna, la riduzione i $P_c$ processi di una riga. Anche
i messaggi si accorciano: ogni processo riceve $(N/P_c)\,k$ elementi di $X$ e
invia $(M/P_r)\,k$ elementi di $Y$, invece di $Nk$ o $Mk$.

\paragraph{Volume.} Sommando su tutta la griglia:
\begin{itemize}
  \item nel \emph{broadcast}, ogni colonna invia $X_c$ ai suoi $P_r - 1$
        processi che non lo possiedono: in tutto si muovono
        $(P_r - 1)\, N k$ elementi;
  \item nella \emph{riduzione}, ogni riga riceve i contributi dei suoi
        $P_c - 1$ processi fuori dalla colonna~0: in tutto si muovono
        $(P_c - 1)\, M k$ elementi.
\end{itemize}
Il volume totale è quindi
\[
  V = (P_r - 1)\, N k + (P_c - 1)\, M k .
\]
Le due divisioni a una dimensione sono i casi estremi: con $P_c = 1$ non c'è
riduzione e $V = (P - 1)\, N k$, con $P_r = 1$ non c'è broadcast e
$V = (P - 1)\, M k$. Una griglia quadrata ripartisce invece il volume tra le
due comunicazioni: con 16 processi su una matrice quadrata, la griglia
$4 \times 4$ muove $6\, N k$ elementi, contro i $15\, N k$ della $16 \times 1$,
cioè 2,5 volte meno.

A parità di numero di processi, la griglia che muove meno dati dipende però
dalla forma della matrice. \Cref{tab:volume} riporta il volume con 4 processi, in
unità di $N k$:
\begin{itemize}
  \item matrice quadrata: conviene la $2 \times 2$, che ripartisce il volume
        tra broadcast e riduzione;
  \item matrice alta: conviene la $4 \times 1$, che evita di ridurre una $Y$
        lunga;
  \item matrice larga: la $2 \times 2$ e la $1 \times 4$ si equivalgono.
\end{itemize}

\begin{table}[H]
  \centering
  \caption{Volume delle comunicazioni con 4 processi, in unità di $N k$.}\label{tab:volume}
  \begin{tabular}{lccc}
    \toprule
    Griglia & $M = N$ & $M = 3N$ & $N = 2M$ \\
    \midrule
    $4 \times 1$ & 3          & \textbf{3} & 3            \\
    $2 \times 2$ & \textbf{2} & 4          & \textbf{1,5} \\
    $1 \times 4$ & 3          & 9          & \textbf{1,5} \\
    \bottomrule
  \end{tabular}
\end{table}

\paragraph{Dal volume al tempo.} Il volume totale non è ancora il tempo. MPI non
invia i dati a un processo per volta, cosa che richiederebbe $p - 1$ passi
successivi: usa algoritmi ad albero, in cui a ogni passo il numero di processi
coinvolti raddoppia, e bastano circa $\log_2 p$ passi, con $p$ pari a $P_r$ nel
broadcast e a $P_c$ nella riduzione. Il tempo dipende quindi soprattutto da due
fattori: \textbf{la dimensione del messaggio di ogni processo e il numero di
processi del gruppo}. La griglia li riduce entrambi rispetto alle divisioni a una
dimensione. A parità di messaggio, inoltre, la riduzione costa più del
broadcast, in parte perché oltre a spostare i dati deve sommarli: l'effetto è
misurato in \cref{sec:ris-mpi}.

\paragraph{Perché non si divide lungo $k$.} $X$ e $Y$ si dividono solo per
righe, e ogni processo ha sempre tutte le $k$ colonne. Dividerle lungo $k$
costringerebbe a duplicare $A$, perché ogni colonna di $Y$ richiede righe
intere di $A$. Ridurrebbe inoltre il riuso di $A$ nei kernel. Infine
produrrebbe messaggi di pochi valori, dominati dalla latenza.

\subsection{Assegnazione delle GPU ai processi}\label{sec:gpu-assign}

Nella versione CUDA ogni processo deve scegliere quale GPU usare. Su una
macchina con più nodi, ciascuno con una o più GPU, i processi dello stesso
nodo vanno distribuiti sulle GPU di quel nodo, e il rank globale non basta,
perché non dice quali processi condividono il nodo.

Per questo si usa \texttt{MPI\_Comm\_split\_type} con il tipo
\texttt{MPI\_COMM\_TYPE\_SHARED}, che raggruppa i processi che condividono la
memoria, cioè lo stesso nodo. Il rank di un processo in questo comunicatore è
il suo \emph{rank locale} sul nodo, e il backend sceglie la GPU con
\[
  \text{GPU} = \text{rank locale} \bmod \text{(numero di GPU del nodo)},
\]
così che i processi di un nodo si distribuiscano a turno sulle sue GPU.

Sul server c'è un solo nodo con una sola GPU: il risultato dell'operazione è
sempre 0, e tutti i processi usano la stessa GPU, contendendosela
(\cref{sec:misure}). Su questa macchina il meccanismo non ha quindi effetto:
serve però alla \emph{portabilità}. Il codice è pensato per usare correttamente
più GPU, su uno o più nodi, senza modifiche, e quindi per poter scalare su una
macchina diversa. Lo scenario con una GPU per processo, qui non disponibile, è invece simulato
dalla modalità a nodi emulati (\cref{sec:misure}).

% =====================================================================
\section{Implementazione su GPU (CUDA)}\label{sec:gpu}

\paragraph{Come sono misurati i kernel.} Tutti i kernel di questo capitolo sono
stati misurati nelle stesse condizioni: matrice quadrata $16384 \times 16384$,
un solo processo, i cinque valori di $k$ della traccia (campagna~B,
\cref{sec:gruppi}). I risultati sono riportati alla fine della sezione di
ciascun kernel e riassunti in \cref{sec:gpu-riepilogo}, dove ogni kernel è
confrontato con i tetti della roofline (\cref{sec:roofline}). I GFLOPS sono
calcolati con il tempo del nucleo, la metrica richiesta dalla traccia
(\cref{sec:metrica}). Con un solo processo broadcast e riduzione non
comunicano nulla, e il tempo del nucleo supera quello del solo calcolo
dell'\SIrange{1}{3}{\percent} nella maggior parte dei casi. La differenza
arriva all'\SI{8}{\percent} solo per cuBLAS in \texttt{float} a $k = 32$,
dove il calcolo dura appena \SI{4}{ms} e la riduzione, che con un solo
processo si limita a copiare $Y$, pesa di più.

\subsection{Modello di programmazione CUDA}\label{sec:cuda-modello}

In CUDA la CPU (\emph{host}) e la GPU (\emph{device}) hanno memorie separate:
i dati vanno copiati esplicitamente dall'una all'altra, e un puntatore alla
memoria della GPU non si può usare dalla CPU.

Un \emph{kernel} è una funzione eseguita sulla GPU da molti thread in
parallelo. I thread sono raggruppati in \emph{blocchi}, fino a 1024 thread
ciascuno, e i blocchi formano la \emph{griglia}. La GPU assegna ogni blocco a
uno SM, che lo esegue a gruppi di 32 thread, i \emph{warp}: i 32 thread di un
warp eseguono la stessa istruzione nello stesso momento, ciascuno sui propri
dati. Per questo il warp, e non il singolo thread, è l'unità su cui ragionare
per gli accessi in memoria (\cref{sec:coalescenza}) e per la suddivisione del
lavoro (\crefrange{sec:thread-row}{sec:warp-smem}).

Il \emph{lancio} di un kernel è \emph{asincrono}: la CPU prosegue subito,
mentre la GPU lavora, e per sapere quando il kernel è terminato serve una
sincronizzazione esplicita. Ne seguono due conseguenze per il progetto:
\begin{itemize}
  \item \texttt{be\_compute} attende la fine del kernel prima di ritornare,
        così che i dati siano pronti per la riduzione MPI;
  \item un cronometro sulla CPU attorno al solo lancio misurerebbe il tempo di
        accodamento e non quello di calcolo: per questo il tempo sulla GPU si
        misura con gli \emph{eventi CUDA} (\cref{sec:h2d}).
\end{itemize}

\subsection{Accessi alla memoria globale}\label{sec:coalescenza}

Quando un warp esegue una lettura dalla memoria globale, i 32 thread chiedono
ciascuno un indirizzo. L'hardware non serve le richieste una per una: le
raggruppa per \emph{settori} da \SI{32}{byte} e legge ogni settore toccato una
sola volta. Il costo di una lettura dipende quindi da quanti settori distinti
il warp tocca, non solo da quanti byte servono davvero.

Si parla di \emph{coalescenza} quando le richieste dei 32 thread di un warp si
raggruppano in pochi settori, cosa che dipende da quanto sono vicini gli
indirizzi letti:

\paragraph{Accesso coalescente:} Se i 32 thread leggono elementi consecutivi,
ad esempio 32 \texttt{double} contigui di una riga di $A$, gli indirizzi coprono
\SI{256}{byte}, cioè 8 settori: è il minimo possibile, e ogni byte letto viene
usato. Se la riga è allineata (\cref{sec:padding}), i \SI{256}{byte} cadono
anche in esattamente 2 \emph{linee} da \SI{128}{byte}, l'unità con cui la cache
L1 conserva i dati.

\paragraph{Accesso non coalescente:} Se i 32 thread leggono elementi distanti
tra loro, ad esempio lo stesso elemento $j$ di 32 righe diverse di $A$, ogni
thread tocca un settore diverso: il warp legge 32 settori, cioè
\SI{1024}{byte}, per usarne 256, con un'efficienza del \SI{25}{\percent} in
\texttt{double} e del \SI{12.5}{\percent} in \texttt{float}. Poiché a $k$
piccolo il nucleo è limitato dalla banda, leggere quattro volte i byte
necessari significa andare fino a quattro volte più lenti.

Il criterio che guida la suddivisione del lavoro tra i thread
(\crefrange{sec:thread-row}{sec:warp-smem}) è quindi che i thread di uno
stesso warp leggano elementi adiacenti per garantire la coalescenza.

\subsection{Versione di partenza: \texttt{thread\_row}}\label{sec:thread-row}

Il primo kernel è la traduzione più diretta del prodotto: un thread per ogni
riga di $Y$. Il thread $i$ calcola da solo l'intera riga $i$ di $Y$, con $k$
prodotti scalari tra la riga $i$ di $A$ e le colonne di $X$. I thread sono
raggruppati in blocchi da 128, e la griglia ha tanti blocchi quanti ne servono
a coprire le $M$ righe.

Il kernel è volutamente semplice e fa da riferimento per misurare il guadagno
delle versioni successive. Presenta tre difetti, che le versioni seguenti
correggono uno alla volta:
\begin{itemize}
  \item \emph{$A$ riletta $k$ volte}: per ogni colonna $c$ di $Y$ il thread
        scorre di nuovo l'intera riga $i$ di $A$, quindi $A$, la matrice
        grande, attraversa la memoria $k$ volte invece di una;
  \item \emph{accessi ad $A$ non coalescenti}: nello stesso istante i 32
        thread di un warp leggono l'elemento $j$ di 32 righe diverse, cioè il caso non coalescente di \cref{sec:coalescenza}, con il \SI{25}{\percent}
        dei byte letti effettivamente usati in \texttt{double};
  \item \emph{poco parallelismo con $M$ piccolo}: i thread sono tanti quante le
        righe del blocco locale, e con matrici basse, o con molti processi tra
        cui le righe si dividono, i thread sono pochi e la GPU resta in gran
        parte inattiva.
\end{itemize}

\paragraph{Risultato.} \texttt{thread\_row} si ferma intorno a 37~GFLOPS,
qualunque sia $k$ (\cref{tab:ris-kernel}). Non migliora con $k$ perché rilegge
la matrice $A$ una volta per ogni colonna di $X$: più cresce $k$, più lavoro
fa, ma anche più dati legge, e il rapporto tra i due non cambia. In più le sue
letture non sono coalescenti.

\subsection{Un warp per riga: \texttt{warp\_row}}\label{sec:warp-row}

Il secondo kernel assegna un warp intero, cioè 32 thread, a ogni riga di $A$, e ogni blocco contiene 8 warp, quindi 256 thread e 8 righe. Il kernel corregge i
tre difetti di \texttt{thread\_row}.

\paragraph{Lettura coalescente di $A$.} I 32 thread del warp, detti
\emph{lane} (da 0 a 31), si dividono la riga: al giro $t$ la lane $l$ legge
l'elemento $l + 32t$. Nello stesso istante il warp legge quindi 32 elementi
consecutivi della riga, cioè l'accesso coalescente di \cref{sec:coalescenza}.

\paragraph{Riuso di $A$.} Ogni elemento di $A$ letto viene messo in un
registro e usato subito per tutte le $k$ colonne: la lane moltiplica $a_{ij}$
per la riga $j$ di $X$ e aggiorna $k$ \emph{accumulatori}, uno per colonna di
$Y$, tenuti nei registri. $A$ viene così letta una sola volta, invece di $k$
volte.

\paragraph{Riduzione nel warp.} Al termine ogni lane ha, per ciascuna colonna,
solo una somma parziale, relativa agli elementi della riga che ha letto, e le
32 somme parziali vanno sommate tra loro. Lo fa l'istruzione
\texttt{\_\_shfl\_down\_sync}, che permette a una lane di leggere un registro
di un'altra lane dello stesso warp senza passare dalla memoria. La somma
procede ad albero in 5 passi: a ogni passo la lane $l$ aggiunge il valore della
lane $l + 16$, poi $l + 8$, $4$, $2$, $1$ (\cref{fig:shfl}). Alla fine la lane~0
contiene il totale e scrive la riga di $Y$.

\begin{figure}[H]
  \centering
  \resizebox{\textwidth}{!}{%
\begin{tikzpicture}[
  cell/.style={draw, minimum width=0.41cm, minimum height=0.4cm, inner sep=0pt, font=\scriptsize},
  act/.style={cell, fill=blue!10}, idle/.style={cell, fill=gray!12, draw=gray!45},
  ar/.style={-{Stealth[length=1.2mm]}, orange!75!black, line width=0.4pt},
  st/.style={-{Stealth[length=1.2mm]}, black!40, line width=0.4pt},
  lab/.style={font=\footnotesize, anchor=east}]
\node[font=\tiny, text=black!60] at (0.000,0.380) {0};
\node[font=\tiny, text=black!60] at (0.410,0.380) {1};
\node[font=\tiny, text=black!60] at (0.820,0.380) {2};
\node[font=\tiny, text=black!60] at (1.230,0.380) {3};
\node[font=\tiny, text=black!60] at (1.640,0.380) {4};
\node[font=\tiny, text=black!60] at (2.050,0.380) {5};
\node[font=\tiny, text=black!60] at (2.460,0.380) {6};
\node[font=\tiny, text=black!60] at (2.870,0.380) {7};
\node[font=\tiny, text=black!60] at (3.280,0.380) {8};
\node[font=\tiny, text=black!60] at (3.690,0.380) {9};
\node[font=\tiny, text=black!60] at (4.100,0.380) {10};
\node[font=\tiny, text=black!60] at (4.510,0.380) {11};
\node[font=\tiny, text=black!60] at (4.920,0.380) {12};
\node[font=\tiny, text=black!60] at (5.330,0.380) {13};
\node[font=\tiny, text=black!60] at (5.740,0.380) {14};
\node[font=\tiny, text=black!60] at (6.150,0.380) {15};
\node[font=\tiny, text=black!60] at (6.560,0.380) {16};
\node[font=\tiny, text=black!60] at (6.970,0.380) {17};
\node[font=\tiny, text=black!60] at (7.380,0.380) {18};
\node[font=\tiny, text=black!60] at (7.790,0.380) {19};
\node[font=\tiny, text=black!60] at (8.200,0.380) {20};
\node[font=\tiny, text=black!60] at (8.610,0.380) {21};
\node[font=\tiny, text=black!60] at (9.020,0.380) {22};
\node[font=\tiny, text=black!60] at (9.430,0.380) {23};
\node[font=\tiny, text=black!60] at (9.840,0.380) {24};
\node[font=\tiny, text=black!60] at (10.250,0.380) {25};
\node[font=\tiny, text=black!60] at (10.660,0.380) {26};
\node[font=\tiny, text=black!60] at (11.070,0.380) {27};
\node[font=\tiny, text=black!60] at (11.480,0.380) {28};
\node[font=\tiny, text=black!60] at (11.890,0.380) {29};
\node[font=\tiny, text=black!60] at (12.300,0.380) {30};
\node[font=\tiny, text=black!60] at (12.710,0.380) {31};
\node[lab] at (-0.3,0.38) {lane};
\node[act] (n0-0) at (0.000,0.000) {2};
\node[act] (n0-1) at (0.410,0.000) {-1};
\node[act] (n0-2) at (0.820,0.000) {3};
\node[act] (n0-3) at (1.230,0.000) {-3};
\node[act] (n0-4) at (1.640,0.000) {-2};
\node[act] (n0-5) at (2.050,0.000) {5};
\node[act] (n0-6) at (2.460,0.000) {-2};
\node[act] (n0-7) at (2.870,0.000) {2};
\node[act] (n0-8) at (3.280,0.000) {-3};
\node[act] (n0-9) at (3.690,0.000) {5};
\node[act] (n0-10) at (4.100,0.000) {0};
\node[act] (n0-11) at (4.510,0.000) {-3};
\node[act] (n0-12) at (4.920,0.000) {-2};
\node[act] (n0-13) at (5.330,0.000) {3};
\node[act] (n0-14) at (5.740,0.000) {3};
\node[act] (n0-15) at (6.150,0.000) {-2};
\node[act] (n0-16) at (6.560,0.000) {0};
\node[act] (n0-17) at (6.970,0.000) {-2};
\node[act] (n0-18) at (7.380,0.000) {5};
\node[act] (n0-19) at (7.790,0.000) {3};
\node[act] (n0-20) at (8.200,0.000) {-3};
\node[act] (n0-21) at (8.610,0.000) {-2};
\node[act] (n0-22) at (9.020,0.000) {0};
\node[act] (n0-23) at (9.430,0.000) {-3};
\node[act] (n0-24) at (9.840,0.000) {3};
\node[act] (n0-25) at (10.250,0.000) {-3};
\node[act] (n0-26) at (10.660,0.000) {0};
\node[act] (n0-27) at (11.070,0.000) {-3};
\node[act] (n0-28) at (11.480,0.000) {5};
\node[act] (n0-29) at (11.890,0.000) {-1};
\node[act] (n0-30) at (12.300,0.000) {1};
\node[act] (n0-31) at (12.710,0.000) {3};
\node[lab] at (-0.3,0.000) {valori iniziali};
\node[act] (n1-0) at (0.000,-1.250) {2};
\node[act] (n1-1) at (0.410,-1.250) {-3};
\node[act] (n1-2) at (0.820,-1.250) {8};
\node[act] (n1-3) at (1.230,-1.250) {0};
\node[act] (n1-4) at (1.640,-1.250) {-5};
\node[act] (n1-5) at (2.050,-1.250) {3};
\node[act] (n1-6) at (2.460,-1.250) {-2};
\node[act] (n1-7) at (2.870,-1.250) {-1};
\node[act] (n1-8) at (3.280,-1.250) {0};
\node[act] (n1-9) at (3.690,-1.250) {2};
\node[act] (n1-10) at (4.100,-1.250) {0};
\node[act] (n1-11) at (4.510,-1.250) {-6};
\node[act] (n1-12) at (4.920,-1.250) {3};
\node[act] (n1-13) at (5.330,-1.250) {2};
\node[act] (n1-14) at (5.740,-1.250) {4};
\node[act] (n1-15) at (6.150,-1.250) {1};
\node[idle] (n1-16) at (6.560,-1.250) {};
\node[idle] (n1-17) at (6.970,-1.250) {};
\node[idle] (n1-18) at (7.380,-1.250) {};
\node[idle] (n1-19) at (7.790,-1.250) {};
\node[idle] (n1-20) at (8.200,-1.250) {};
\node[idle] (n1-21) at (8.610,-1.250) {};
\node[idle] (n1-22) at (9.020,-1.250) {};
\node[idle] (n1-23) at (9.430,-1.250) {};
\node[idle] (n1-24) at (9.840,-1.250) {};
\node[idle] (n1-25) at (10.250,-1.250) {};
\node[idle] (n1-26) at (10.660,-1.250) {};
\node[idle] (n1-27) at (11.070,-1.250) {};
\node[idle] (n1-28) at (11.480,-1.250) {};
\node[idle] (n1-29) at (11.890,-1.250) {};
\node[idle] (n1-30) at (12.300,-1.250) {};
\node[idle] (n1-31) at (12.710,-1.250) {};
\node[lab] at (-0.3,-1.250) {passo 1 (\texttt{off} = 16)};
\node[act] (n2-0) at (0.000,-2.500) {2};
\node[act] (n2-1) at (0.410,-2.500) {-1};
\node[act] (n2-2) at (0.820,-2.500) {8};
\node[act] (n2-3) at (1.230,-2.500) {-6};
\node[act] (n2-4) at (1.640,-2.500) {-2};
\node[act] (n2-5) at (2.050,-2.500) {5};
\node[act] (n2-6) at (2.460,-2.500) {2};
\node[act] (n2-7) at (2.870,-2.500) {0};
\node[idle] (n2-8) at (3.280,-2.500) {};
\node[idle] (n2-9) at (3.690,-2.500) {};
\node[idle] (n2-10) at (4.100,-2.500) {};
\node[idle] (n2-11) at (4.510,-2.500) {};
\node[idle] (n2-12) at (4.920,-2.500) {};
\node[idle] (n2-13) at (5.330,-2.500) {};
\node[idle] (n2-14) at (5.740,-2.500) {};
\node[idle] (n2-15) at (6.150,-2.500) {};
\node[idle] (n2-16) at (6.560,-2.500) {};
\node[idle] (n2-17) at (6.970,-2.500) {};
\node[idle] (n2-18) at (7.380,-2.500) {};
\node[idle] (n2-19) at (7.790,-2.500) {};
\node[idle] (n2-20) at (8.200,-2.500) {};
\node[idle] (n2-21) at (8.610,-2.500) {};
\node[idle] (n2-22) at (9.020,-2.500) {};
\node[idle] (n2-23) at (9.430,-2.500) {};
\node[idle] (n2-24) at (9.840,-2.500) {};
\node[idle] (n2-25) at (10.250,-2.500) {};
\node[idle] (n2-26) at (10.660,-2.500) {};
\node[idle] (n2-27) at (11.070,-2.500) {};
\node[idle] (n2-28) at (11.480,-2.500) {};
\node[idle] (n2-29) at (11.890,-2.500) {};
\node[idle] (n2-30) at (12.300,-2.500) {};
\node[idle] (n2-31) at (12.710,-2.500) {};
\node[lab] at (-0.3,-2.500) {passo 2 (\texttt{off} = 8)};
\node[act] (n3-0) at (0.000,-3.750) {0};
\node[act] (n3-1) at (0.410,-3.750) {4};
\node[act] (n3-2) at (0.820,-3.750) {10};
\node[act] (n3-3) at (1.230,-3.750) {-6};
\node[idle] (n3-4) at (1.640,-3.750) {};
\node[idle] (n3-5) at (2.050,-3.750) {};
\node[idle] (n3-6) at (2.460,-3.750) {};
\node[idle] (n3-7) at (2.870,-3.750) {};
\node[idle] (n3-8) at (3.280,-3.750) {};
\node[idle] (n3-9) at (3.690,-3.750) {};
\node[idle] (n3-10) at (4.100,-3.750) {};
\node[idle] (n3-11) at (4.510,-3.750) {};
\node[idle] (n3-12) at (4.920,-3.750) {};
\node[idle] (n3-13) at (5.330,-3.750) {};
\node[idle] (n3-14) at (5.740,-3.750) {};
\node[idle] (n3-15) at (6.150,-3.750) {};
\node[idle] (n3-16) at (6.560,-3.750) {};
\node[idle] (n3-17) at (6.970,-3.750) {};
\node[idle] (n3-18) at (7.380,-3.750) {};
\node[idle] (n3-19) at (7.790,-3.750) {};
\node[idle] (n3-20) at (8.200,-3.750) {};
\node[idle] (n3-21) at (8.610,-3.750) {};
\node[idle] (n3-22) at (9.020,-3.750) {};
\node[idle] (n3-23) at (9.430,-3.750) {};
\node[idle] (n3-24) at (9.840,-3.750) {};
\node[idle] (n3-25) at (10.250,-3.750) {};
\node[idle] (n3-26) at (10.660,-3.750) {};
\node[idle] (n3-27) at (11.070,-3.750) {};
\node[idle] (n3-28) at (11.480,-3.750) {};
\node[idle] (n3-29) at (11.890,-3.750) {};
\node[idle] (n3-30) at (12.300,-3.750) {};
\node[idle] (n3-31) at (12.710,-3.750) {};
\node[lab] at (-0.3,-3.750) {passo 3 (\texttt{off} = 4)};
\node[act] (n4-0) at (0.000,-5.000) {10};
\node[act] (n4-1) at (0.410,-5.000) {-2};
\node[idle] (n4-2) at (0.820,-5.000) {};
\node[idle] (n4-3) at (1.230,-5.000) {};
\node[idle] (n4-4) at (1.640,-5.000) {};
\node[idle] (n4-5) at (2.050,-5.000) {};
\node[idle] (n4-6) at (2.460,-5.000) {};
\node[idle] (n4-7) at (2.870,-5.000) {};
\node[idle] (n4-8) at (3.280,-5.000) {};
\node[idle] (n4-9) at (3.690,-5.000) {};
\node[idle] (n4-10) at (4.100,-5.000) {};
\node[idle] (n4-11) at (4.510,-5.000) {};
\node[idle] (n4-12) at (4.920,-5.000) {};
\node[idle] (n4-13) at (5.330,-5.000) {};
\node[idle] (n4-14) at (5.740,-5.000) {};
\node[idle] (n4-15) at (6.150,-5.000) {};
\node[idle] (n4-16) at (6.560,-5.000) {};
\node[idle] (n4-17) at (6.970,-5.000) {};
\node[idle] (n4-18) at (7.380,-5.000) {};
\node[idle] (n4-19) at (7.790,-5.000) {};
\node[idle] (n4-20) at (8.200,-5.000) {};
\node[idle] (n4-21) at (8.610,-5.000) {};
\node[idle] (n4-22) at (9.020,-5.000) {};
\node[idle] (n4-23) at (9.430,-5.000) {};
\node[idle] (n4-24) at (9.840,-5.000) {};
\node[idle] (n4-25) at (10.250,-5.000) {};
\node[idle] (n4-26) at (10.660,-5.000) {};
\node[idle] (n4-27) at (11.070,-5.000) {};
\node[idle] (n4-28) at (11.480,-5.000) {};
\node[idle] (n4-29) at (11.890,-5.000) {};
\node[idle] (n4-30) at (12.300,-5.000) {};
\node[idle] (n4-31) at (12.710,-5.000) {};
\node[lab] at (-0.3,-5.000) {passo 4 (\texttt{off} = 2)};
\node[act] (n5-0) at (0.000,-6.250) {8};
\node[idle] (n5-1) at (0.410,-6.250) {};
\node[idle] (n5-2) at (0.820,-6.250) {};
\node[idle] (n5-3) at (1.230,-6.250) {};
\node[idle] (n5-4) at (1.640,-6.250) {};
\node[idle] (n5-5) at (2.050,-6.250) {};
\node[idle] (n5-6) at (2.460,-6.250) {};
\node[idle] (n5-7) at (2.870,-6.250) {};
\node[idle] (n5-8) at (3.280,-6.250) {};
\node[idle] (n5-9) at (3.690,-6.250) {};
\node[idle] (n5-10) at (4.100,-6.250) {};
\node[idle] (n5-11) at (4.510,-6.250) {};
\node[idle] (n5-12) at (4.920,-6.250) {};
\node[idle] (n5-13) at (5.330,-6.250) {};
\node[idle] (n5-14) at (5.740,-6.250) {};
\node[idle] (n5-15) at (6.150,-6.250) {};
\node[idle] (n5-16) at (6.560,-6.250) {};
\node[idle] (n5-17) at (6.970,-6.250) {};
\node[idle] (n5-18) at (7.380,-6.250) {};
\node[idle] (n5-19) at (7.790,-6.250) {};
\node[idle] (n5-20) at (8.200,-6.250) {};
\node[idle] (n5-21) at (8.610,-6.250) {};
\node[idle] (n5-22) at (9.020,-6.250) {};
\node[idle] (n5-23) at (9.430,-6.250) {};
\node[idle] (n5-24) at (9.840,-6.250) {};
\node[idle] (n5-25) at (10.250,-6.250) {};
\node[idle] (n5-26) at (10.660,-6.250) {};
\node[idle] (n5-27) at (11.070,-6.250) {};
\node[idle] (n5-28) at (11.480,-6.250) {};
\node[idle] (n5-29) at (11.890,-6.250) {};
\node[idle] (n5-30) at (12.300,-6.250) {};
\node[idle] (n5-31) at (12.710,-6.250) {};
\node[lab] at (-0.3,-6.250) {passo 5 (\texttt{off} = 1)};
\draw[st] (n0-0.south) -- (n1-0.north);
\draw[ar] (n0-16.south) -- (n1-0.north);
\draw[st] (n0-1.south) -- (n1-1.north);
\draw[ar] (n0-17.south) -- (n1-1.north);
\draw[st] (n0-2.south) -- (n1-2.north);
\draw[ar] (n0-18.south) -- (n1-2.north);
\draw[st] (n0-3.south) -- (n1-3.north);
\draw[ar] (n0-19.south) -- (n1-3.north);
\draw[st] (n0-4.south) -- (n1-4.north);
\draw[ar] (n0-20.south) -- (n1-4.north);
\draw[st] (n0-5.south) -- (n1-5.north);
\draw[ar] (n0-21.south) -- (n1-5.north);
\draw[st] (n0-6.south) -- (n1-6.north);
\draw[ar] (n0-22.south) -- (n1-6.north);
\draw[st] (n0-7.south) -- (n1-7.north);
\draw[ar] (n0-23.south) -- (n1-7.north);
\draw[st] (n0-8.south) -- (n1-8.north);
\draw[ar] (n0-24.south) -- (n1-8.north);
\draw[st] (n0-9.south) -- (n1-9.north);
\draw[ar] (n0-25.south) -- (n1-9.north);
\draw[st] (n0-10.south) -- (n1-10.north);
\draw[ar] (n0-26.south) -- (n1-10.north);
\draw[st] (n0-11.south) -- (n1-11.north);
\draw[ar] (n0-27.south) -- (n1-11.north);
\draw[st] (n0-12.south) -- (n1-12.north);
\draw[ar] (n0-28.south) -- (n1-12.north);
\draw[st] (n0-13.south) -- (n1-13.north);
\draw[ar] (n0-29.south) -- (n1-13.north);
\draw[st] (n0-14.south) -- (n1-14.north);
\draw[ar] (n0-30.south) -- (n1-14.north);
\draw[st] (n0-15.south) -- (n1-15.north);
\draw[ar] (n0-31.south) -- (n1-15.north);
\draw[st] (n1-0.south) -- (n2-0.north);
\draw[ar] (n1-8.south) -- (n2-0.north);
\draw[st] (n1-1.south) -- (n2-1.north);
\draw[ar] (n1-9.south) -- (n2-1.north);
\draw[st] (n1-2.south) -- (n2-2.north);
\draw[ar] (n1-10.south) -- (n2-2.north);
\draw[st] (n1-3.south) -- (n2-3.north);
\draw[ar] (n1-11.south) -- (n2-3.north);
\draw[st] (n1-4.south) -- (n2-4.north);
\draw[ar] (n1-12.south) -- (n2-4.north);
\draw[st] (n1-5.south) -- (n2-5.north);
\draw[ar] (n1-13.south) -- (n2-5.north);
\draw[st] (n1-6.south) -- (n2-6.north);
\draw[ar] (n1-14.south) -- (n2-6.north);
\draw[st] (n1-7.south) -- (n2-7.north);
\draw[ar] (n1-15.south) -- (n2-7.north);
\draw[st] (n2-0.south) -- (n3-0.north);
\draw[ar] (n2-4.south) -- (n3-0.north);
\draw[st] (n2-1.south) -- (n3-1.north);
\draw[ar] (n2-5.south) -- (n3-1.north);
\draw[st] (n2-2.south) -- (n3-2.north);
\draw[ar] (n2-6.south) -- (n3-2.north);
\draw[st] (n2-3.south) -- (n3-3.north);
\draw[ar] (n2-7.south) -- (n3-3.north);
\draw[st] (n3-0.south) -- (n4-0.north);
\draw[ar] (n3-2.south) -- (n4-0.north);
\draw[st] (n3-1.south) -- (n4-1.north);
\draw[ar] (n3-3.south) -- (n4-1.north);
\draw[st] (n4-0.south) -- (n5-0.north);
\draw[ar] (n4-1.south) -- (n5-0.north);
\node[font=\footnotesize, anchor=west] at (0.250,-6.250) {$\leftarrow$ somma dell'intera riga, nella lane 0};
\end{tikzpicture}%
  }
  \caption{Esempio di riduzione delle somme parziali in un warp con
    \texttt{\_\_shfl\_down\_sync}. A ogni passo la lane $l$ somma il valore
    della lane $l + \mathit{off}$, con $\mathit{off}$ dimezzato da 16 a 1, e le somme utili restano nelle prime $\mathit{off}$ lane (in azzurro). Le lane
    in grigio eseguono anch'esse l'istruzione, ma il loro risultato non viene
    usato. Dopo 5 passi la lane~0 contiene la somma delle 32
    lane.}\label{fig:shfl}
\end{figure}

Il \cref{lst:warp-row} riporta il ciclo principale e la riduzione.

\begin{lstlisting}[caption={Ciclo principale e riduzione nel warp di \texttt{warp\_row} (\texttt{kernel\_warp\_row.cu}).}, label={lst:warp-row}]
for (int j = lane; j < N; j += WARP_SIZE) {
    const real_t a = arow[j];                        /* lettura coalescente  */
    const real_t *xrow = xcol + (size_t)j * k;       /* riga j di X (fetta)  */
#pragma unroll
    for (int c = 0; c < KC; c++)
        acc[c] += a * xrow[c];                       /* a riusato KC volte  */
}

#pragma unroll
for (int off = WARP_SIZE / 2; off > 0; off /= 2) {
#pragma unroll
    for (int c = 0; c < KC; c++)
        acc[c] += __shfl_down_sync(FULL_MASK, acc[c], off);
}
\end{lstlisting}

\paragraph{Più parallelismo.} I thread sono ora 32 per riga invece di uno: a
parità di $M$ la GPU dispone di 32 volte più lavoro in parallelo.

\paragraph{$k$ noto solo a tempo di esecuzione.} Gli accumulatori devono stare
nei registri, ma un array con dimensione variabile non può: il compilatore lo
sposterebbe nella \emph{local memory}, che fisicamente è memoria globale, e
quindi lenta. Il kernel è perciò un \emph{template} sul numero di colonne
$K_C$, costante di compilazione, e calcola una ``fetta'' di $K_C$ colonne di
$Y$:
\begin{itemize}
  \item per i valori di $k$ richiesti dalla traccia (3, 6, 8, 20, 32) esiste
        una versione con $K_C = k$, e $A$ viene letta una sola volta;
  \item per ogni altro $k$ si usano $\lfloor k/8 \rfloor$ fette da 8 colonne più una fetta finale di dimensione $k \bmod 8$ (se diversa da 0): il codice funziona per qualunque $k$, al costo di leggere $A$
        una volta per fetta.
\end{itemize}

\paragraph{Il limite: le letture di $X$.} Ogni lane legge la riga di $X$
corrispondente al proprio elemento di $A$. Per una colonna $c$ fissata, le 32
lane leggono gli elementi $c$ di 32 righe consecutive di $X$, che in
row-major distano $k$ posizioni l'uno dall'altro: la lettura non è
coalescente, e lo è sempre meno al crescere di $k$. Per questo le prestazioni
di \texttt{warp\_row} calano all'aumentare di $k$.

\paragraph{Risultato.} \texttt{warp\_row} fa un salto di quasi 10 volte
rispetto a \texttt{thread\_row}: 296~GFLOPS a $k = 3$ e 348 a $k = 6$, molto
vicino ai tetti (305 e 374). Poi però crolla al crescere di $k$: 178 a $k = 8$
e 91 a $k = 32$. È l'effetto delle letture non coalescenti di $X$ descritte
sopra, che diventano tanto più costose quanto più $X$ è larga. Il kernel
successivo risolve il problema.


\subsection{X trasposta: \texttt{warp\_row\_xt}}\label{sec:warp-xt}

Il terzo kernel mantiene l'intero schema di \texttt{warp\_row} (un warp per
riga, $A$ letta una volta e in modo coalescente, accumulatori nei registri,
riduzione con \texttt{\_\_shfl\_down\_sync}, template su $K_C$) e ne cambia un
solo aspetto: il modo in cui si legge $X$.

Prima del calcolo $X$ viene trasposta in una copia
$X^T$ di dimensione $k \times N$, cioè $X$ in column-major:
$X^T[c][j] = X[j][c]$. Per una colonna $c$ fissata, le 32 lane leggono ora
$X^T[c][j], X^T[c][j+1], \dots, X^T[c][j+31]$, elementi consecutivi: la
lettura di $X$ diventa coalescente come quella di $A$.

Rispetto a \texttt{warp\_row} cambia solo la lettura di $X$ nel ciclo interno
(\cref{lst:warp-xt}).

\begin{lstlisting}[caption={Ciclo principale di \texttt{warp\_row\_xt} (\texttt{kernel\_warp\_row\_xt.cu}).}, label={lst:warp-xt}]
for (int j = lane; j < N; j += WARP_SIZE) {
    const real_t a = arow[j];                       /* coalescente               */
#pragma unroll
    for (int c = 0; c < KC; c++)
        acc[c] += a * xt0[(size_t)c * ldxt + j];    /* coalescente: j consecutivi */
}
\end{lstlisting}

\paragraph{Perché non tenere $X$ trasposta ovunque.} Il layout migliore
dipende da chi legge. La parte MPI e la CPU usano $X$ row-major: le righe
consecutive formano un blocco contiguo per il broadcast (\cref{sec:layout}), e
sulla CPU le righe contigue servono ai registri vettoriali (\cref{sec:cpu}).
La trasposta si crea quindi solo sulla GPU, dove serve.

\paragraph{Dove e quanto costa.} La trasposizione è una preparazione dei dati
di ingresso, come la copia di $X$ sulla GPU: il backend la esegue in
\texttt{be\_set\_x}, subito dopo la copia, quindi è esclusa dal tempo del
nucleo e misurata insieme alla copia (\cref{sec:metrica}). Il suo costo è
piccolo, perché legge e scrive $N k$ elementi, contro gli $M N$ elementi di $A$
letti dal calcolo. Il kernel di trasposizione usa un thread per elemento:
thread consecutivi scrivono elementi consecutivi di $X^T$, quindi le scritture sono coalescenti, mentre le letture da $X$ hanno passo $k$, piccolo, e restano in
poche linee di cache.

\paragraph{Allineamento di $X^T$.} Le righe di $X^T$ sono lunghe $N$, come
quelle di $A$, e ricevono lo stesso padding a \SI{128}{byte}
(\cref{sec:padding}).

\paragraph{Risultato.} \texttt{warp\_row\_xt} risolve proprio questo
problema e non crolla più: 298~GFLOPS a $k = 3$, cioè il \SI{98}{\percent} del
tetto della memoria, e tra 350 e 354 per tutti i $k$ da 6 a 32, cioè il
\SIrange{94}{95}{\percent} del picco di calcolo. È il kernel migliore e lavora
praticamente al limite della GPU.

\subsection{Memoria shared: \texttt{warp\_row\_smem}}\label{sec:warp-smem}

Il quarto kernel mantiene lo schema di \texttt{warp\_row\_xt} e cambia il punto
da cui i warp leggono $X$.

\paragraph{Il problema.} Un blocco contiene 8 warp, cioè 8 righe diverse di
$A$, ma la $X$ che leggono è la stessa: allo stesso passo ognuno degli 8 warp
chiede gli stessi elementi di $X^T$ (\cref{fig:smem}). Gli stessi valori
vengono richiesti 8
volte, e a servirli è la cache, se li trova. Il traffico verso $X$ non è
piccolo: a ogni giro un warp legge 32 elementi di $A$ e $32 K_C$ elementi di
$X^T$, quindi le richieste per $X$ sono $K_C$ volte quelle per $A$. Con
$k = 32$ in \texttt{double} e $N = 16384$, $X^T$ occupa \SI{4}{MiB}, quanto
l'intera cache L2.

\begin{figure}[H]
  \centering
  \begin{tikzpicture}[
    font=\footnotesize,
    arow/.style={draw, fill=gray!10, minimum width=1.5cm, minimum height=0.5cm, inner sep=1pt},
    warp/.style={draw, fill=white, minimum width=1.3cm, minimum height=0.5cm, inner sep=1pt},
    xg/.style={draw, fill=blue!12, minimum width=1.2cm, minimum height=2.6cm, align=center},
    sh/.style={draw, fill=orange!18, minimum width=1.2cm, minimum height=2.6cm, align=center},
    ra/.style={-{Stealth[length=1.8mm]}, gray!70},
    rx/.style={-{Stealth[length=1.8mm]}, thick, blue!60!black},
    rs/.style={-{Stealth[length=1.8mm]}, thick, orange!70!black},
    hd/.style={font=\scriptsize, text=black!70},
    lab/.style={font=\scriptsize, align=center}]
    % ---------- warp_row_xt ----------
    \begin{scope}
      \foreach \y/\l/\w in {0/i/0, -0.7/i{+}1/1, -2.1/i{+}7/7} {
        \node[arow] (a\w) at (0, \y) {riga $\l$};
        \node[warp] (w\w) at (2.3, \y) {warp \w};
        \draw[ra] (a\w.east) -- (w\w.west);
      }
      \node at (0, -1.35) {$\vdots$};
      \node at (2.3, -1.35) {$\vdots$};
      \node[xg] (x) at (4.9, -1.05) {$X^T$};
      \foreach \w in {0,1,7} \draw[rx] (x.west) -- (w\w.east);
      \node[hd] at (0, 0.6) {$A$ (globale)};
      \node[hd] at (2.3, 0.6) {blocco};
      \node[hd] at (4.9, 0.6) {globale};
      \node[lab, text=blue!60!black] at (3.6, -3.0) {8 letture degli\\stessi elementi};
      \node[font=\footnotesize] at (2.45, 1.25) {\texttt{warp\_row\_xt}};
    \end{scope}
    % ---------- warp_row_smem ----------
    \begin{scope}[xshift=7.6cm]
      \foreach \y/\l/\w in {0/i/0, -0.7/i{+}1/1, -2.1/i{+}7/7} {
        \node[arow] (a\w) at (0, \y) {riga $\l$};
        \node[warp] (w\w) at (2.3, \y) {warp \w};
        \draw[ra] (a\w.east) -- (w\w.west);
      }
      \node at (0, -1.35) {$\vdots$};
      \node at (2.3, -1.35) {$\vdots$};
      \node[sh] (s) at (4.7, -1.05) {tile di\\$X^T$};
      \node[xg] (x) at (6.9, -1.05) {$X^T$};
      \foreach \w in {0,1,7} \draw[rs] (s.west) -- (w\w.east);
      \draw[rx, line width=1.1pt] (x.west) -- (s.east);
      \node[hd] at (0, 0.6) {$A$ (globale)};
      \node[hd] at (2.3, 0.6) {blocco};
      \node[hd] at (4.7, 0.6) {shared};
      \node[hd] at (6.9, 0.6) {globale};
      \node[lab, text=orange!70!black] at (3.5, -3.0) {8 letture\\dalla shared};
      \node[lab, text=blue!60!black] at (5.8, -3.0) {1 copia\\per blocco};
      \node[font=\footnotesize] at (3.45, 1.25) {\texttt{warp\_row\_smem}};
    \end{scope}
  \end{tikzpicture}
  \caption{Letture di $X$ in un blocco di 8 warp. In \texttt{warp\_row\_xt}
    ogni warp legge dalla memoria globale gli stessi elementi di $X^T$, mentre in \texttt{warp\_row\_smem} il blocco li copia una sola volta in memoria
    shared, e gli 8 warp li leggono da lì. Le frecce indicano il verso dei
    dati.}\label{fig:smem}
\end{figure}

\paragraph{La soluzione.} Il blocco carica $X$ una sola volta nella
\emph{memoria shared}, una memoria interna allo SM, condivisa dai thread del
blocco e veloce quanto la L1, e gli 8 warp la leggono da lì. La riga viene
percorsa a \emph{tile} di $T_C$ colonne. Per ogni tile:
\begin{enumerate}
  \item i 256 thread copiano insieme la porzione di $X^T$ del tile in memoria
        shared, con letture coalescenti, una sola volta per blocco invece che
        una per warp;
  \item \texttt{\_\_syncthreads()}: nessun thread usa la shared prima che il
        caricamento sia completo;
  \item ogni warp calcola la propria riga sul tile, leggendo $A$ dalla memoria
        globale e $X$ dalla shared;
  \item \texttt{\_\_syncthreads()}: nessun thread sovrascrive la shared con il
        tile successivo mentre altri la stanno ancora leggendo.
\end{enumerate}
Le richieste di $X$ verso la cache diminuiscono così di 8 volte.
Il ciclo sui tile è riportato nel \cref{lst:warp-smem}.

\begin{lstlisting}[caption={Ciclo sui tile di \texttt{warp\_row\_smem} (\texttt{kernel\_warp\_row\_smem.cu}, commenti abbreviati).}, label={lst:warp-smem}]
for (int j0 = 0; j0 < N; j0 += TC) {
    const int w = min(TC, N - j0);            /* colonne valide nel tile */

    /* 1. caricamento cooperativo del tile in shared */
    constexpr int LOADS = (KC * TC + THREADS - 1) / THREADS;
#pragma unroll 4
    for (int it = 0; it < LOADS; it++) {
        const int e = it * THREADS + tid;
        const int c = e / TC, jj = e % TC;
        if (e < KC * TC && jj < w) xs[c][jj] = xt0[(size_t)c * ldxt + j0 + jj];
    }
    __syncthreads();                          /* 2. tile completo */

    /* 3. calcolo sul tile, con X letta dalla shared */
    if (active) {
#pragma unroll 2
        for (int jj = lane; jj < w; jj += WARP_SIZE) {
            const real_t a = arow[j0 + jj];   /* coalescente */
#pragma unroll
            for (int c = 0; c < KC; c++)
                acc[c] += a * xs[c][jj];      /* shared: lane su elementi consecutivi */
        }
    }
    __syncthreads();                          /* 4. xs libera per il tile successivo */
}
\end{lstlisting}

\paragraph{Dimensione del tile.} Un tile è la porzione di $X^T$ che il blocco
tiene in memoria shared in un dato momento: $K_C$ righe da $T_C$ colonne, cioè
$K_C\,T_C$ elementi. Conviene che sia grande, perché ogni tile richiede un
caricamento e due sincronizzazioni, ma la memoria shared è limitata, e la
dimensione va scelta di conseguenza.

Il limite è fissato a \SI{16}{KiB} per blocco. Uno SM dispone di
\SI{64}{KiB} di shared e ospita al più 4 blocchi da 256 thread: con
\SI{16}{KiB} ciascuno i 4 blocchi ci stanno tutti, e la shared non riduce il
numero di blocchi attivi (\cref{sec:registri}).

In \texttt{double}, \SI{16}{KiB} corrispondono a 2048 elementi. $T_C$ è il
valore più grande per cui $K_C\,T_C$ non supera questo limite, arrotondato per
difetto a una potenza di 2:
\begin{itemize}
  \item per $k = 3$: $2048 / 3 \approx 682$, quindi $T_C = 512$;
  \item per $k = 32$: $2048 / 32 = 64$, quindi $T_C = 64$.
\end{itemize}

La potenza di 2 serve nel caricamento del tile: ogni thread ricava dalla
posizione $e$ dell'elemento che copia la sua riga $e / T_C$ e la sua colonna
$e \bmod T_C$. Con $T_C$ potenza di 2 queste due operazioni diventano uno
\emph{shift} e una maschera di bit, molto più veloci di una divisione.

\paragraph{Nessun conflitto di banco.} La memoria shared è veloce, ma va usata
rispettando una regola. È divisa in 32 \emph{banchi} che lavorano in parallelo,
e ogni banco serve un solo indirizzo alla volta. Se due lane dello stesso warp
chiedono nello stesso momento indirizzi diversi dello stesso banco si ha un
\emph{conflitto}: le richieste vengono servite una dopo l'altra, e la lettura
rallenta.

Elementi consecutivi in memoria appartengono a banchi consecutivi. In questo
kernel le 32 lane di un warp leggono 32 elementi consecutivi della stessa riga
del tile (\cref{lst:warp-smem}): le richieste si distribuiscono in modo
uniforme su tutti i banchi, e non ci sono conflitti evitabili. In
\texttt{double} un warp chiede \SI{256}{byte}, il doppio di quanto i banchi
servono in una passata: ne servono quindi due, che è il minimo possibile.

\paragraph{Nessuna uscita anticipata.} Negli altri kernel i warp dell'ultimo
blocco senza una riga da calcolare escono subito. Qui non possono:
\texttt{\_\_syncthreads()} attende tutti i thread del blocco, e se alcuni
fossero già usciti il comportamento sarebbe indefinito. I warp senza riga
partecipano quindi al caricamento e alle barriere, e saltano solo il calcolo.

\paragraph{Cosa ci si aspetta.} La L1 e la shared sono in effetti la stessa
memoria fisica: se la L1 cattura già il riuso di $X$ tra i warp, il guadagno è
piccolo, a maggior ragione perchè le barriere hanno un costo.

\paragraph{Risultato.} \texttt{warp\_row\_smem} va quasi come
\texttt{warp\_row\_xt}, ma sempre un po' meno, fino a 331 contro 352 a
$k = 32$. Tenere $X$ nella memoria shared non porta vantaggi, probabilmente
perché la cache L1 riusava già $X$ tra i warp e le barriere hanno un costo.
% NOTA: ipotesi coerente con i dati, non misurata: possibile verifica con il
% profiler (tasso di hit della cache L1 di warp_row_xt).

\subsection{Registri e occupazione}\label{sec:registri}

\paragraph{Perché servono molti warp?} A $k$ piccolo i kernel sono limitati
dalla memoria: un warp che chiede un dato deve attendere centinaia di cicli
prima di riceverlo. Lo SM nasconde questa attesa facendo lavorare nel
frattempo altri warp: più warp sono \emph{residenti} su uno SM, più l'attesa
viene coperta. Il rapporto tra i warp residenti e il massimo possibile (32 per
SM) viene chiamato \emph{occupazione}.

\paragraph{I registri limitano l'occupazione.} Uno SM ha \num{65536} registri
da 32 bit, divisi tra tutti i thread residenti. Per avere 4 blocchi da 256
thread, cioè il massimo di 1024 thread, ogni thread può usare al più 64 registri: se ne usa di più, sullo SM entrano meno blocchi. Nei kernel del
progetto la voce principale sono gli accumulatori: $K_C$ valori per thread,
ciascuno di due registri in \texttt{double}. Con $k = 32$ in \texttt{double} i
soli accumulatori occupano già 64 registri.

\paragraph{\texttt{L'utilizzo di \#pragma unroll:}} I registri non si possono
indicizzare con una variabile: un array sta nei registri solo se, per ogni
accesso, il compilatore conosce l'indice già in compilazione. Il ciclo
\begin{center}
  \texttt{for (c = 0; c < KC; c++) acc[c] += \dots}
\end{center}
usa invece un indice variabile, $c$. \texttt{\#pragma unroll} lo sostituisce
con $K_C$ istruzioni separate, \texttt{acc[0] += \dots},
\texttt{acc[1] += \dots}, e così via, ciascuna con un indice costante. Per
questo tutti i cicli sulle colonne sono srotolati completamente.

\paragraph{Il costo: il caso \texttt{warp\_row\_smem}.} Srotolare può però
aumentare i registri usati. In \texttt{warp\_row\_smem} due cicli si ripetono
per ogni tile:
\begin{itemize}
  \item il \emph{caricamento}, in cui ogni thread copia in shared più elementi
        di $X^T$;
  \item il \emph{calcolo}, in cui ogni warp scorre il tile.
\end{itemize}
Se anche questi vengono srotolati completamente, il compilatore si accorge che
le posizioni da cui ogni thread legge sono le stesse in ogni tile, spostate
solo dell'inizio del tile: calcola allora tutti quegli indirizzi una volta
sola, prima di cominciare, e li tiene in registri per tutta l'esecuzione. Con
$k = 32$ in \texttt{double} si arriva a 190 registri per thread, quindi a un
solo blocco per SM.

\paragraph{La soluzione: srotolare solo in parte.} Con
\texttt{\#pragma unroll 4} sul caricamento e \texttt{\#pragma unroll 2} sul
calcolo, il compilatore srotola 4 e 2 iterazioni alla volta: restano
abbastanza letture in corso per tenere occupata la memoria, ma gli indirizzi
da conservare nei registri sono pochi. Si resta così entro 64 registri, cioè 4
blocchi per SM, per ogni $k$ tranne $k = 32$ in \texttt{double}, dove i soli
accumulatori occupano già 64 registri: servono 124 registri e si hanno 2
blocchi per SM, come in \texttt{warp\_row\_xt}.

\paragraph{Spill.} Se un thread avesse bisogno di più registri di quanti
l'hardware ne permette, il compilatore sposterebbe alcuni valori nella
\emph{local memory} (\emph{spill}), con lo stesso costo della memoria globale.
I registri usati da ogni kernel, ed eventualispill, sono riportati 
dall'assemblatore della GPU con l'opzione \texttt{-v}: 
nel progetto si ottengono con \texttt{make NVFLAGS\_EXTRA="-Xptxas -v"}
(\cref{sec:compilazione}). I valori citati sono stati ottenuti con CUDA 12.8
per \texttt{sm\_75}.

\subsection{Confronto con cuBLAS}\label{sec:cublas}

Come termine di paragone si usa \emph{cuBLAS}, la libreria di algebra lineare
di NVIDIA, ottimizzata per ogni architettura: il prodotto $Y = A\,X$ è un caso
particolare della sua funzione \texttt{gemm}, il prodotto tra matrici
generico. cuBLAS non fa parte della soluzione: serve a capire quanto i kernel
scritti a mano si avvicinano, o si allontanano, da un'implementazione
professionale, selezionabile con \texttt{-v -1}.

cuBLAS lavora in column-major, mentre le
matrici del progetto sono row-major. Non serve però convertirle: una matrice
row-major $R \times C$, letta come column-major, è la sua trasposta
$C \times R$. Passando a cuBLAS gli stessi puntatori, quindi:
\begin{itemize}
  \item $A$ row-major ($M \times N$) viene vista come $A^T$ ($N \times M$);
  \item $X$ row-major ($N \times k$) viene vista come $X^T$ ($k \times N$);
  \item $Y$ row-major ($M \times k$) viene vista come $Y^T$ ($k \times M$).
\end{itemize}
Poiché
\[
  Y = A\,X \iff Y^T = X^T A^T,
\]
basta chiedere a cuBLAS il prodotto $X^T A^T$, senza trasposizioni, con
dimensioni $m = k$, $n = M$ e dimensione interna $N$, passando $X$ prima di
$A$. Il risultato, scritto in column-major come $Y^T$, è esattamente $Y$ in
row-major. Non si copia nulla: cambia solo l'interpretazione dei dati. La
leading dimension di $A$ si passa così com'è, quindi anche il padding delle
righe è gestito.

\paragraph{Risultato.} cuBLAS, sorprendentemente, è il più lento di tutti in
\texttt{double} a $k = 3$ e $k = 6$, con 16 e 31~GFLOPS, meno anche di \texttt{thread\_row}, e lo supera solo da $k = 8$ e arriva a 162 a $k = 32$. Il
suo tempo resta quasi costante, circa \SI{0.1}{s}, qualunque sia $k$: in
pratica cuBLAS lavora come se $k$ fosse sempre grande, ed è pensato per il
prodotto tra matrici, non per multivettori con poche colonne. In
\texttt{float}, a $k$ grande, è invece il più veloce (\cref{sec:gpu-riepilogo}).

\subsection{Trasferimenti host--device ed eventi CUDA}\label{sec:h2d}

\paragraph{Trasferimenti.} Per tutti i kernel, il backend CUDA sposta i dati tra CPU e GPU in tre
momenti:
\begin{itemize}
  \item $A$, una sola volta, in \texttt{be\_create}, con
        \texttt{cudaMemcpy2D}, che copia riga per riga solo i dati e non il
        padding (\cref{sec:padding});
  \item $X$, a ogni invocazione, in \texttt{be\_set\_x}, con eventuale
        trasposizione;
  \item $Y$, a ogni invocazione, in \texttt{be\_get\_y}, verso la CPU.
\end{itemize}
$X$ e $Y$ devono passare per la CPU a ogni invocazione perché MPI non è
CUDA-aware (\cref{sec:calcpar}). Come chiede la traccia, questi trasferimenti
sono esclusi dal tempo del nucleo e misurati a parte (\cref{sec:metrica}). Le
copie usano \texttt{cudaMemcpy}, che ritorna solo a copia completata: quando
\texttt{be\_set\_x} termina, $X$ è già sulla GPU.

\paragraph{Eventi CUDA.} Poiché il lancio di un kernel è asincrono
(\cref{sec:cuda-modello}), il tempo del calcolo sulla GPU si misura con gli
\emph{eventi CUDA}:
\begin{enumerate}
  \item un evento viene inserito nella coda della GPU prima del lancio, e un
        secondo evento dopo;
  \item la GPU annota l'istante in cui raggiunge ciascun evento;
  \item la CPU attende il secondo evento (\texttt{cudaEventSynchronize}), cioè
        la fine del calcolo;
  \item la differenza tra i due istanti (\texttt{cudaEventElapsedTime}) è il
        tempo trascorso sulla GPU tra l'inizio e la fine del kernel.
\end{enumerate}
Per un $k$ generico il calcolo è diviso in più kernel, uno per fetta di colonne
(\cref{sec:warp-row}): i due eventi racchiudono l'intera sequenza, quindi il
tempo misurato è quello di tutte le fette. Questa misura è riportata dal
programma accanto al tempo misurato con l'orologio della CPU, e le differenze tra
le due, rilevanti quando più processi condividono la GPU, sono discusse in
\cref{sec:misure}.

\subsection{Riepilogo dei kernel GPU}\label{sec:gpu-riepilogo}

La \cref{tab:ris-kernel} e la \cref{fig:kernel-k} riassumono i GFLOPS
di tutti i kernel GPU.

\begin{table}[H]
  \centering
  \caption{GFLOPS del nucleo dei kernel GPU (campagna~B,
    $16384 \times 16384$, un processo).}\label{tab:ris-kernel}
  \begin{tabular}{llrrrrr}
    \toprule
    Kernel & Precisione & $k=3$ & $k=6$ & $k=8$ & $k=20$ & $k=32$ \\
    \midrule
    \texttt{thread\_row}    & \texttt{double} &  38 &   38 &   38 &   36 &   36 \\
    \texttt{warp\_row}      & \texttt{double} & 296 &  348 &  178 &  177 &   91 \\
    \texttt{warp\_row\_xt}  & \texttt{double} & 298 &  352 &  354 &  350 &  352 \\
    \texttt{warp\_row\_smem}& \texttt{double} & 296 &  346 &  348 &  348 &  331 \\
    cuBLAS                  & \texttt{double} &  16 &   31 &   41 &  102 &  162 \\
    \midrule
    \texttt{warp\_row\_xt}  & \texttt{float}  & 600 & 1174 & 1538 & 1698 & 1883 \\
    \texttt{warp\_row\_smem}& \texttt{float}  & 597 & 1156 & 1468 & 1605 & 1674 \\
    cuBLAS                  & \texttt{float}  & 429 &  841 & 1101 & 2759 & 4285 \\
    \bottomrule
  \end{tabular}
\end{table}

\begin{figure}[H]
  \centering
\begin{tikzpicture}
\begin{groupplot}[
  group style={group size=1 by 2, vertical sep=2.2cm},
  width=0.9\textwidth, height=0.42\textwidth,
  xmin=0, xmax=34, xtick={3,6,8,20,32},
  xlabel={$k$}, ylabel={GFLOPS},
  xlabel style={font=\footnotesize}, ylabel style={font=\footnotesize},
  ticklabel style={font=\scriptsize},
  title style={font=\footnotesize},
  grid=major, grid style={gray!20},
  every axis plot/.append style={thick, mark size=1.7pt},
  /pgf/number format/1000 sep={\,},
]
\nextgroupplot[title={\texttt{double}}, ymin=0, ymax=420,
  legend to name=legkernel,
  legend style={font=\scriptsize, legend columns=5, cells={anchor=west}, column sep=6pt}]
\addplot[black!50, dashed, thin, domain=3:32, samples=200, forget plot] {min(407*x/4,374)};
\addplot[gray, mark=x] coordinates {(3,38) (6,38) (8,38) (20,36) (32,36)};
\addlegendentry{\texttt{thread\_row}}
\addplot[blue!45, mark=triangle*] coordinates {(3,296) (6,348) (8,178) (20,177) (32,91)};
\addlegendentry{\texttt{warp\_row}}
\addplot[blue!70!black, mark=*] coordinates {(3,298) (6,352) (8,354) (20,350) (32,352)};
\addlegendentry{\texttt{warp\_row\_xt}}
\addplot[teal, mark=diamond*] coordinates {(3,296) (6,346) (8,348) (20,348) (32,331)};
\addlegendentry{\texttt{warp\_row\_smem}}
\addplot[red!70!black, mark=square*] coordinates {(3,16) (6,31) (8,41) (20,102) (32,162)};
\addlegendentry{cuBLAS}
\node[font=\scriptsize, text=black!60, anchor=south east] at (axis cs:33,376) {tetto};
\nextgroupplot[title={\texttt{float}}, ymin=0, ymax=7000]
\addplot[black!50, dashed, thin, domain=3:32, samples=2, forget plot] {407*x/2};
\addplot[blue!70!black, mark=*] coordinates {(3,600) (6,1174) (8,1538) (20,1698) (32,1883)};
\addplot[teal, mark=diamond*] coordinates {(3,597) (6,1156) (8,1468) (20,1605) (32,1674)};
\addplot[red!70!black, mark=square*] coordinates {(3,429) (6,841) (8,1101) (20,2759) (32,4285)};
\node[font=\scriptsize, text=black!60, anchor=south east] at (axis cs:31,6550) {tetto};
\end{groupplot}
\end{tikzpicture}

  \vspace{2pt}\ref{legkernel}
  \caption{GFLOPS del nucleo dei kernel GPU al variare di $k$ (campagna~B,
    $16384 \times 16384$, un processo). La linea tratteggiata è il tetto della
    roofline: in \texttt{double} $\min(407 \cdot k/4,\ 374)$, in \texttt{float}
    $407 \cdot k/2$ (\cref{sec:roofline}).}\label{fig:kernel-k}
\end{figure}

\paragraph{GPU in \texttt{float}.} In singola precisione il tetto del calcolo è
molto più alto, e per tutti i $k$ della traccia il limite teorico è la memoria
(\cref{sec:roofline}).
\begin{itemize}
  \item Fino a $k = 8$, \texttt{warp\_row\_xt} segue il tetto della memoria quasi
        alla perfezione: 600~GFLOPS a $k = 3$, 1174 a $k = 6$ e 1538 a $k = 8$,
        tra il \SI{94}{\percent} e il \SI{98}{\percent} del tetto.
  \item Da $k = 20$ in poi, invece, si stacca dal tetto: 1698~GFLOPS a $k = 20$
        e 1883 a $k = 32$, contro tetti teorici di 4070 e 6512, cioè solo il
        \SIrange{29}{42}{\percent}. Qui il limite, probabilmente, non è più la
        memoria globale ma le letture dentro il chip: per ogni moltiplicazione-somma il kernel deve
        leggere un valore di $X$ dalla cache, e quando $k$ è grande le operazioni
        sono così tante che la cache non riesce a fornire i dati abbastanza in
        fretta. In \texttt{double} questo limite non si vede, perché lì il
        calcolo si ferma prima, a 374~GFLOPS.
% NOTA: ipotesi coerente con i dati, non misurata: possibile verifica con il
% profiler (occupazione della cache L1 di warp_row_xt in float a k grande).
  \item \texttt{warp\_row\_smem} va di nuovo un po' meno di
        \texttt{warp\_row\_xt}, come in \texttt{double}.
  \item cuBLAS parte sotto \texttt{warp\_row\_xt} (429~GFLOPS a $k = 3$), ma a
        $k$ grande lo supera nettamente: 2759~GFLOPS a $k = 20$ e 4285 a
        $k = 32$, cioè 1,6 e 2,3 volte di più. È il caso in cui conviene usare
        la libreria. Il codice interno di cuBLAS non è noto, ma un'ipotesi
        plausibile è che riusi ogni valore di $X$ per più righe di $A$, come fa
        \texttt{cpu\_block} sulla CPU, leggendo così meno dati dalla cache per
        ogni operazione. Sarebbe anche un possibile sviluppo per
        \texttt{warp\_row\_xt}: calcolare più righe per warp.
\end{itemize}

% =====================================================================
\section{Implementazione CPU-based}\label{sec:cpu}

La traccia chiede una parallelizzazione interna al processo in OpenMP o in CUDA
e non prevede una versione che usi solo la CPU. Questa versione è stata
sviluppata comunque, sia per averre un termine di paragone con la GPU, 
sia per una questione di \emph{hardware del server},
che ha una sola GPU ma 20 core fisici: una versione che usa solo MPI 
permette di sfruttare tutti i core, mentre nella versione
 CUDA i processi sono limitati dalla condivisione dell'unica GPU.

In questa versione ogni processo calcola il proprio blocco su un solo core, e
il parallelismo è dato interamente dal numero di processi: sul server si lancia
un processo per core fisico, fino a 20 (\cref{sec:binding}). Rispetto al
backend CUDA naturalmente non servono trasferimenti tra CPU e GPU: $A$, $X$ e
$Y$ si trovano già nella memoria della CPU. Il backend legge $X$ direttamente
dal buffer di \texttt{mxm.c}, senza copiarla, mentre $Y$ viene solo copiata nel
buffer della riduzione.

\paragraph{Come sono misurati i kernel.} I kernel CPU sono misurati sulla
stessa matrice dei kernel GPU, $16384 \times 16384$, con 20 processi su una
griglia $4 \times 5$ (campagna~B). Il riferimento è il tetto della memoria,
$71 \cdot k/4$: 53~GFLOPS a $k = 3$, 142 a $k = 8$ e 568 a $k = 32$. Il picco di calcolo della CPU non è stato misurato.

\subsection{Versione di partenza: \texttt{cpu\_row}}\label{sec:cpu-row}

Il primo kernel calcola una riga di $Y$ alla volta, con lo stesso ordine dei
cicli del riferimento seriale (\cref{sec:collaudo}): per ogni riga $i$ scorre
la riga $i$ di $A$, e ogni elemento $a_{ij}$, letto una sola volta, viene
moltiplicato per l'intera riga $j$ di $X$, aggiornando i $k$ accumulatori della
riga $i$ di $Y$, in questo modo quindi l'intera matrice $A$ attraversa la memoria una sola volta.

Come \texttt{thread\_row} sulla GPU, il kernel è volutamente semplice e fa da
riferimento. Presenta due limiti, risolti dal prossimo kernel:
\begin{itemize}
  \item \emph{accumulatori in memoria}: $k$ è noto solo a tempo di esecuzione,
        quindi gli accumulatori stanno in un array in memoria, cioè nella cache
        L1, e non nei registri: a ogni moltiplicazione-somma il core deve
        leggere l'accumulatore e riscriverlo;
  \item \emph{$X$ non riusata}: ogni valore di $X$ letto serve per una sola
        riga di $A$. Per ogni operazione servono quindi due letture ($x$ e
        accumulatore) e una scrittura, il core passa più tempo ad attendere
        la memoria che a calcolare.
\end{itemize}

\paragraph{Risultato.} \texttt{cpu\_row} cresce con $k$, ma sempre meno: da
46~GFLOPS a $k = 3$ fino a soli 124 a $k = 32$. Il limite sono gli
accumulatori tenuti in memoria: a ogni operazione il core deve leggere e
riscrivere un accumulatore, e questo lavoro in più cresce con $k$.

\subsection{Register blocking: \texttt{cpu\_block}}\label{sec:cpu-block}

\texttt{cpu\_row} riusa $A$, perché ogni elemento serve
$k$ colonne, ma non $X$, perché ogni valore di $X$ letto serve una sola riga.
\texttt{cpu\_block} calcola invece $R_B$ righe di $Y$ insieme:

\begin{center}
\fbox{\begin{minipage}{0.85\textwidth}\small
\begin{tabbing}
\hspace{1.5em}\=\hspace{1.5em}\=\hspace{1.5em}\=\kill
per ogni $j$ da $0$ a $N-1$:\\
\>leggi la riga $j$ di $X$, cioè $k$ valori, \emph{una volta}\\
\>per ognuna delle $R_B$ righe $r$ del blocco:\\
\>\>leggi $a = A[r][j]$\\
\>\>aggiorna i $k$ accumulatori della riga $r$:
  $\mathit{acc}[r][c] \mathrel{+}= a \cdot x[c]$
\end{tabbing}
\end{minipage}}
\end{center}

Si ottengono due forme di riuso, una per ciascun operando:
\begin{itemize}
  \item \emph{riuso di $A$}: ogni elemento $a$ viene usato $k$ volte, come in
        \texttt{cpu\_row};
  \item \emph{riuso di $X$}: ogni valore $x$ viene usato $R_B$ volte, quindi
        le letture di $X$ per operazione si riducono di un fattore $R_B$.
\end{itemize}
$A$ viene comunque letta una sola volta, perché le $R_B$ righe si scorrono
insieme. È lo stesso principio di \texttt{warp\_row\_smem}
(\cref{sec:warp-smem}): lì $X$ si riusa tra 8 righe tenendola in memoria
shared, qui tra $R_B$ righe tenendola nei registri.

\paragraph{Gli accumulatori.} Un \emph{accumulatore} è la variabile in cui si
somma, termine dopo termine, un elemento di $Y$:
$y_{rc} = \sum_j a_{rj}\, x_{jc}$. Un blocco ne richiede $R_B \times k$, uno per
ogni elemento del blocco di $Y$ in calcolo. Perché il riuso sia efficace
devono stare tutti nei registri, da cui il nome \emph{register blocking}: in
memoria, ogni aggiornamento richiederebbe una lettura e una scrittura, come in
\texttt{cpu\_row}.

\paragraph{Dimensioni note in compilazione.} Esattamente come in GPU, gli accumulatori
stanno nei registri solo se il compilatore conosce $R_B$ e $k$.
Si usa quindi una funzione \texttt{always\_inline}, che il
compilatore deve espandere, chiamata con argomenti
costanti. Per ogni coppia $(R_B, k)$ il compilatore genera così una copia
dedicata della funzione, con i cicli sulle righe e sulle colonne completamente
srotolati. Il ciclo corrispondente nel codice è riportato nel
\cref{lst:cpu-block}: i due \texttt{\#pragma GCC unroll} srotolano
completamente i cicli sulle righe e sulle colonne (al più $8 \times 32$
iterazioni). Per i $k$ della traccia esiste una versione con $A$ letta una volta, mentre per ogni altro $k$ si procede a fette di 8 colonne, come in
\texttt{warp\_row} (\cref{sec:warp-row}).

\begin{lstlisting}[caption={Ciclo principale di \texttt{cpu\_block} (\texttt{kernel\_cpu\_block.c}).}, label={lst:cpu-block}]
for (int j = 0; j < N; j++) {
    const real_t *xrow = X + (size_t)j * k;          /* riga j di X (fetta) */
#pragma GCC unroll 8
    for (int r = 0; r < rb; r++) {
        const real_t a = A[(size_t)r * lda + j];     /* a riusato kc volte  */
#pragma GCC unroll 32
        for (int c = 0; c < kc; c++)
            acc[r][c] += a * xrow[c];                /* x riusato rb volte  */
    }
}
\end{lstlisting}

\subsection{Registri vettoriali e vettorizzazione}\label{sec:vettorizzazione}

\paragraph{Registri vettoriali.} Un registro vettoriale contiene più valori, e
un'unica istruzione può operare contemporaneamente su tutti i valori del registro
(\emph{SIMD}, \emph{Single Instruction, Multiple Data}). Lo Xeon del server ha
32 registri vettoriali (AVX-512). Il compilatore, con le opzioni del progetto,
li usa da 256 bit, cioè 4 \texttt{double} o 8 \texttt{float} l'uno. In tutto i
registri contengono quindi 128 \texttt{double} o 256 \texttt{float}.

\paragraph{Quante righe per blocco.} $R_B$, il numero di righe di $Y$ calcolate
insieme (\cref{sec:cpu-block}), è vincolato proprio da questa capacità: gli
$R_B \times k$ accumulatori devono stare nei registri, quindi $R_B \times k$
non può superare un \emph{budget} \texttt{ACC\_BUDGET}, pari a , come appena detto, 128
accumulatori in \texttt{double} e 256 in \texttt{float}, cioè \SI{1024}{byte}
di registri. Ne segue:
\[
  R_B = \left\lfloor \frac{\texttt{ACC\_BUDGET}}{k} \right\rfloor ,
\]
limitato tra 1 e 8 (\cref{lst:rb}). Il limite inferiore è una protezione: il
budget si può cambiare in compilazione, e se fosse più piccolo di $k$ la
divisione darebbe 0 righe. Il limite superiore serve innanzitutto al
compilatore: gli accumulatori sono un array di dimensioni fisse,
\texttt{acc[RB\_MAX][32]}, e il ciclo sulle righe viene srotolato al più 8
volte. Senza questo limite, a $k$ piccolo si arriverebbe a $128/3 \approx 42$
righe, ma non ci si attende un guadagno perché si arriverebbe a 42 flussi
distinti, più di quanti cache e prefetch gestiscano bene.

\begin{lstlisting}[caption={Calcolo di $R_B$ (\texttt{kernel\_cpu\_block.c}).}, label={lst:rb}]
#define ACC_BUDGET (1024 / (int)sizeof(real_t))
#define RB_MAX 8
#define RB_FOR(kc) ((ACC_BUDGET / (kc)) > RB_MAX ? RB_MAX : \
                    (ACC_BUDGET / (kc)) < 1 ? 1 : (ACC_BUDGET / (kc)))
\end{lstlisting}

\begin{table}[H]
  \centering
  \caption{Righe per blocco $R_B$ di \texttt{cpu\_block} per i $k$ della traccia.}\label{tab:rb}
  \begin{tabular}{lccccc}
    \toprule
    $k$ & 3 & 6 & 8 & 20 & 32 \\
    \midrule
    $R_B$ in \texttt{double} & 8 & 8 & 8 & 6 & 4 \\
    $R_B$ in \texttt{float}  & 8 & 8 & 8 & 8 & 8 \\
    \bottomrule
  \end{tabular}
\end{table}

A $k = 32$ in \texttt{double} i 128 accumulatori occupano tutti i registri, e il
compilatore deve tenerne qualcuno nella cache L1 (\emph{spill}). Conviene
comunque, perché leggere $X$ la metà delle volte fa risparmiare più di quanto
costi lo spill. La scelta del budget è stata fatta misurando sul server i valori
32, 64 e 128 (\cref{sec:ris-padding}), e il budget si può cambiare in compilazione
con \texttt{CFLAGS\_EXTRA=-DACC\_BUDGET=\dots} (\cref{sec:compilazione}).

\paragraph{Vettorizzazione nella direzione giusta.} Dopo lo srotolamento dei
cicli su $r$ e $c$ resta un solo ciclo, quello su $j$, e il vettorizzatore di
GCC tende a vettorizzare quello, perché $A$ è contigua lungo $j$: carica 4
valori consecutivi di $A$ e deve poi rimescolarli con permutazioni per
combinarli con $X$. Le permutazioni esauriscono i registri e gli accumulatori
finiscono in memoria: in alcuni casi le prestazioni scendevano a meno di un
quarto. La vettorizzazione giusta è invece lungo $c$, cioè su colonne consecutive di $Y$:
\begin{itemize}
  \item un registro contiene 4 colonne consecutive della riga di $X$;
  \item il valore $a$ viene replicato in tutte le posizioni di un registro
        (\emph{broadcast});
  \item un'unica FMA aggiorna 4 accumulatori della riga.
\end{itemize}
Per ottenerla, nel file di \texttt{cpu\_block} si disattiva la vettorizzazione
dei cicli e si lascia attiva la \emph{vettorizzazione SLP} (\emph{superword-level
parallelism}, \cref{lst:novec}), che raggruppa in un'unica istruzione
vettoriale le operazioni scalari identiche e indipendenti prodotte dallo
srotolamento, come $\mathit{acc}[r][0..3] \mathrel{+}= a \cdot x[0..3]$, e
produce esattamente lo schema broadcast + FMA.
La fusione in FMA è resa possibile dall'opzione
\texttt{-ffp-contract=fast} (\cref{sec:compilazione}).

\begin{lstlisting}[caption={Disattivazione della vettorizzazione dei cicli in \texttt{kernel\_cpu\_block.c}.}, label={lst:novec}]
#pragma GCC optimize ("no-tree-loop-vectorize")
\end{lstlisting}

\paragraph{Risultato.} \texttt{cpu\_block} segue il tetto della memoria molto
da vicino fino a $k = 20$: 49, 96, 129 e 299~GFLOPS, tra l'\SI{84}{\percent} e
il \SI{92}{\percent} del tetto. A $k = 20$ va 2,6 volte più veloce di
\texttt{cpu\_row}: è l'effetto del riuso di $X$ tra più righe e degli
accumulatori nei registri. A $k = 32$ scende a 271~GFLOPS, sotto il valore di
$k = 20$ e al \SI{48}{\percent} del tetto: il blocco ha solo 4 righe, i 128
accumulatori riempiono tutti i registri e il compilatore deve tenerne qualcuno
nella cache L1. Resta comunque più del doppio di \texttt{cpu\_row}.

\subsection{Binding dei processi e NUMA}\label{sec:binding}

La versione CPU si lancia con un processo per core, legato al proprio core:
\begin{lstlisting}[language=bash,numbers=none,frame=none]
mpirun -np 20 --map-by core --bind-to core bin/mxm_cpu_double ...
\end{lstlisting}
\begin{itemize}
  \item \texttt{-{}-map-by core}: assegna i rank a core consecutivi (i rank
        0--9 finiscono sui core del socket~0 e i rank 10--19 su quelli del
        socket~1);
  \item \texttt{-{}-bind-to core}: impedisce al sistema operativo di migrare un
        processo su un altro core durante l'esecuzione.
\end{itemize}

Questo perché in caso di migrazione il processo perde il contenuto delle cache L1 e L2 del
core che lascia e, soprattutto, il suo blocco di $A$ resta nella memoria del
nodo NUMA in cui è stato generato.
Se il processo finisce sull'altro socket, ogni
lettura di $A$ diventa un accesso remoto, circa due volte più costoso
(\cref{sec:cpu-hw}). Con il binding ogni processo genera e legge il proprio
blocco sempre dallo stesso socket: i due socket lavorano ciascuno sulla
propria memoria, e la banda complessiva è quella di entrambi.

Nelle misure i processi CPU sono lanciati attraverso lo script
\texttt{mem\_locale.sh}, che con \texttt{numactl -{}-membind} vincola la
memoria di ogni processo al proprio socket: senza il vincolo, se la memoria
locale è occupata dalla cache dei file di altri utenti, Linux alloca
sull'altro socket, e la banda misurata scende da 71 a \SI{52}{GB/s}.

\paragraph{Niente hyper-threading.} Si usano solo i core fisici 0--19, uno per processo, mentre i core logici 20--39 condividono con questi unità di calcolo e cache
(\cref{sec:cpu-hw}). A $k$ piccolo un guadagno è escluso, perché 20 processi
saturano già la banda di memoria (\cref{sec:ris-mpi}).
% NOTA: a k grande l'effetto dell'hyper-threading non e' stato misurato
% (cpu_block e' a circa meta' del tetto: potrebbe guadagnare o perdere, per la
% L1/L2 condivise e i blocchi piu' piccoli con 40 processi). Prova possibile:
% mpirun -np 40 --use-hwthread-cpus --bind-to hwthread ...

\paragraph{Rank prevedibili.} Il binding funziona come previsto anche grazie a
\texttt{reorder = 0} nella creazione della griglia (\cref{sec:griglia}): MPI
non rinumera i processi, quindi il rank nella griglia coincide con quello
assegnato da \texttt{mpirun}, e il legame tra posizione nella griglia e core
resta noto.

Nella versione CUDA invece i processi si lanciano senza opzioni di binding e passano
quasi tutto il tempo ad attendere la GPU, e la posizione sui core della CPU ha
un effetto trascurabile sulle prestazioni.

\subsection{Confronto con la GPU}\label{sec:cpu-gpu}

Nello stesso caso \texttt{warp\_row\_xt} in \texttt{double} va da circa 6 volte
più veloce di \texttt{cpu\_block} a $k = 3$ (298 contro 49~GFLOPS) a 1,3 volte a
$k = 32$ (352 contro 271). Il rapporto si riduce al crescere di $k$, perché la
GPU in \texttt{double} è già ferma al picco di calcolo, mentre la CPU continua a
salire con la memoria.

% =====================================================================
\section{Metodologia di misura}\label{sec:misure}

\subsection{Fasi misurate}\label{sec:fasi}

A ogni invocazione del nucleo, ogni processo misura separatamente le fasi
seguenti con l'orologio di MPI (\texttt{MPI\_Wtime}). Tra parentesi è
indicata la colonna corrispondente nel CSV prodotto dal programma.
\begin{itemize}
  \item \emph{Broadcast}: \texttt{MPI\_Bcast} di $X$ lungo la colonna (CSV:
        \texttt{t\_bcast}).
  \item \emph{Preparazione di $X$}: \texttt{be\_set\_x}, che sulla GPU comprende la copia verso
        il dispositivo ed eventuale trasposizione (CSV: \texttt{t\_h2d}).
  \item \emph{Calcolo}: \texttt{be\_compute}, dalla chiamata al ritorno (CSV:
        \texttt{t\_compute\_mean}).
  \item \emph{Restituzione di $Y$}: \texttt{be\_get\_y}, che sulla GPU comprende la copia verso
        la CPU (CSV: \texttt{t\_d2h}).
  \item \emph{Riduzione}: \texttt{MPI\_Reduce} di $Y$ lungo la riga (CSV:
        \texttt{t\_reduce}).
\end{itemize}

Il \emph{tempo del nucleo} di un processo è la somma delle tre fasi incluse
dalla traccia,
\[
  t_{\text{nucleo}} = t_{\text{broadcast}} + t_{\text{calcolo}} + t_{\text{riduzione}}
  \qquad (\text{CSV: } \texttt{t\_core\_mean}),
\]
ed è il tempo con cui si calcolano i GFLOPS (\cref{sec:metrica}, CSV:
\texttt{gflops\_core}). Il programma ne calcola anche una versione sul solo
calcolo, per valutare i kernel indipendentemente dalle comunicazioni (CSV:
\texttt{gflops\_compute}).

Vengono inoltre riportate tre misure accessorie:
\begin{itemize}
  \item il \emph{totale delle fasi locali}: la durata complessiva di
        preparazione di $X$, calcolo e restituzione di $Y$, cioè di tutto ciò
        che avviene tra il broadcast e la riduzione (CSV:
        \texttt{t\_local\_wall});
  \item il \emph{calcolo misurato dal backend}: sulla GPU con gli eventi CUDA
        (\cref{sec:h2d}), sulla CPU con l'orologio. Il confronto con il calcolo
        a orologio è discusso in \cref{sec:eventi-orologio} (CSV:
        \texttt{t\_compute\_dev});
  \item il \emph{preprocessamento}, misurato una sola volta prima delle
        invocazioni (\cref{sec:dati}, CSV: \texttt{t\_pre}).
\end{itemize}

Come già introdotto, il calcolo si misura con l'orologio della CPU prima e dopo
\texttt{be\_compute}, che ritorna solo a calcolo terminato: l'intervallo
comprende così tutto ciò che il processo attende per avere il risultato,
compresa l'eventuale attesa della GPU occupata da altri processi
(\cref{sec:eventi-orologio}). Ogni invocazione inizia con una barriera
(\texttt{MPI\_Barrier}), così che tutti i processi partano insieme e i loro
tempi siano confrontabili.

\subsection{Aggregazione statistiche}\label{sec:statistiche}

Con più processi, ogni invocazione produce un tempo per ciascun processo, e
occorre decidere come ridurli a un solo numero.

\paragraph{Prima il massimo, poi le statistiche.} Un'invocazione del nucleo è
conclusa solo quando ha terminato il processo più lento: per ogni invocazione,
e separatamente per ogni fase, si prende quindi il \emph{massimo tra i
processi}. Le statistiche si calcolano dopo, sulla serie di questi massimi, e
la loro media è il tempo medio per invocazione richiesto dalla traccia.

L'ordine conta: facendo prima la media di ogni processo e poi il massimo tra le
medie, si sottostimerebbe il tempo quando il processo più lento cambia da
un'invocazione all'altra. Per esempio, con due processi e due invocazioni.

\begin{table}[H]
  \centering
  \caption{Ordine delle operazioni nell'aggregazione dei tempi: il massimo per
    invocazione, poi la media.}\label{tab:aggregazione}
  \begin{tabular}{lccc}
    \toprule
     & Invocazione 1 & Invocazione 2 & Media del processo \\
    \midrule
    Processo A & 10 & 6  & 8 \\
    Processo B & 6  & 10 & 8 \\
    \midrule
    \textbf{Massimo dell'invocazione} & \textbf{10} & \textbf{10} & \\
    \bottomrule
  \end{tabular}
\end{table}

Il massimo delle medie darebbe 8, la media dei massimi 10, che è il tempo
reale.

Il massimo si calcola con un'unica \texttt{MPI\_Reduce} con operazione
\texttt{MPI\_MAX}, applicata elemento per elemento all'intera tabella dei tempi
(fasi $\times$ invocazioni) di ogni processo.

\begin{lstlisting}[caption={Massimo tra i processi dell'intera tabella dei tempi (\texttt{main.c}).}, label={lst:reduce-max}]
MPI_Reduce(g.rank == 0 ? MPI_IN_PLACE : ph, g.rank == 0 ? ph : NULL,
           MXM_NTIMES * o.reps, MPI_DOUBLE, MPI_MAX, 0, g.grid);
\end{lstlisting}
\texttt{ph} (da \emph{phases}, fasi) è la tabella dei tempi del processo, con
una riga per ogni fase e una colonna per ogni invocazione. Sul processo~0
\texttt{MPI\_IN\_PLACE} fa usare la stessa tabella come ingresso e come
risultato: non serve un secondo buffer.

\paragraph{Statistiche riportate.} Sulla serie delle $R$ invocazioni misurate
si calcolano:
\begin{itemize}
  \item la \emph{media}, usata per i GFLOPS;
  \item la \emph{mediana}, meno sensibile della media a singole invocazioni
        anomale;
  \item il \emph{minimo}, il miglior tempo ottenuto;
  \item la \emph{deviazione standard}, campionaria, cioè con divisore $R - 1$;
  \item il \emph{coefficiente di variazione} (CV), rapporto percentuale tra
        deviazione standard e media: misura quanto le invocazioni sono stabili
        ed è confrontabile tra configurazioni con tempi molto diversi.
\end{itemize}
Le statistiche complete si riportano per il tempo del nucleo e per il calcolo, mentre per le altre fasi basta la media.

\subsection{Eventi CUDA e orologio}\label{sec:eventi-orologio}

Sulla GPU il tempo di calcolo si può misurare in due modi, e il programma li
riporta entrambi:
\begin{itemize}
  \item \emph{a orologio}: \texttt{MPI\_Wtime} prima e dopo
        \texttt{be\_compute}, sulla CPU, cioè il tempo che il processo attende per
        avere il risultato (CSV: \texttt{t\_compute\_mean});
  \item \emph{con gli eventi CUDA}: il tempo trascorso sulla GPU tra l'inizio e
        la fine del kernel (\cref{sec:h2d}, CSV: \texttt{t\_compute\_dev}).
\end{itemize}
Queste misure differiscono a seconda del caso di utilizzo della GPU.

\paragraph{GPU in uso esclusivo.} Con un solo processo le due misure
differiscono solo per il costo del lancio e della sincronizzazione, pochi
microsecondi.

\paragraph{GPU condivisa: il time-slicing.} Con più processi sulla stessa GPU, i
kernel di processi diversi non vengono eseguiti insieme: la GPU li esegue a
turno, passando dall'uno all'altro (\emph{time-slicing}). Un processo che lancia
il proprio kernel mentre la GPU è occupata deve quindi attendere il proprio
turno, e le due misure si separano:
\begin{itemize}
  \item l'orologio comprende l'attesa del turno, perché parte quando il
        processo chiede il calcolo;
  \item gli eventi CUDA non la vedono, perché il primo evento viene registrato
        solo quando la GPU arriva al processo.
\end{itemize}

Per questo il tempo del nucleo usa la misura a orologio: descrive ciò che
accade davvero su questo server, dove i processi si contendono un'unica GPU.
Gli eventi CUDA restano una misura informativa: escludono l'attesa prima che
il kernel cominci, ma non i turni ceduti agli altri processi durante
l'esecuzione, e quindi non corrispondono al tempo che il calcolo avrebbe con la
GPU tutta per sé, che si stima invece con la modalità a nodi emulati
(\cref{sec:modalita}). La differenza tra le due misure dà una stima diretta del
tempo perso nell'attesa iniziale.

Un possibile miglioramento ci sarebbe adattando \emph{Multi-Process Service} (MPS) di NVIDIA,
che permette ai kernel di processi diversi di essere eseguiti
contemporaneamente sulla stessa GPU, invece che a turno, riducendo l'attesa dei processi. In questo progetto non è stato utilizzato perché non disponibile
sul server.

\subsection{Riscaldamento}\label{sec:riscaldamento}

Le prime invocazioni del nucleo vengono eseguite ma non registrate. Questo
riscaldamento serve a due scopi:
\begin{itemize}
  \item \emph{pagare i costi una tantum}: cache vuote, caricamento del codice
        dei kernel sulla GPU, inizializzazione dei canali di comunicazione di MPI: sono costi che un'applicazione reale paga una sola volta, non a
        ogni invocazione;
  \item \emph{portare la macchina a regime}: GPU e CPU adattano la frequenza al
        carico, e la alzano solo dopo un periodo di lavoro continuo
        (\cref{sec:gpu-hw}).
\end{itemize}

\paragraph{Un numero fisso di invocazioni non basta.} Le dimensioni misurate
vanno da $2048 \times 2048$ a $16384 \times 16384$: alla più piccola due
invocazioni durano una frazione di millisecondo, alla più grande cento volte di
più. Un numero fisso di invocazioni di riscaldamento per qualsiasi dimensione non sarebbe la best-practice.
Sul server la stessa
configurazione dava \SI{286}{GFLOPS} o \SI{340}{GFLOPS} con due invocazioni di
riscaldamento, e sempre circa \SI{339}{GFLOPS} con circa due secondi di
riscaldamento.

Per questo il riscaldamento dura almeno un certo numero di invocazioni
(opzione \texttt{-W}) \emph{e} almeno un certo tempo (opzione \texttt{-T}, un
secondo per default): ogni configurazione parte così dallo stesso stato della
macchina. In pratica il riscaldamento esegue tutte le invocazioni che stanno in
\texttt{-T} secondi, mentre \texttt{-W} garantisce un numero minimo quando una
singola invocazione è lunga.

\paragraph{Una decisione unica per tutti.} È \emph{il processo~0 a decidere}
quando il riscaldamento è finito, e \emph{comunica la decisione agli altri con
un \texttt{MPI\_Bcast}} dopo ogni invocazione di riscaldamento
(\cref{lst:warmup}). Il riscaldamento prosegue finché non sono state eseguite
almeno \texttt{-W} invocazioni e non sono trascorsi almeno \texttt{-T} secondi,
misurati sull'orologio del processo~0. Nonostante la condizione è calcolata da tutti i
processi, il broadcast la sostituisce con quella del processo~0: tutti
escono dal riscaldamento alla stessa invocazione. Questo è necessario perché
ogni invocazione contiene comunicazioni collettive, e tutti i processi devono
eseguirne lo stesso numero. Il numero di invocazioni di riscaldamento
effettivamente eseguite è riportato nell'output (CSV: \texttt{n\_warmup}).

\begin{lstlisting}[caption={Fine del riscaldamento decisa dal processo 0 (\texttt{main.c}).}, label={lst:warmup}]
if (warming) {
    n_warm++;
    int more = n_warm < o.warmup || MPI_Wtime() - tw0 < o.warm_time;
    MPI_Bcast(&more, 1, MPI_INT, 0, g.grid);
    warming = more;
    continue;
}
\end{lstlisting}

\subsection{Modalità concorrente e nodi emulati}\label{sec:modalita}

La fase locale si può eseguire in due modi, scelti con l'opzione \texttt{-S}.

\paragraph{Concorrente (predefinita).} Tutti i processi calcolano insieme e si
contendono le risorse condivise della macchina: sulla CPU la banda di memoria e
le cache, nella versione CUDA l'unica GPU (\cref{sec:eventi-orologio}). Questa
modalità misura le prestazioni reali di questo server.

\paragraph{Serializzata, o a nodi emulati (\texttt{-S}).} I processi calcolano a
turno, uno alla volta, mentre gli altri attendono a una barriera. Durante il
proprio turno ogni processo ha la macchina tutta per sé, come se si trovasse su
un nodo separato: il tempo del nucleo misura quindi il tempo che lo stesso
algoritmo impiegherebbe su $P$ nodi indipendenti, senza contesa. Serve perché sul
server c'è una sola GPU, che i processi si dividono (\cref{sec:gpu-assign}): la
modalità \texttt{-S} mostra come scalerebbe la stessa soluzione con una GPU per
processo, e quanto conta la forma della griglia senza l'effetto della contesa.
Lo stesso vale per la CPU,
dove i processi si dividono la banda di memoria.

Nel codice (\cref{lst:turni}), in modalità concorrente tutti i processi
eseguono insieme il prodotto locale. In modalità \texttt{-S} calcolano invece a
turno, in ordine di rank: al primo giro calcola il processo~0 mentre gli altri
attendono alla barriera, al secondo il processo~1, e così via fino all'ultimo.
Ogni processo calcola una sola volta per invocazione, e il suo turno dura
quanto il suo prodotto locale. È simile a un giro di \emph{round robin}, ma
senza un quanto di tempo: il turno non viene interrotto.

\begin{lstlisting}[caption={Esecuzione a turni della fase locale (\texttt{mxm.c}).}, label={lst:turni}]
if (!h->serial) {
    local_product(h, Xc, t);                 /* tutti insieme */
} else {
    /* A turno: calcola il processo di rank "turn", gli altri lo
       aspettano alla barriera; poi tocca al successivo. */
    for (int turn = 0; turn < g->size; turn++) {
        if (g->rank == turn)
            local_product(h, Xc, t);
        MPI_Barrier(g->grid);
    }
}
\end{lstlisting}

La stima rimane comunque otttimistica broadcast e riduzione passano
comunque per la memoria condivisa del server, molto più veloce di una rete tra
nodi veri.

\subsection{Precisione e indici di prestazione}\label{sec:modelli}

\paragraph{Confronto tra \texttt{float} e \texttt{double}.} Il \texttt{float} è più veloce,
ma meno preciso: conserva circa 7 cifre significative, contro le circa 16 del
\texttt{double}. L'errore relativo atteso del prodotto è
\[
  \text{errore} \approx \varepsilon \sqrt{N} .
\]
$\varepsilon$ (\emph{epsilon di macchina}) è il più piccolo errore relativo che
si commette arrotondando un numero: \num{1.2e-7} in \texttt{float},
\num{2.2e-16} in \texttt{double}. Ogni elemento di $Y$ è una somma di $N$
termini, e a ogni somma si aggiunge un piccolo errore di arrotondamento ma questi
errori hanno segni casuali e in parte si compensano, quindi crescono come
$\sqrt{N}$ e non come $N$. Con $N = 16384$, $\sqrt{N} = 128$, e l'errore atteso
è circa \num{1.5e-5} in \texttt{float} e \num{2.8e-14} in \texttt{double}
(\cref{sec:verifica}).
Quale delle due precisioni usare è a discrezione dell'utente: 
il programma offre entrambe le versioni. Le velocità misurate sono confrontate
in \cref{sec:gpu-riepilogo}, gli errori in \cref{sec:verifica}.

\paragraph{Speedup ed efficienza.} Per capire quanto conviene aggiungere
processi si confronta il tempo con $P$ processi con quello con un solo
processo:
\[
  S(P) = \frac{T(1)}{T(P)}, \qquad E(P) = \frac{S(P)}{P} .
\]
$T(1)$ è il tempo del nucleo con un solo processo e $T(P)$ quello con $P$
processi. Lo \emph{speedup} $S(P)$ dice quante volte si va più veloci.
L'\emph{efficienza} $E(P)$ dice quanto bene si sfruttano i
processi: è lo speedup diviso per il numero di processi.

Se la versione a un processo è scritta o configurata male, è lenta,
e qualunque versione parallela sembra ottima al confronto, si
può persino ottenere uno speedup superiore a $P$, che sembra impossibile. Per
evitarlo, il riferimento a un processo usa lo stesso kernel, configurato al
meglio: ad esempio, per il budget dei registri di \texttt{cpu\_block} si è
scelto 128 invece di 64 questo perchè, nonostante con 20 processi i due valori si equivalgano
con un solo processo 128 è nettamente più veloce (\cref{sec:ris-padding}).
Uno speedup superlineare (>$P$) può comunque avere cause reali come il caso in cui con
più processi ogni blocco è più piccolo e può stare meglio in
cache.Ad ogni modo il caso osservato è riportato in \cref{sec:ris-mpi}.

% =====================================================================
\section{Verifica della correttezza}\label{sec:verifica}

\subsection{Riferimento seriale e verifica distribuita}\label{sec:verifica-distr}

Il riferimento è la funzione \texttt{serial\_mxm} (\texttt{serial.c}): tre
cicli annidati, con gli accumulatori sempre in \texttt{double}
(\cref{sec:collaudo}).

Come anticipato in \cref{sec:collaudo}, la verifica è \emph{distribuita}: non
raccoglie mai $A$ o $Y$ su un solo processo, perché per le matrici più grandi
non entrerebbero nella memoria di un processo solo. La esegue
\texttt{verify} in \texttt{main.c}: ogni processo della colonna~0, che alla
fine possiede le proprie righe di $Y$, le ricontrolla da sé, rigenerando $A$
una riga alla volta grazie al fatto che il generatore dipende solo dagli indici
globali (\cref{sec:generazione}). In memoria restano così solo una riga di $A$ e
$X$ intera (\cref{lst:verify}).

\begin{lstlisting}[caption={Verifica delle righe di $Y$ possedute dal processo (\texttt{main.c}).}, label={lst:verify}]
for (int r = 0; r < b->Mloc; r++) {
    fill_block(arow, o->N, b->m_off + r, 1, 0, o->N, SEED_A);
    serial_mxm(1, o->N, o->k, arow, o->N, Xfull, o->k, yref, o->k);
    for (int c = 0; c < o->k; c++) {
        double d = (double)Y[IDX(r, c, o->k)] - yref[c];
        loc[0] += d * d;
        loc[1] += yref[c] * yref[c];
    }
}
\end{lstlisting}

Per ogni riga:
\begin{enumerate}
  \item si rigenera la riga di $A$ con lo stesso generatore usato per il
        calcolo;
  \item si calcola la riga corretta di $Y$ con il riferimento seriale;
  \item si accumulano due somme: quella dei quadrati delle differenze tra il
        risultato del programma e quello di riferimento, e quella dei quadrati
        dei valori di riferimento.
\end{enumerate}
Alla fine una sola \texttt{MPI\_Allreduce} somma i contributi di tutti i
processi (quelli fuori dalla colonna~0 contribuiscono con zero), e dalle due
somme si ricava l'errore relativo (\cref{sec:errore}).

\subsection{Misura dell'errore}\label{sec:errore}

\paragraph{Errore relativo in norma di Frobenius.} Il programma riassume la
differenza tra il risultato $Y$ e il riferimento $Y_{\text{ref}}$ in un solo
numero:
\[
  \text{errore} = \frac{\lVert Y - Y_{\text{ref}} \rVert_F}{\lVert Y_{\text{ref}} \rVert_F} .
\]
La \emph{norma di Frobenius} di una matrice è la radice quadrata della somma dei
quadrati di tutti i suoi elementi.

\paragraph{Tolleranza ed Errori:} Si è deciso di accettare il risultato se
\[
  \text{errore} \le 10\, \varepsilon \sqrt{N} ,
\]
dove si ricorda che $\varepsilon\sqrt{N}$ è l'errore atteso per una somma di $N$ termini
(\cref{sec:modelli}).Per istanza con $N = 16384$ la
tolleranza vale circa \num{2.8e-13} in \texttt{double} e \num{1.5e-4} in
\texttt{float}.

Oltre all'esito, il programma stampa il rapporto $\text{errore} / (\varepsilon\sqrt{N})$, cioè quante volte l'errore
misurato supera quello atteso, lo script di collaudo ne riporta il massimo su
tutte le prove (\cref{sec:batteria}).
Sulla batteria di test l'errore misurato non supera
$0{,}161\,\varepsilon\sqrt{N}$ in doppia precisione e
$0{,}150\,\varepsilon\sqrt{N}$ in singola: la tolleranza
$10\,\varepsilon\sqrt{N}$ ha quindi un margine di oltre 60 volte. La costante 10
è dunque, sicuramente molto prudente, ma lascia spazio a configurazioni con errori più alti di quelle
provate senza dare falsi allarmi.

Il modo in cui si sommano i termini influisce sull'errore:
\begin{itemize}
  \item in \texttt{thread\_row} ogni thread somma gli $N$ termini di una riga
        in fila, uno dopo l'altro, e ogni nuovo termine si aggiunge a un totale
        sempre più grande, accumulando insieme anche gli errori di arrotondamento;
  \item nei kernel a warp la riga è divisa tra 32 lane: ognuna somma solo
        $N/32$ termini, e le 32 somme parziali vengono poi combinate ad albero,
        in 5 passi (\cref{sec:warp-row}).
\end{itemize}
Nel secondo caso le somme sono più corte e combinano numeri di grandezza simile: ci si aspetta
quindi un errore di arrotondamento più piccolo. I dati della batteria di test
lo confermano in \texttt{float}, dove il riferimento in \texttt{double} è più
accurato dei kernel: il rapporto $\text{errore}/(\varepsilon\sqrt{N})$ vale
circa 0,149 per \texttt{thread\_row} e 0,030 per i kernel a warp. In
\texttt{double} il confronto non è possibile: \texttt{thread\_row} somma nello
stesso ordine del riferimento seriale e ne riproduce esattamente il risultato,
con errore misurato nullo.

\subsection{Batteria di test}\label{sec:batteria}

Lo script \texttt{test/run\_tests.sh} esegue l'intero collaudo funzionale con un
solo comando. Ogni prova lancia il programma con la verifica attiva (opzione
\texttt{-C}) e ne controlla l'esito, e alla fine lo script stampa il numero di
prove riuscite e fallite e, per ciascuna precisione, il massimo del rapporto
$\text{errore}/(\varepsilon\sqrt{N})$ (\cref{sec:errore}).

Le prove sono 166, in quattro fasi:
\begin{enumerate}
  \item \emph{compilazione} delle quattro versioni (CPU/GPU $\times$
        \texttt{double}/\texttt{float});
  \item \emph{versione CPU, solo MPI} (81 prove): entrambi i kernel, griglie da
        1 a 6 processi di forme diverse (comprese $4 \times 1$ e $1 \times 4$),
        $k = 3, 8, 13, 32$, entrambe le precisioni, la modalità a nodi emulati
        e l'esecuzione senza padding;
  \item \emph{versione GPU con un processo} (65 prove): tutti i kernel,
        compreso cuBLAS, entrambe le precisioni, tutti i $k$ della traccia più
        $k = 13$, con e senza padding;
  \item \emph{versione GPU con MPI} (20 prove): i tre kernel a warp su griglie
        da 2 a 6 processi, anche in modalità a nodi emulati.
\end{enumerate}

\paragraph{Dimensioni scelte per mettere in difficoltà gli indici.} Le matrici
di prova hanno dimensioni volutamente ``scomode'':
\begin{itemize}
  \item $301 \times 257$ per la CPU;
  \item $2000 \times 1500$ e $3001 \times 2003$ per la GPU.
\end{itemize}
Non sono multipli di 32 (la dimensione del warp), né del numero di processi, né
della dimensione dei blocchi e dei tile dei kernel. Ogni prova esercita così i
casi limite:
\begin{itemize}
  \item l'ultimo blocco più piccolo degli altri;
  \item i warp senza riga dell'ultimo blocco (\cref{sec:warp-smem});
  \item la ripartizione non esatta tra i processi (\cref{sec:distribuzione});
  \item il padding che non coincide con la fine della riga (\cref{sec:padding});
  \item il percorso per $k$ generico, con fette da 8 colonne più il resto
        (\cref{sec:warp-row}), controllato con $k = 13$, che non è tra i
        valori della traccia.
\end{itemize}

Lo script verifica la correttezza e non le prestazioni per questo disattiva il
riscaldamento a tempo (\texttt{-T 0}), che allungherebbe solo il collaudo.
Tutte le 166 prove sono superate.

Il \cref{lst:test} riporta un estratto dell'output della batteria, eseguita sul
server il 6 novembre 2026.

\begin{lstlisting}[language={}, basicstyle=\ttfamily\scriptsize, numbers=none, caption={Estratto dell'output di \texttt{run\_tests.sh} (\texttt{...} = righe omesse; spazi di allineamento ridotti).}, label={lst:test}]
== 1. compilazione
  OK  make BACKEND=cpu PREC=double
  ...
== 2. versione CPU, solo MPI (M=301, N=257)
  OK  kernel=0 P=1 griglia 1x1 k=3  errore relativo 0.000e+00 (tolleranza 3.6e-14, errore/(eps*sqrt(N)) = 0.000) -> OK
  OK  kernel=0 P=4 griglia 1x4 k=3  errore relativo 5.721e-16 (tolleranza 3.6e-14, errore/(eps*sqrt(N)) = 0.161) -> OK
  OK  float kernel=0 P=4 griglia 2x2 k=3  errore relativo 2.110e-07 (tolleranza 1.9e-05, errore/(eps*sqrt(N)) = 0.110) -> OK
  ...
== 3. versione GPU, 1 processo (M=2000, N=1500)
  OK  double kernel=2 k=3  errore relativo 1.329e-15 (tolleranza 8.6e-14, errore/(eps*sqrt(N)) = 0.154) -> OK
  OK  float kernel=0 k=13  errore relativo 6.912e-07 (tolleranza 4.6e-05, errore/(eps*sqrt(N)) = 0.150) -> OK
  ...
== 4. versione GPU + MPI, warp_row, warp_row_xt, warp_row_smem (M=3001, N=2003, k=8)
  OK  kernel=2 P=4 griglia 2x2  errore relativo 1.514e-15 (tolleranza 9.9e-14, errore/(eps*sqrt(N)) = 0.152) -> OK
  OK  modalità -S P=4 griglia 2x2  errore relativo 1.514e-15 (tolleranza 9.9e-14, errore/(eps*sqrt(N)) = 0.152) -> OK
  ...
== riepilogo: 166 OK, 0 FALLITE
  errore/(eps*sqrt(N)) massimo: double 0.161, float 0.150 (tolleranza: 10)
\end{lstlisting}

% =====================================================================
\section{Campagne sperimentali}\label{sec:campagne}

\subsection{Gruppi di esperimenti e automazione}\label{sec:gruppi}\label{sec:automazione}

Le misure sono organizzate in gruppi, ciascuno pensato per rispondere a una
domanda. Ogni gruppo è descritto da un \emph{file di configurazione} in
\texttt{bench/config/} ed eseguito da un unico script,
\texttt{bench/campagna.sh}: una campagna si ripete così con un solo comando,
sempre con gli stessi parametri. La \cref{tab:conf} riporta i gruppi
nell'ordine in cui vengono eseguiti, insieme a due verifiche mirate che hanno
uno script proprio.

\begin{table}[H]
  \centering
  \small
  \caption{Gruppi di esperimenti, file che li descrivono e numero di esecuzioni.}\label{tab:conf}
  \renewcommand{\arraystretch}{1.25}
  \begin{tabularx}{\textwidth}{|c|l|X|r|}
    \hline
    \textbf{Gruppo} & \textbf{File} & \textbf{Descrizione} & \textbf{Esec.} \\
    \hline
    B & \texttt{B\_kernel.conf} &
      \textbf{Quanto vale ogni ottimizzazione?} $16384 \times 16384$, tutti i
      $k$ della traccia. In \texttt{double} tutti i kernel GPU, cuBLAS e i due
      kernel CPU con 20 processi. In \texttt{float} \texttt{warp\_row\_xt},
      \texttt{warp\_row\_smem} e cuBLAS. Si lancia per primo, perché decide il
      kernel GPU usato dagli altri (risultati nelle \cref{sec:gpu,sec:cpu}). & 50 \\
    \hline
    A & \texttt{A\_finale.conf} &
      \textbf{Quanto va il progetto sulle dimensioni della traccia?} Kernel GPU
      migliore, \texttt{double}, dimensioni da 2048 a 16384, le tre forme e
      tutti i $k$. & 60 \\
    \hline
    C & \texttt{banda.sh}, \texttt{picco\_gpu.cu} &
      \textbf{Quali sono i tetti della roofline?} Banda di lettura di GPU e CPU e
      picco di calcolo della GPU in \texttt{double} e \texttt{float}
      (\cref{sec:banda}). & 4 \\
    \hline
    D & \texttt{D\_mpi\_cpu.conf} &
      \textbf{Come scala la versione CPU?} $16384 \times 16384$,
      \texttt{cpu\_block}, $k = 3$ e $32$, 11 griglie da 1 a 20 processi
      (comprese $20 \times 1$ e $1 \times 20$), in modalità concorrente e a
      nodi emulati. & 42 \\
    \hline
    E & \texttt{E\_mpi\_gpu.conf} &
      \textbf{Come scalerebbe la versione MPI + CUDA con nodi emulati, e quanto
      conta la forma della griglia?} $16384 \times 16384$, kernel migliore,
      $k = 8$, griglie da 1 a 4 processi, in modalità concorrente e a nodi
      emulati. & 11 \\
    \hline
    F & \texttt{F\_padding.conf} &
      \textbf{Quanto aiuta allineare le righe?} Dimensioni 8191 e 16383, scelte
      perché senza padding le righe partano disallineate, con e senza padding,
      su GPU e CPU. & 24 \\
    \hline
    --- & \texttt{budget\_cpu.sh} &
      Verifica mirata: budget dei registri 32, 64 e 128 di \texttt{cpu\_block},
      con 1 e con 20 processi (\cref{sec:vettorizzazione}). & 30 \\
    \hline
    --- & \texttt{frequenza.sh} &
      Verifica mirata: la stessa configurazione con riscaldamento breve e lungo,
      per l'effetto della frequenza variabile (\cref{sec:riscaldamento}). & 18 \\
    \hline
    --- & \texttt{migliore.conf} &
      Kernel GPU migliore, aggiornato dopo il gruppo B e letto da A, E e F. & --- \\
    \hline
  \end{tabularx}
\end{table}

\begin{lstlisting}[language=bash,numbers=none,frame=none]
bash bench/campagna.sh bench/config/A_finale.conf
\end{lstlisting}

\paragraph{Il file di configurazione.} Contiene una lista di valori per ogni
parametro:
\begin{itemize}
  \item precisioni;
  \item dimensioni;
  \item forme;
  \item valori di $k$;
  \item kernel;
  \item griglie;
  \item modalità;
  \item padding.
\end{itemize}
Ad esempio il \cref{lst:conf-a} riporta il cuore della configurazione del gruppo
A. Le esecuzioni sono tutte le combinazioni dei valori elencati: per il gruppo
A, 4 dimensioni $\times$ 3 forme $\times$ 5 valori di $k$, cioè 60.

\begin{lstlisting}[language=bash, caption={Parametri della campagna A (\texttt{A\_finale.conf}, estratto).}, label={lst:conf-a}]
BACKENDS="cuda"
PRECS="double"
SIZES="2048 4096 8192 16384"
SHAPES="square m3n n2m"
KS="3 6 8 20 32"
KERNELS_cuda="$KERNEL_GPU_MIGLIORE"
LAUNCH_cuda="1:1:1"          # processi : righe : colonne della griglia
REPS=30
\end{lstlisting}

\paragraph{Cosa fa lo script.} Dato il file di configurazione,
\texttt{campagna.sh}:
\begin{enumerate}
  \item elenca tutte le combinazioni da eseguire. Con \texttt{ELENCO=1} si
        ferma qui: stampa l'elenco senza compilare né eseguire nulla, e serve a
        controllare una configurazione prima di lanciarla;
  \item compila solo gli eseguibili necessari, cioè le combinazioni di backend
        e precisione presenti nella configurazione;
  \item esegue le misure una alla volta. Se una fallisce, annota nel log l'errore e passa alla successiva, senza
        interrompere la campagna;
  \item salva i risultati in un unico CSV in \texttt{results/}, con il nome
        della campagna e la data, insieme a una copia della configurazione e al
        log degli errori.
\end{enumerate}

\paragraph{Il kernel GPU migliore} usato dai gruppi
A, E ed F, non è hard-coded in ciascuna configurazione, ma letto da un file unico,
\texttt{migliore.conf}, che viene aggiornato dopo l'esecuzione del gruppo B. La copia della
configurazione salvata con i risultati riporta anche il kernel effettivamente
usato: ogni CSV resta così riproducibile, anche se in seguito
\texttt{migliore.conf} dovesse cambiare per motivi diversi.

Le campagne più lunghe si lanciano con
\texttt{nohup \dots\ \&}, così da proseguire anche se la connessione SSH cade.

\subsection{Scelta delle dimensioni}\label{sec:dimensioni}

I gruppi B, D ed E usano tutti
la matrice quadrata $16384 \times 16384$, perché è molto più grande delle cache:
occupa \SI{2}{GiB}, contro i \SI{4}{MiB} della cache L2 della GPU e i
\SI{13.75}{MiB} della L3 di ogni processore. Ogni invocazione è quindi costretta a leggere $A$
davvero dalla memoria, così si misura il comportamento reale del nucleo, senza la facilitazione di una matrice che sta in cache.

Il gruppo A fa crescere la dimensione di base $S$ raddoppiandola (2048, 4096,
8192 e 16384). Si coprono così quasi due ordini di grandezza nel numero di
elementi di $A$ (da circa 4 a 268 milioni nella forma quadrata): in questo modo
si osserva come cambiano le prestazioni passando da matrici piccole, con poco
lavoro per la GPU, a matrici grandi. Le tre forme si ottengono da $S$:
\begin{itemize}
  \item quadrata: $M = N = S$;
  \item alta: $M = 3S$ e $N = S$;
  \item larga: $M = S$ e $N = 2S$.
\end{itemize}

Il tempo usato è la media
delle invocazioni misurate. Se una sola invocazione viene disturbata, ad
esempio da un altro programma sul server o dal sistema operativo, il ritardo
viene assorbito nella media, mentre con le matrici grandi il problema quasi non si pone,
questo può invece sfalzare le run per le dimensioni piccole,
con invocazioni di pochi millisecondi, per far fronte a questo si misurano 30
invocazioni per configurazione invece di 10.

Il gruppo F usa invece 8191
e 16383, cioè $2^p - 1$: con queste dimensioni, senza padding, le righe
cominciano a tutti i possibili disallineamenti, mentre con 8192 o 16384
sarebbero già allineate per conto loro e il padding non avrebbe effetto
(\cref{sec:padding}).

% =====================================================================
\section{Risultati}\label{sec:risultati}

Durante lo studio si è constatato, in sintesi:
\begin{itemize}
  \item sulla GPU il kernel migliore è \texttt{warp\_row\_xt}: in
        \texttt{double} raggiunge il \SI{98}{\percent} del tetto della memoria a
        $k = 3$ e il \SIrange{94}{95}{\percent} del picco di calcolo per $k \ge 6$. In \texttt{float} segue il tetto della memoria fino a $k = 8$, mentre a $k = 20$ e $k = 32$ cuBLAS è 1,6 e 2,3 volte più
        veloce;
  \item sulla CPU, come prevedibile, \texttt{cpu\_block}, con 20 processi, arriva tra
        l'\SI{84}{\percent} e il \SI{92}{\percent} del tetto della memoria fino
        a $k = 20$.
\end{itemize}
In questo capitolo vengono presentate le prestazioni della soluzione finale, la scalabilità
con MPI, l'effetto del padding e del budget dei registri, e gli effetti sulla
misura.

\subsection{La soluzione finale}\label{sec:ris-finale}

La campagna~A misura la soluzione (\emph{reale}) ottimale trovata sull'hardware fornito, in questo caso \texttt{warp\_row\_xt} in
\texttt{double} con un processo, sulle dimensioni richieste dalla traccia: la
dimensione di base $S$ cresce da 2048 a 16384, nelle tre forme, per tutti i $k$
(\cref{tab:ris-finale,fig:ris-finale}). I valori sono i GFLOPS del tempo del
nucleo.

La soluzione finale usa la doppia precisione, con un errore relativo atteso
dell'ordine di $10^{-14}$. La singola è da 2 a 5 volte più veloce
(\cref{tab:ris-kernel}), ma con un errore atteso dell'ordine di $10^{-5}$
(\cref{sec:modelli}).

\begin{table}[H]
  \centering
  \caption{GFLOPS della soluzione finale (\texttt{warp\_row\_xt},
    \texttt{double}, un processo, campagna~A).}\label{tab:ris-finale}
  \begin{tabular}{rllrrrrr}
    \toprule
    $S$ & Forma & $M \times N$ & $k=3$ & $k=6$ & $k=8$ & $k=20$ & $k=32$ \\
    \midrule
    2048  & quadrata & $2048 \times 2048$   & 222 & 261 & 265 & 274 & 277 \\
          & larga    & $2048 \times 4096$   & 249 & 289 & 291 & 294 & 299 \\
          & alta     & $6144 \times 2048$   & 260 & 302 & 305 & 310 & 310 \\
    \midrule
    4096  & quadrata & $4096 \times 4096$   & 274 & 319 & 320 & 323 & 325 \\
          & larga    & $4096 \times 8192$   & 287 & 335 & 334 & 337 & 340 \\
          & alta     & $12288 \times 4096$  & 286 & 332 & 333 & 335 & 334 \\
    \midrule
    8192  & quadrata & $8192 \times 8192$   & 289 & 335 & 336 & 336 & 337 \\
          & larga    & $8192 \times 16384$  & 294 & 342 & 344 & 342 & 345 \\
          & alta     & $24576 \times 8192$  & 296 & 346 & 344 & 347 & 345 \\
    \midrule
    16384 & quadrata & $16384 \times 16384$ & 294 & 346 & 347 & 347 & 350 \\
          & larga    & $16384 \times 32768$ & 295 & 352 & 352 & 351 & 354 \\
          & alta     & $49152 \times 16384$ & 300 & 351 & 352 & 348 & 350 \\
    \bottomrule
  \end{tabular}
\end{table}

Le prestazioni crescono con la dimensione e poi si stabilizzano:
\begin{itemize}
  \item con la matrice più piccola, $2048 \times 2048$, il kernel fa
        222--277~GFLOPS;
  \item dalla dimensione 8192 in poi si assesta tra 289 e 300~GFLOPS a $k = 3$
        e tra 335 e 354 per $k \ge 6$, vicino ai tetti della roofline (305 e
        374).
\end{itemize}

Le matrici piccole vanno più lente per due motivi:
\begin{itemize}
  \item con 2048 righe ci sono solo 2048 warp, cioè 256 blocchi, contro i 192
        che la GPU può ospitare contemporaneamente ($48$ SM $\times$ 4; 96 a
        $k = 32$, dove entrano 2 blocchi per SM), quindi
        i blocchi vengono eseguiti in due ``turni'' e nel secondo la GPU è
        occupata solo per un terzo;
  \item inoltre un'invocazione dura appena un quarto di millisecondo, e i costi
        fissi di lancio e sincronizzazione, di pochi microsecondi, pesano di più.
\end{itemize}

La forma, invece, conta poco: la differenza tra la forma migliore e la peggiore
arriva al \SI{17}{\percent} a $S = 2048$, scende al \SI{5}{\percent} a 4096 e
resta sotto il \SI{3.5}{\percent} da 8192 in su. Anche questa differenza dipende
in gran parte dalla dimensione: a parità di $S$ le forme rettangolari hanno 2
(larga) o 3 (alta) volte gli elementi della quadrata, e alle dimensioni piccole
il lavoro in più tiene la GPU più occupata.

\begin{figure}[H]
  \centering
  {\small\bfseries Soluzione finale: GFLOPS rispetto al tetto della roofline, al variare di dimensione e forma\par}\medskip
\begin{tikzpicture}
\begin{groupplot}[
  group style={group size=2 by 3, horizontal sep=1.5cm, vertical sep=1.8cm},
  width=0.5\textwidth, height=0.36\textwidth,
  xmode=log, log basis x=2,
  xmin=1600, xmax=21000, ymin=200, ymax=400,
  xtick={2048,4096,8192,16384}, xticklabels={2048,4096,8192,16384},
  xticklabel style={font=\scriptsize}, yticklabel style={font=\scriptsize},
  xlabel={dimensione di base $S$}, ylabel={GFLOPS},
  xlabel style={font=\footnotesize}, ylabel style={font=\footnotesize},
  title style={font=\footnotesize},
  grid=major, grid style={gray!20},
  every axis plot/.append style={thick, mark size=1.6pt},
]
\nextgroupplot[title={$k = 3$ \quad (tetto: 305)}, legend style={font=\scriptsize, at={(0.98,0.03)}, anchor=south east, cells={anchor=west}}]
\addplot[gray!60, dashed, domain=1600:21000, samples=2, forget plot] {305};
\addplot[blue!70!black, mark=*] coordinates {(2048,222) (4096,274) (8192,289) (16384,294)};
\addlegendentry{quadrata}
\addplot[orange!85!black, mark=square*] coordinates {(2048,249) (4096,287) (8192,294) (16384,295)};
\addlegendentry{larga ($N = 2M$)}
\addplot[green!50!black, mark=triangle*] coordinates {(2048,260) (4096,286) (8192,296) (16384,300)};
\addlegendentry{alta ($M = 3N$)}
\nextgroupplot[title={$k = 6$ \quad (tetto: 374)}]
\addplot[gray!60, dashed, domain=1600:21000, samples=2, forget plot] {374};
\addplot[blue!70!black, mark=*] coordinates {(2048,261) (4096,319) (8192,335) (16384,346)};
\addplot[orange!85!black, mark=square*] coordinates {(2048,289) (4096,335) (8192,342) (16384,352)};
\addplot[green!50!black, mark=triangle*] coordinates {(2048,302) (4096,332) (8192,346) (16384,351)};
\nextgroupplot[title={$k = 8$ \quad (tetto: 374)}]
\addplot[gray!60, dashed, domain=1600:21000, samples=2, forget plot] {374};
\addplot[blue!70!black, mark=*] coordinates {(2048,265) (4096,320) (8192,336) (16384,347)};
\addplot[orange!85!black, mark=square*] coordinates {(2048,291) (4096,334) (8192,344) (16384,352)};
\addplot[green!50!black, mark=triangle*] coordinates {(2048,305) (4096,333) (8192,344) (16384,352)};
\nextgroupplot[title={$k = 20$ \quad (tetto: 374)}]
\addplot[gray!60, dashed, domain=1600:21000, samples=2, forget plot] {374};
\addplot[blue!70!black, mark=*] coordinates {(2048,274) (4096,323) (8192,336) (16384,347)};
\addplot[orange!85!black, mark=square*] coordinates {(2048,294) (4096,337) (8192,342) (16384,351)};
\addplot[green!50!black, mark=triangle*] coordinates {(2048,310) (4096,335) (8192,347) (16384,348)};
\nextgroupplot[title={$k = 32$ \quad (tetto: 374)}]
\addplot[gray!60, dashed, domain=1600:21000, samples=2, forget plot] {374};
\addplot[blue!70!black, mark=*] coordinates {(2048,277) (4096,325) (8192,337) (16384,350)};
\addplot[orange!85!black, mark=square*] coordinates {(2048,299) (4096,340) (8192,345) (16384,354)};
\addplot[green!50!black, mark=triangle*] coordinates {(2048,310) (4096,334) (8192,345) (16384,350)};
\end{groupplot}
\end{tikzpicture}
  \caption{GFLOPS della soluzione finale (\texttt{warp\_row\_xt},
    \texttt{double}, un processo) al crescere della dimensione di base $S$,
    nelle tre forme, per ciascun $k$ della traccia. La linea tratteggiata è il
    tetto della roofline: 305~GFLOPS a $k = 3$ (limite di memoria), 374 per
    $k \ge 6$ (limite di calcolo).}\label{fig:ris-finale}
\end{figure}

\subsection{Scalabilità con MPI}\label{sec:ris-mpi}

La scalabilità con MPI è misurata su entrambe le versioni: prima quella solo
CPU, in cui ogni processo calcola su un proprio core (campagna~D), poi quella
MPI + CUDA, in cui i processi condividono l'unica GPU del server (campagna~E).

\subsubsection{Versione CPU (solo MPI)}

\paragraph{La velocità cresce con i processi.} La campagna~D misura
\texttt{cpu\_block} sulla matrice $16384 \times 16384$, da 1 a 20 processi, per
$k = 3$ e $k = 32$ (\cref{fig:speedup-cpu}). Con 20 processi su griglia
$4 \times 5$ lo speedup rispetto a un processo è:
\begin{itemize}
  \item a $k = 3$, 10,9 in modalità concorrente, cioè un'efficienza del
        \SI{55}{\percent}, e 21,1 a nodi emulati, il \SI{106}{\percent};
  \item a $k = 32$, 17,8 in modalità concorrente, l'\SI{89}{\percent}, e 17,7 a
        nodi emulati, praticamente uguale.
\end{itemize}

\begin{figure}[H]
  \centering
\begin{tikzpicture}
\begin{groupplot}[
  group style={group size=2 by 1, horizontal sep=1.5cm},
  width=0.5\textwidth, height=0.45\textwidth,
  xmin=0, xmax=21, ymin=0, ymax=23,
  xtick={1,2,4,8,10,16,20}, ytick={0,5,10,15,20},
  xticklabel style={font=\scriptsize}, yticklabel style={font=\scriptsize},
  xlabel={processi $P$}, ylabel={speedup $S(P)$},
  xlabel style={font=\footnotesize}, ylabel style={font=\footnotesize},
  title style={font=\footnotesize},
  grid=major, grid style={gray!20},
  every axis plot/.append style={thick, mark size=1.7pt},
]
\nextgroupplot[title={$k = 3$}, legend style={font=\scriptsize, at={(0.02,0.98)}, anchor=north west, cells={anchor=west}}]
\addplot[gray!60, dashed, domain=0:21, samples=2] {x};
\addlegendentry{lineare, $S(P) = P$}
\addplot[blue!70!black, mark=*] coordinates {(1,1.00) (2,1.98) (4,3.29) (8,4.85) (10,5.32) (16,8.49) (20,10.91)};
\addlegendentry{concorrente}
\addplot[orange!85!black, mark=square*] coordinates {(1,1.00) (2,1.93) (4,3.39) (8,6.57) (10,8.37) (16,14.29) (20,21.13)};
\addlegendentry{nodi emulati (\texttt{-S})}
\nextgroupplot[title={$k = 32$}]
\addplot[gray!60, dashed, domain=0:21, samples=2] {x};
\addplot[blue!70!black, mark=*] coordinates {(1,1.00) (2,1.89) (4,3.54) (8,5.35) (10,5.31) (16,13.34) (20,17.83)};
\addplot[orange!85!black, mark=square*] coordinates {(1,1.00) (2,1.88) (4,3.50) (8,6.92) (10,8.38) (16,13.43) (20,17.65)};
\end{groupplot}
\end{tikzpicture}
  \caption{Speedup della versione CPU (\texttt{cpu\_block},
    $16384 \times 16384$) al variare del numero di processi, in modalità
    concorrente e a nodi emulati, per $k = 3$ e $k = 32$. Griglie usate:
    $2 \times 1$, $2 \times 2$, $4 \times 2$, $5 \times 2$, $4 \times 4$,
    $4 \times 5$. La linea tratteggiata è lo speedup lineare.}\label{fig:speedup-cpu}
\end{figure}

Memoria condivisa: conta a $k = 3$, non a $k = 32$. A $k = 3$, con
20 processi, la versione concorrente fa 49~GFLOPS, il \SI{92}{\percent} del tetto
della memoria della CPU (53): i processi non possono andare più veloci perché si
dividono la stessa banda. Non è un problema di MPI, ma della macchina: lo
conferma la modalità a nodi emulati, dove ogni processo ha la memoria tutta per
sé e lo speedup è lineare, anzi leggermente superiore (\SI{106}{\percent}). A
$k = 32$, invece, con 20 processi si arriva a 269~GFLOPS, meno della metà del
tetto della memoria (568): la banda non è un collo di bottiglia, e infatti la
modalità concorrente e quella a nodi emulati danno lo stesso risultato.

\paragraph{La forma della griglia conta molto a $k$ grande.} A $k = 32$, con 20
processi:
\begin{itemize}
  \item con la griglia $4 \times 5$ un'invocazione dura \SI{64}{ms}, con
        speedup 17,8;
  \item con la griglia $20 \times 1$ dura \SI{216}{ms}, con speedup solo 5,3.
        Ogni processo ha bisogno di tutta $X$, che a $k = 32$ occupa
        \SI{4}{MiB}: molto più della cache L2 di un core (\SI{1}{MiB}), e 10
        processi per socket insieme superano anche la L3. $X$ deve quindi essere
        riletta continuamente dalla memoria, in concorrenza con $A$. A nodi
        emulati, dove il processo ha la L3 tutta per sé, la stessa griglia scende
        a \SI{74}{ms}. Con la griglia $4 \times 5$, invece, ogni processo usa
        solo un quinto di $X$, circa \SI{0.8}{MiB}, che sta nella L2. Lo stesso
        effetto si vede con 8 e 10 processi (griglie $4 \times 2$ e
        $5 \times 2$): con 2 sole colonne ogni processo usa metà di $X$,
        \SI{2}{MiB}, e in modalità concorrente lo speedup si ferma intorno
        a 5;
  \item con la griglia $1 \times 20$ dura \SI{74}{ms}, con speedup 15,4: qui il
        costo principale è la riduzione di $Y$, circa \SI{14}{ms}, perché con
        20 processi su una riga ognuno contribuisce a tutta $Y$, \SI{4}{MiB}.
\end{itemize}
% NOTA: possibile aggiunta di esperimenti per giustificare il ragionamento sulla
% griglia 20x1 e sulle griglie a 2 colonne (X fuori dalla cache L2/L3): ad es.
% contatori di cache (perf stat -e cache-misses) o una prova con X piu' piccola.

\subsubsection{Versione GPU (MPI + CUDA)}

Si passa ora alla
versione con backend GPU, la soluzione finale. La campagna~E misura \texttt{warp\_row\_xt} in \texttt{double} sulla matrice
$16384 \times 16384$, con $k = 8$, da 1 a 4 processi
(\cref{tab:ris-mpi-gpu}). Con un processo il nucleo fa 349~GFLOPS, mentre in modalità
concorrente aggiungere processi non migliora nulla, e nella maggior parte
delle griglie peggiora leggermente. È il comportamento atteso: c'è una sola GPU, e i processi la usano
a turno (\cref{sec:eventi-orologio}), quindi il lavoro totale è lo stesso, ma si
aggiungono le comunicazioni e i cambi di turno. Lo speedup resta intorno a 1 (0,84--1,02).

\begin{table}[H]
  \centering
  \caption{GFLOPS del nucleo con più processi sulla GPU (campagna~E,
    \texttt{warp\_row\_xt}, \texttt{double}, $16384 \times 16384$, $k = 8$).}\label{tab:ris-mpi-gpu}
  \begin{tabular}{lrrrrrr}
    \toprule
    Griglia & $1 \times 1$ & $2 \times 1$ & $1 \times 2$ & $2 \times 2$ & $4 \times 1$ & $1 \times 4$ \\
    \midrule
    concorrente  & 349 & 327 & 303 & 297  & 356  & 292  \\
    nodi emulati & --- & 664 & 650 & 1253 & 1186 & 1040 \\
    \bottomrule
  \end{tabular}
\end{table}

\paragraph{A nodi emulati la soluzione scala.} Se ogni processo avesse una GPU
per sé, la stessa soluzione andrebbe quasi in proporzione al numero di
processi: speedup 1,9 con 2 processi (\SI{95}{\percent}) e 3,6 con 4 sulla
griglia $2 \times 2$ (\SI{90}{\percent}), fino a 1253~GFLOPS. È la stima di
come andrebbe su una macchina con più GPU (\cref{sec:gpu-assign}).

\paragraph{La forma della griglia, misurata.} A nodi emulati le comunicazioni
non subiscono la contesa per la GPU, e si possono quindi confrontare con il
volume previsto in \cref{tab:volume} (la tabella dei dati scambiati da ogni
griglia con 4 processi, \cref{sec:costo-comm}). Il confronto è riportato in
\cref{tab:comm-misurate}.

\begin{table}[H]
  \centering
  \caption{Costo di broadcast e riduzione con 4 processi a nodi emulati
    (campagna~E, matrice quadrata, $k = 8$), confrontato con il volume previsto
    in \cref{tab:volume}.}\label{tab:comm-misurate}
  \begin{tabular}{lccc}
    \toprule
    Griglia & Volume previsto & Tempo misurato & \% del tempo del nucleo \\
    \midrule
    $2 \times 2$ & $2\, N k$                    & \SI{0.24}{ms} & \SI{7}{\percent}  \\
    $4 \times 1$ & $3\, N k$ (solo broadcast)  & \SI{0.43}{ms} & \SI{12}{\percent} \\
    $1 \times 4$ & $3\, N k$ (solo riduzione)  & \SI{0.94}{ms} & \SI{23}{\percent} \\
    \bottomrule
  \end{tabular}
\end{table}

L'ordine è quello previsto dal volume: la $2 \times 2$ muove meno dati ed è la
più veloce. C'è però una differenza che il solo volume non prevede: la
$4 \times 1$ e la $1 \times 4$ muovono lo stesso volume (\SI{1}{MiB} per
messaggio), ma la riduzione (\SI{0.94}{ms}) costa più del doppio del broadcast
(\SI{0.39}{ms}; gli \SI{0.43}{ms} della tabella comprendono anche la
riduzione), perché oltre a trasferire i dati deve sommarli.

\paragraph{In conclusione: CPU e GPU a confronto.} Le due versioni scalano con
MPI per motivi diversi. Sulla CPU i processi calcolano davvero in parallelo, uno
per core: a $k = 32$ con 20 processi lo speedup è 17,8, e a $k = 3$ il limite è
la banda di memoria condivisa, non MPI. Sulla GPU, con un'unica scheda, più
processi non aiutano, mentre a nodi emulati, cioè con una GPU per processo, la
soluzione scalerebbe invece quasi linearmente (speedup 3,6 con 4 processi), con
una stima ottimistica (\cref{sec:modalita}). A nodi emulati il peso delle
comunicazioni è però molto diverso: sulla CPU restano intorno al \SI{3}{\percent} del tempo del
nucleo (griglia $4 \times 4$, $k = 32$), sulla GPU vanno dal \SI{7}{\percent}
al \SI{23}{\percent} a seconda della griglia (\cref{tab:comm-misurate}). La
differenza viene dal calcolo, che
sulla GPU è molto più rapido e lascia quindi alle comunicazioni una quota
maggiore del tempo.

\subsection{Padding e budget dei registri}\label{sec:ris-padding}
La campagna~F confronta
\texttt{-p 1}, con tutte le righe allineate, e \texttt{-p 0}, che disattiva
padding e allocazione allineata, sulle dimensioni 8191 e 16383
(\cref{sec:dimensioni}), in \texttt{double} con $k = 3$, 8 e 32: sulla GPU con il
kernel migliore e un processo, sulla CPU con \texttt{cpu\_block} e 20 processi
($4 \times 5$). Qui una riga occupa $8N$ byte, non multiplo né di 64 né
di 128, quindi senza padding quasi nessuna riga parte allineata (sulla GPU solo
la prima, perché \texttt{cudaMalloc} allinea a \SI{256}{byte}). Le differenze
sono minime: sulla GPU da \SI{-0.1}{\percent} a \SI[retain-explicit-plus]{+1.1}{\percent}, sulla CPU da
\SI{-3.4}{\percent} a \SI[retain-explicit-plus]{+1.6}{\percent}, cioè dentro la variabilità delle misure (in un caso il coefficiente
di variazione arriva all'\SI{11.8}{\percent}).

 Il motivo più probabile per cui
l'effetto non si vede è che le righe di $A$ sono consecutive in memoria: il
settore ``in più'' letto all'inizio di una riga disallineata contiene la fine
della riga precedente, che è già stata letta o lo sarà a breve. La cache L2 lo
serve quindi senza leggerlo di nuovo dalla memoria, e il traffico verso la
memoria globale resta quasi lo stesso. Per la CPU vale lo stesso discorso, con
le linee di cache da \SI{64}{byte}.
% NOTA: la spiegazione (settore in piu' servito dalla L2 perche' appartiene alla
% riga vicina) e' un'ipotesi coerente con i dati, non misurata: possibile verifica
% con il profiler (byte letti dalla DRAM con e senza padding).

Il padding resta comunque attivo per default: non costa nulla di misurabile e
garantisce l'allineamento anche in condizioni diverse da quelle provate, ad
esempio con blocchi non consecutivi in memoria.

\paragraph{La scelta del budget dei registri.} Lo script
\texttt{budget\_cpu.sh} confronta i budget 32, 64 e 128 di \texttt{cpu\_block}
(\cref{sec:vettorizzazione}) sulla matrice $16384 \times 16384$
(\cref{tab:budget}).

\begin{table}[H]
  \centering
  \caption{GFLOPS di \texttt{cpu\_block} al variare del budget dei registri
    (\texttt{double}, $16384 \times 16384$).}\label{tab:budget}
  \begin{tabular}{lrrrr}
    \toprule
     & \multicolumn{2}{c}{1 processo} & \multicolumn{2}{c}{20 processi ($4 \times 5$)} \\
    \cmidrule(lr){2-3}\cmidrule(lr){4-5}
    Budget & $k=20$ & $k=32$ & $k=20$ & $k=32$ \\
    \midrule
    32 ($R_B = 1$) & 5,0           & 5,0           & 147          & 148          \\
    64             & 13,0          & 10,2          & \textbf{298} & 251          \\
    \textbf{128}   & \textbf{17,8} & \textbf{15,2} & 294          & \textbf{265} \\
    \bottomrule
  \end{tabular}
\end{table}

\begin{itemize}
  \item Passare da 32 a 64 raddoppia le prestazioni con 20 processi, e le
        moltiplica per 2--3 con un processo solo. A $k \le 8$ i tre budget
        danno invece risultati quasi uguali, perché $R_B$ è comunque alto.
  \item Tra 64 e 128, con 20 processi, la differenza è piccola: appena ogni
        valore di $X$ serve almeno due righe, il limite diventa la banda di
        memoria condivisa, e altre righe non aiutano.
  \item Con un solo processo 128 è nettamente migliore: \SI[retain-explicit-plus]{+37}{\percent} a $k = 20$ e
        \SI[retain-explicit-plus]{+49}{\percent} a $k = 32$. Un core da solo ha banda in abbondanza, e più riuso
        di $X$ significa più velocità. È questo che ha deciso la scelta: il
        riferimento a un processo, usato per lo speedup, deve sfruttare il
        kernel al meglio (\cref{sec:modelli}).
\end{itemize}

\subsection{Effetti sulla misura}\label{sec:ris-effetti}
\begin{itemize}
  \item \textbf{La frequenza della GPU.} \texttt{frequenza.sh} ripete tre volte
        la stessa configurazione ($8192 \times 8192$, $k = 6, 8, 20$) con un
        riscaldamento breve (2 invocazioni ciascuno) e uno lungo (200 ciascuno). A $k = 20$, con
        quello breve le prime due ripetizioni danno 286 e 285~GFLOPS e la terza
        340, con quello lungo tutte circa 339: un \SI{19}{\percent} di
        differenza dovuto solo alla frequenza, che sale dopo un periodo di
        lavoro continuo. Da qui il riscaldamento a tempo
        (introdotto in \cref{sec:riscaldamento}).
  \item \textbf{Eventi CUDA e orologio.} Con un processo, o a nodi emulati, le
        due misure coincidono (\SI{12.17}{ms} e \SI{12.16}{ms} con un processo,
        \SI{3.19}{ms} e \SI{3.18}{ms} sulla griglia $2 \times 2$ a nodi
        emulati). In concorrenza, sulla $2 \times 2$, si separano: il lavoro
        vero dura \SI{3.2}{ms}, gli eventi misurano \SI{10.2}{ms} perché
        includono i turni ceduti agli altri processi, e l'orologio
        \SI{14.1}{ms} perché aggiunge circa \SI{3.9}{ms} di attesa prima del
        kernel. Con gli eventi il nucleo sembrerebbe più veloce di quanto sia:
        per questo si usa l'orologio (\cref{sec:eventi-orologio}).
  \item \textbf{I trasferimenti esclusi dal tempo.} Le copie di $X$ verso la
        GPU e di $Y$ verso la CPU, escluse dal tempo del nucleo come chiede la
        traccia, pesano circa il \SI{5}{\percent} del
        nucleo con la matrice $16384 \times 16384$ (\SI{0.6}{ms} su
        \SI{12.4}{ms} a $k = 8$) e tra il \SI{31}{\percent} e il
        \SI{39}{\percent} con la $2048 \times 2048$.
  \item \textbf{Il preprocessamento.} Generare $A$ e copiarla sulla GPU costa
        circa \SI{2.3}{s} con la $16384 \times 16384$, quanto circa 190
        invocazioni a $k = 8$. Si ammortizza solo se il nucleo viene invocato
        molte volte sulla stessa $A$, come nelle misure di questo progetto.
\end{itemize}


% =====================================================================
\section{Conclusioni}\label{sec:conclusioni}

\paragraph{Sintesi dei risultati.} Il progetto realizza il prodotto
$Y \leftarrow A\,X$ tra una matrice densa e un multivettore, con MPI su una
griglia bidimensionale di processi e due versioni per il calcolo locale: CUDA
sulla GPU e solo MPI sulla CPU.
\begin{itemize}
  \item Sulla GPU il kernel finale \texttt{warp\_row\_xt} lavora praticamente al
        limite dell'hardware: in \texttt{double} arriva al \SI{98}{\percent} del
        tetto della memoria a $k = 3$ e al \SIrange{94}{95}{\percent} del picco
        di calcolo per $k \ge 6$, mentre in \texttt{float} segue il tetto della memoria fino a $k = 8$. Alla dimensione $S = 16384$ fa tra 294 e
        354~GFLOPS in \texttt{double}, nelle tre forme richieste dalla traccia.
  \item Sulla CPU, con 20 processi, \texttt{cpu\_block} arriva a circa il
        \SI{90}{\percent} del tetto della memoria fino a $k = 20$.
  \item La GPU è da 6 volte più veloce della CPU a $k = 3$ a 1,3 volte a
        $k = 32$.
  \item Con MPI, sulla CPU la soluzione scala bene quando la memoria non è il
        limite (speedup 17,8 su 20 processi a $k = 32$). A nodi emulati la
        versione GPU scalerebbe quasi linearmente (speedup 3,6 su 4 processi),
        mentre sul server reale, con una sola GPU, più processi non portano
        vantaggi.
\end{itemize}

\paragraph{Cosa si è imparato.}
\begin{itemize}
  \item Per questo nucleo conta soprattutto come si leggono i dati. Ogni
        miglioramento dei kernel GPU, da \texttt{thread\_row} a
        \texttt{warp\_row\_xt}, non ha aggiunto calcoli, ma ha cambiato il modo
        di leggere $A$ e $X$: una volta sola invece di $k$ volte, in modo
        coalescente invece che sparso. Lo stesso vale sulla CPU, dove il
        guadagno di \texttt{cpu\_block} viene dal riuso di $X$ tra più righe.
  \item Non tutte le ottimizzazioni rendono. La memoria shared di
        \texttt{warp\_row\_smem} e il padding delle righe non hanno portato
        vantaggi misurabili, probabilmente perché le cache della GPU svolgevano già quel
        lavoro: misurare è stato l'unico modo per saperlo.
  \item Le misure vanno controllate quanto il codice. Il riscaldamento a tempo e
        la scelta dell'orologio al posto degli eventi CUDA: entrambe le scelte
        sono nate da un risultato che non tornava. Lo stesso vale per il picco di calcolo, che
        con la frequenza dichiarata dalla scheda tecnica risultava più basso
        delle prestazioni misurate.
  \item Il miglior risultato dipende dall'hardware a disposizione. Su un server
        con una sola GPU la versione CUDA non scala con i processi, mentre
        quella solo MPI usa tutti i 20 core, e la modalità a nodi emulati ha
        permesso di valutare entrambe anche nello scenario che il server non
        offre.
\end{itemize}

\paragraph{Sviluppi futuri.}
\begin{itemize}
  \item \emph{Più righe per warp sulla GPU.} In \texttt{float}, a $k$ grande,
        \texttt{warp\_row\_xt} è probabilmente limitato dalle letture di $X$
        dalla cache, e cuBLAS lo supera (\cref{sec:gpu-riepilogo}). Se ogni warp
        calcolasse $R$ righe insieme, ogni valore di $X$ letto servirebbe $R$
        righe, come in \texttt{cpu\_block}. In cambio i warp scendono a $M/R$ e
        ogni thread tiene $R \times k$ accumulatori, quindi sullo SM entrano
        meno warp: conviene solo con $k$ grande. Non è stato implementato perché
        la soluzione finale usa il \texttt{double}, dove \texttt{warp\_row\_xt}
        è già al \SIrange{94}{95}{\percent} del picco di calcolo: servirebbe
        solo in \texttt{float}.
  \item \emph{MPI CUDA-aware.} Con un'installazione di MPI capace di leggere
        direttamente la memoria della GPU, le copie di $X$ e $Y$ attraverso la
        CPU sparirebbero (\cref{sec:calcpar}), e su matrici piccole pesano fino a
        circa il \SI{40}{\percent} del tempo (\cref{sec:ris-effetti}).
  \item \emph{Più processi sulla stessa GPU con MPS.} Il servizio MPS di NVIDIA
        permetterebbe ai kernel di processi diversi di essere eseguiti insieme
        invece che a turno (\cref{sec:eventi-orologio}).
  \item \emph{Verifiche più dirette.} Alcune spiegazioni, come l'effetto della
        cache con la griglia $20 \times 1$, l'effetto ridotto del padding e il
        limite di \texttt{warp\_row\_xt} in \texttt{float} a $k$ grande, sono
        ipotesi coerenti con
        i dati ma non misurate: si potrebbero confermare con i contatori
        dell'hardware o con il profiler.
\end{itemize}

% =====================================================================
\appendix
\section{Opzioni del programma}\label{app:opzioni}

Il programma si lancia con \texttt{mpirun} e si configura interamente da riga
di comando. Le opzioni \texttt{-m}, \texttt{-n} e \texttt{-k} sono
obbligatorie, le altre hanno un valore di default (\cref{tab:uso}).
Un'opzione mancante o non valida fa stampare al programma un riepilogo d'uso e
lo termina. Il processo~0 stampa su standard output una riga CSV con i
risultati (\cref{sec:fasi}).

\begin{table}[H]
  \centering
  \small
  \caption{Opzioni del programma.}\label{tab:uso}
  \begin{tabularx}{\textwidth}{@{}lX@{}}
    \toprule
    \textbf{Opzione} & \textbf{Significato} \\
    \midrule
    \multicolumn{2}{@{}l}{\emph{Obbligatorie}} \\
    \texttt{-m M} & Righe di $A$ e di $Y$. \\
    \texttt{-n N} & Colonne di $A$ e righe di $X$. \\
    \texttt{-k K} & Colonne di $X$ e di $Y$. \\
    \midrule
    \multicolumn{2}{@{}l}{\emph{Con un valore}} \\
    \texttt{-r Pr} & Righe della griglia di processi (default 1). \\
    \texttt{-c Pc} & Colonne della griglia; $P_r \cdot P_c$ deve essere
      uguale al numero di processi lanciati (default 1). \\
    \texttt{-v n} & Kernel da usare (\cref{tab:kernel}); $-1$ = cuBLAS
      (default: il più veloce del backend, \texttt{warp\_row\_xt} o
      \texttt{cpu\_block}). \\
    \texttt{-p 0|1} & Padding e allocazione allineata delle righe di $A$
      (default 1). \\
    \texttt{-R n} & Invocazioni misurate (default 20). \\
    \texttt{-W n} & Numero minimo di invocazioni di riscaldamento
      (default 3). \\
    \texttt{-T s} & Durata minima del riscaldamento in secondi; con
      \texttt{-W 0 -T 0} non c'è riscaldamento (default 1). \\
    \midrule
    \multicolumn{2}{@{}l}{\emph{Senza valore}} \\
    \texttt{-S} & Modalità a nodi emulati: i processi calcolano a turno. \\
    \texttt{-C} & Verifica con il riferimento seriale. Stampa l'esito su
      standard error e riempie la colonna \texttt{relerr} del CSV, che
      altrimenti vale $-1$. \\
    \texttt{-H} & Stampa prima della riga dei risultati la riga di
      intestazione del CSV. \\
    \bottomrule
  \end{tabularx}
\end{table}

\noindent\begin{minipage}{\textwidth}
Ad esempio:
\begin{lstlisting}[language=bash,numbers=none,frame=none]
mpirun -np 4 bin/mxm_cuda_double -m 16384 -n 16384 -k 8 -r 2 -c 2 -H -C
\end{lstlisting}
\end{minipage}

% =====================================================================
\clearpage
\phantomsection
\addcontentsline{toc}{section}{\refname}
\begin{thebibliography}{9}

\bibitem{salmon2011}
J.~K. Salmon, M.~A. Moraes, R.~O. Dror, D.~E. Shaw.
\newblock Parallel random numbers: as easy as 1, 2, 3.
\newblock In \emph{Proceedings of the International Conference for High
Performance Computing, Networking, Storage and Analysis (SC '11)}, 2011.

\bibitem{steele2014}
G.~L. Steele~Jr., D.~Lea, C.~H. Flood.
\newblock Fast splittable pseudorandom number generators.
\newblock In \emph{Proceedings of OOPSLA '14}, ACM, pp.~453--472, 2014.
\newblock doi:10.1145/2660193.2660195.

\bibitem{stafford2011}
D.~Stafford.
\newblock Better bit mixing -- improving on MurmurHash3's 64-bit finalizer.
\newblock 2011.
\newblock \url{https://zimbry.blogspot.com/2011/09/better-bit-mixing-improving-on.html}

\bibitem{vigna}
S.~Vigna.
\newblock xoshiro / xoroshiro generators and the PRNG shootout.
\newblock \url{https://prng.di.unimi.it/}

\bibitem{williams2009}
S.~Williams, A.~Waterman, D.~Patterson.
\newblock Roofline: an insightful visual performance model for multicore
architectures.
\newblock \emph{Communications of the ACM}, 52(4):65--76, 2009.
\newblock doi:10.1145/1498765.1498785.

\end{thebibliography}

% ---------------------------------------------------------------------
% Esempi pronti da copiare:
%
% \begin{lstlisting}[caption={Esempio di kernel CUDA}, label={lst:kernel}]
% __global__ void scale(float *v, float a, int n) { ... }
% \end{lstlisting}
%
% \begin{table}[H]
%   \centering
%   \caption{Tempi di esecuzione}\label{tab:tempi}
%   \begin{tabular}{lrr}
%     \toprule
%     Versione & Tempo (\si{\milli\second}) & Speedup \\
%     \midrule
%     Seriale  & \num{1250.4} & 1.00 \\
%     \bottomrule
%   \end{tabular}
% \end{table}
%
% \begin{figure}[H]
%   \centering
%   \includegraphics[width=0.8\textwidth]{grafico}   % file in img/
%   \caption{Speedup al variare dei thread}\label{fig:speedup}
% \end{figure}
%
% Riferimenti: \cref{sec:mpi}, \cref{tab:tempi}

\end{document}
